% IV : rédiger un résumé de texte — Français Tle Tech — Cours signé OnBuch+
% Programme officiel MINESEC. Compiler avec : tectonic N16-iv-rediger-resume-texte.tex
\newcommand{\DOCMATIERE}{Français}
\newcommand{\DOCNIVEAU}{Tle Tech}
\newcommand{\DOCMODULE}{Français – classe de Terminale technique}
\newcommand{\DOCLECON}{N16}
\newcommand{\DOCTITRE}{IV : rédiger un résumé de texte}
% OnBuch+ — préambule « cours signés » ENRICHI (compilé avec tectonic / XeLaTeX)
% Système visuel « Pop » d'OnBuch+ : fond crème (jamais blanc), bordures encre
% épaisses, ombres « dures » décalées, titres Archivo Black en capitales.
% Version MATHÉMATIQUES : graphes (pgfplots/tikz), boîtes pédagogiques
% (définition, propriété, méthode, attention, le savais-tu, exercice résolu,
% exercice, TP, bilan), page de garde et signature OnBuch+ ancrée en bas.
%
% Macros attendues dans main.tex AVANT \input{preamble} :
%   \DOCMATIERE  \DOCNIVEAU  \DOCMODULE  \DOCLECON (numéro)  \DOCTITRE
\documentclass[11pt]{article}

\usepackage{fontspec}
\usepackage[french]{babel}
\usepackage[a4paper,margin=2cm,top=3.4cm,bottom=2.8cm,headheight=1.4cm,footskip=1.3cm]{geometry}
\usepackage{amsmath,amssymb,amsfonts}
\usepackage{mathtools} % \xmapsto, \coloneqq... souvent utilisés par les modèles
\usepackage{cancel} % simplifications barrées (bilans de forces)
% \abs{z} (module, valeur absolue) et \norm{u} (norme), utilisés par les modèles
\providecommand{\abs}{}\let\abs\relax\DeclarePairedDelimiter{\abs}{\lvert}{\rvert}
\providecommand{\norm}{}\let\norm\relax\DeclarePairedDelimiter{\norm}{\lVert}{\rVert}
\usepackage[table]{xcolor} % [table] : fournit aussi \rowcolors (alternance de lignes)
\usepackage{pagecolor}
\usepackage{tikz}
\usetikzlibrary{arrows.meta,positioning,calc,shapes.geometric,shapes.misc,shapes.arrows,decorations.pathmorphing,decorations.pathreplacing,decorations.markings,patterns,shadows,mindmap,trees,fit,backgrounds,matrix,intersections,angles,quotes}
\usepackage{pgfplots}
\usepackage[version=4]{mhchem} % équations nucléaires \ce{^{235}_{92}U}
\usepackage[european,siunitx]{circuitikz} % schémas électriques
\pgfplotsset{compat=1.18}
\usepgfplotslibrary{fillbetween,groupplots} % aires entre courbes (calcul intégral)
\usepackage{enumitem}
\usepackage{fancyhdr}
% [most] charge toutes les bibliothèques tcolorbox (listings, minted, raster,
% documentation...) et gonfle inutilement la mémoire TeX sur un document long
% (~300 boîtes) : on ne charge que ce qui est réellement utilisé.
\usepackage[skins,breakable]{tcolorbox}
\usepackage{graphicx}
\usepackage{colortbl}
\usepackage{booktabs}
\usepackage{tabularx}
\usepackage{multirow}
\usepackage{diagbox} % case d'en-tête diagonale des tableaux à double entrée
% Colonnes extensibles centrée / gauche / droite (C, L, R), souvent utilisées par les modèles
\newcolumntype{C}{>{\centering\arraybackslash}X}
\newcolumntype{L}{>{\raggedright\arraybackslash}X}
\newcolumntype{R}{>{\raggedleft\arraybackslash}X}
\usepackage{array}
\usepackage{multicol}
\usepackage{siunitx}
% parse-numbers=false : accepte aussi \SI{2,5 \times 10^{-3}}{...}
\sisetup{locale=FR,output-decimal-marker={,},group-minimum-digits=4,parse-numbers=false}
\DeclareSIUnit\ohm{\ensuremath{\symup{\Omega}}} % U+2126 absent de la police
\DeclareSIUnit\division{div} % réglages d'oscilloscope (ms/div, V/div)
\DeclareSIUnit\tr{tr} % tours (vitesse de rotation, ex. \SI{1500}{\tr\per\minute})
\DeclareSIUnit\molar{M} % concentration molaire (ex. \SI{3}{\molar}, \SI{1}{\micro\molar})
\DeclareSIUnit\an{an} % année (le modèle confond parfois avec \year, un compteur TeX) — ex. \SI{5}{\kilo\meter\per\an}
\DeclareSIUnit\FCFA{FCFA} % franc CFA (ex. \SI{500}{\FCFA\per\kilo\gram})
\DeclareSIUnit\degreeBrix{\textdegree Brix} % teneur en sucre d'un sirop/jus (réfractométrie)
\DeclareSIUnit\month{mois} % le modèle utilise parfois \month, un compteur TeX, comme unité de durée
\DeclareSIUnit\microliter{\micro\liter} % le modèle écrit parfois \microliter en un seul token
\DeclareSIUnit\million{millions} % dénombrement (ex. \SI{15}{\million\per\mL} spermatozoïdes)
\usepackage{qrcode}
% \degree : commande fréquemment utilisée par les modèles pour un angle ou une
% température en dehors de \SI{}{} (non fournie par les paquets chargés).
\providecommand{\degree}{^{\circ}}
% \qed (fin de démonstration), utilisé par les modèles sans amsthm
\providecommand{\qed}{\hfill$\square$}
% \barre : double barre des valeurs interdites dans un tableau de variations (tkz-tab)
\providecommand{\barre}{\ensuremath{\|}}
% environnement proof (amsthm non chargé) utilisé par les modèles
\ifdefined\proof\else\newenvironment{proof}[1][Démonstration]{\par\noindent\textit{#1.}\ }{\qed\par}\fi
\usepackage{lastpage}
\usepackage{hyperref}
\hypersetup{hidelinks}

% ============================================================
% Palette « Pop » OnBuch+ — mêmes hex que la classe P (plus_ui.dart)
% ============================================================
\definecolor{poporange}{HTML}{F59321}
\definecolor{popdark}{HTML}{E07A0C}
\definecolor{popcream}{HTML}{FFF6E8}
\definecolor{popink}{HTML}{1C1714}
\definecolor{poppurple}{HTML}{7A5AE0}
\definecolor{popgreen}{HTML}{2E9E6A}
\definecolor{poppink}{HTML}{E0457B}
\definecolor{popblue}{HTML}{2F7DE1}
\definecolor{popgold}{HTML}{C98A12}
\definecolor{popmuted}{HTML}{8A7F76}
\definecolor{popline}{HTML}{EADFD2}
\definecolor{popgray}{HTML}{8A7F76}
\colorlet{popgrey}{popgray}
% Alias tolérants (le modèle nomme parfois les couleurs différemment)
\colorlet{popwhite}{white}
\colorlet{popblack}{popink}
\colorlet{popred}{poppink}
\colorlet{popyellow}{popgold}
% Teintes claires (fonds de tableaux/encadrés) — le modèle nomme parfois
% les couleurs avec un suffixe L (ex. popblueL) sans les avoir définies.
\colorlet{poporangeL}{poporange!20}
\colorlet{popdarkL}{popdark!20}
\colorlet{poppurpleL}{poppurple!20}
\colorlet{popgreenL}{popgreen!20}
\colorlet{poppinkL}{poppink!20}
\colorlet{popblueL}{popblue!20}
\colorlet{popgoldL}{popgold!20}
\colorlet{popredL}{popred!20}
\colorlet{popgrayL}{popgray!20}
% Teintes claires pour fonds de tableaux / zones
\colorlet{popblueL}{popblue!10!white}
\colorlet{poporangeL}{poporange!14!white}
\colorlet{popgreenL}{popgreen!12!white}
\colorlet{poppurpleL}{poppurple!12!white}
\colorlet{poppinkL}{poppink!12!white}
\colorlet{popgoldL}{popgold!14!white}

\pagecolor{popcream}

% ============================================================
% Typographie
% ============================================================
\defaultfontfeatures{Ligatures=TeX}
\setmainfont{PlusJakartaSans-Medium.ttf}[Path=../../../fonts/,BoldFont=PlusJakartaSans-Bold.ttf]
% Maths en sans-serif (Fira Math) avec chiffres et lettres droites de Plus
% Jakarta, pour que les formules se fondent dans le texte.
\usepackage[math-style=ISO,bold-style=ISO]{unicode-math}
\setmathfont{FiraMath-Regular.otf}[Path=../../../fonts/]
\setmathfont{PlusJakartaSans-Medium.ttf}[Path=../../../fonts/,range={up/num,up/{Latin,latin},\mathup}]
\usepackage{icomma} % virgule décimale sans espace en mode math
% Caractères mathématiques Unicode absents des polices de texte (Plus Jakarta,
% Archivo Black) : rendus par leur équivalent mathématique, en texte comme en
% formule — sinon ils s'affichent en carré vide (ex. « Arithmétique dans ℤ »).
\usepackage{newunicodechar}
\newunicodechar{ℝ}{\ensuremath{\mathbb{R}}}
\newunicodechar{ℤ}{\ensuremath{\mathbb{Z}}}
\newunicodechar{ℂ}{\ensuremath{\mathbb{C}}}
\newunicodechar{ℕ}{\ensuremath{\mathbb{N}}}
\newunicodechar{ℚ}{\ensuremath{\mathbb{Q}}}
\newunicodechar{𝔻}{\ensuremath{\mathbb{D}}}
\newunicodechar{α}{\ensuremath{\alpha}}
\newunicodechar{β}{\ensuremath{\beta}}
\newunicodechar{γ}{\ensuremath{\gamma}}
\newunicodechar{δ}{\ensuremath{\delta}}
\newunicodechar{ε}{\ensuremath{\varepsilon}}
\newunicodechar{θ}{\ensuremath{\theta}}
\newunicodechar{λ}{\ensuremath{\lambda}}
\newunicodechar{μ}{\ensuremath{\mu}}
\newunicodechar{σ}{\ensuremath{\sigma}}
\newunicodechar{τ}{\ensuremath{\tau}}
\newunicodechar{φ}{\ensuremath{\varphi}}
\newunicodechar{ψ}{\ensuremath{\psi}}
\newunicodechar{ω}{\ensuremath{\omega}}
\newunicodechar{Σ}{\ensuremath{\Sigma}}
% majuscules produites par \MakeUppercase dans les titres (α-aminés -> Α)
\newunicodechar{Α}{\ensuremath{\alpha}}
\newunicodechar{Β}{\ensuremath{\beta}}
\newunicodechar{Γ}{\ensuremath{\Gamma}}
\newunicodechar{Δ}{\ensuremath{\Delta}}
\newunicodechar{Ε}{\ensuremath{\varepsilon}}
\newunicodechar{Θ}{\ensuremath{\Theta}}
\newunicodechar{Λ}{\ensuremath{\Lambda}}
\newunicodechar{Μ}{\ensuremath{\mu}}
\newunicodechar{Τ}{\ensuremath{\tau}}
\newunicodechar{Φ}{\ensuremath{\Phi}}
\newunicodechar{Ψ}{\ensuremath{\Psi}}
\newunicodechar{Ω}{\ensuremath{\Omega}}
\newunicodechar{↦}{\ensuremath{\mapsto}}
\newunicodechar{∈}{\ensuremath{\in}}
\newunicodechar{∉}{\ensuremath{\notin}}
\newunicodechar{∧}{\ensuremath{\wedge}}
\newunicodechar{⇔}{\ensuremath{\Leftrightarrow}}
\newunicodechar{⟺}{\ensuremath{\Longleftrightarrow}}
\newunicodechar{⇒}{\ensuremath{\Rightarrow}}
\newunicodechar{⟹}{\ensuremath{\Longrightarrow}}
\newunicodechar{∘}{\ensuremath{\circ}}
\newunicodechar{∪}{\ensuremath{\cup}}
\newunicodechar{∩}{\ensuremath{\cap}}
\newunicodechar{≡}{\ensuremath{\equiv}}
\newunicodechar{⊂}{\ensuremath{\subset}}
\newunicodechar{⊥}{\ensuremath{\perp}}
\newunicodechar{∃}{\ensuremath{\exists}}
\newunicodechar{∀}{\ensuremath{\forall}}
\newunicodechar{∛}{\ensuremath{\sqrt[3]{\,}}}
\newunicodechar{∜}{\ensuremath{\sqrt[4]{\,}}}
\newunicodechar{⃗}{\ensuremath{{}^{\to}}}
\newunicodechar{ⁿ}{\ifmmode^{n}\else\textsuperscript{\ensuremath{n}}\fi}
\newunicodechar{ˣ}{\ifmmode^{x}\else\textsuperscript{\ensuremath{x}}\fi}
\newunicodechar{ᵇ}{\ifmmode^{b}\else\textsuperscript{\ensuremath{b}}\fi}
\newunicodechar{ʳ}{\ifmmode^{r}\else\textsuperscript{\ensuremath{r}}\fi}
\newunicodechar{ᵘ}{\ifmmode^{u}\else\textsuperscript{\ensuremath{u}}\fi}
\newunicodechar{ʸ}{\ifmmode^{y}\else\textsuperscript{\ensuremath{y}}\fi}
\newunicodechar{ᵃ}{\ifmmode^{a}\else\textsuperscript{\ensuremath{a}}\fi}
\newunicodechar{ᵉ}{\ifmmode^{e}\else\textsuperscript{\ensuremath{e}}\fi}
\newunicodechar{ᵏ}{\ifmmode^{k}\else\textsuperscript{\ensuremath{k}}\fi}
\newunicodechar{ᵖ}{\ifmmode^{p}\else\textsuperscript{\ensuremath{p}}\fi}
\newunicodechar{ᴺ}{\ifmmode^{N}\else\textsuperscript{\ensuremath{N}}\fi}
\newunicodechar{ᶜ}{\ifmmode^{c}\else\textsuperscript{\ensuremath{c}}\fi}
\newunicodechar{ʰ}{\ifmmode^{h}\else\textsuperscript{\ensuremath{h}}\fi}
\newunicodechar{ᵠ}{\ifmmode^{q}\else\textsuperscript{\ensuremath{q}}\fi}
\newunicodechar{ₙ}{\ifmmode_{n}\else\textsubscript{\ensuremath{n}}\fi}
\newunicodechar{ₐ}{\ifmmode_{a}\else\textsubscript{\ensuremath{a}}\fi}
\newunicodechar{ᵢ}{\ifmmode_{i}\else\textsubscript{\ensuremath{i}}\fi}
\newunicodechar{ᵣ}{\ifmmode_{r}\else\textsubscript{\ensuremath{r}}\fi}
\newunicodechar{ₖ}{\ifmmode_{k}\else\textsubscript{\ensuremath{k}}\fi}
\newunicodechar{ᵦ}{\ifmmode_{\beta}\else\textsubscript{\ensuremath{\beta}}\fi}
\newunicodechar{ₓ}{\ifmmode_{x}\else\textsubscript{\ensuremath{x}}\fi}
\newfontfamily\popdisplay{ArchivoBlack-Regular.ttf}[Path=../../../fonts/]
\newfontfamily\poplogo{SpaceGrotesk-Bold.ttf}[Path=../../../fonts/]
\newfontfamily\popsemi{PlusJakartaSans-SemiBold.ttf}[Path=../../../fonts/]
\newfontfamily\popxbold{PlusJakartaSans-ExtraBold.ttf}[Path=../../../fonts/]
\newfontfamily\popxboldls{PlusJakartaSans-ExtraBold.ttf}[Path=../../../fonts/,LetterSpace=3.0]

\setlist[enumerate]{leftmargin=1.7em,itemsep=5pt,topsep=4pt}
\setlist[itemize]{leftmargin=1.7em,itemsep=4pt,topsep=4pt,label={\color{poporange}\textbullet}}
\setlength{\parindent}{0pt}
\setlength{\parskip}{5pt}
\renewcommand{\arraystretch}{1.25}

% pgfplots : style Pop par défaut
\pgfplotsset{
  every axis/.append style={axis line style={popink,line width=1pt},tick style={popink},
    grid style={popline,line width=0.5pt},label style={font=\small},tick label style={font=\footnotesize}},
  popcurve/.style={poporange,line width=1.8pt,no markers,smooth},
}
% Même style disponible dans un \draw TikZ simple (hors pgfplots)
\tikzset{popcurve/.style={poporange,line width=1.8pt}}
\tikzset{popgrid/.style={popline,very thin}}
\colorlet{popcurve}{poporange} % le nom du style est parfois utilisé comme couleur

% ============================================================
% Pill / squiggle
% ============================================================
\newcommand{\poppill}[3]{%
  \tikz[baseline=-0.55ex]\node[draw=popink,line width=1.2pt,rounded corners=2.6pt,%
    fill=#2,inner xsep=6.5pt,inner ysep=3pt]{%
    \popxboldls\fontsize{7}{7}\selectfont\color{#3}\MakeUppercase{#1}};%
}
\newcommand{\popsquiggle}[1]{%
  \tikz[baseline=-0.4ex]{
    \draw[#1,line width=2.6pt,line cap=round]
      plot [smooth] coordinates {(0,0.09) (0.34,0.01) (0.68,0.21) (1.02,0.01) (1.36,0.10)};
  }%
}
% Mot-clé surligné (marqueur orange sous le texte)
\newcommand{\cle}[1]{{\popxbold\color{popdark}#1}}

% ============================================================
% En-tête / pied
% ============================================================
\pagestyle{fancy}
\fancyhf{}
\renewcommand{\headrulewidth}{0pt}
\renewcommand{\footrulewidth}{0pt}
\lhead{\raisebox{-0.15ex}{\poplogo\fontsize{13}{13}\selectfont\color{popink}OnBuch\color{poporange}+}}
\rhead{\poppill{\DOCMATIERE\ \textbullet\ \DOCNIVEAU\ \textbullet\ Leçon \DOCLECON}{white}{popink}}
\lfoot{\popxbold\fontsize{7.6}{7.6}\selectfont\color{popmuted}\MakeUppercase{Cours signé OnBuch+}\ \textbullet\ {\popsemi\DOCTITRE}}
\rfoot{\tikz[baseline=-0.6ex]\node[rounded corners=8.5pt,fill=popink,minimum height=17pt,minimum width=17pt,inner xsep=5pt,inner sep=0]{\popxbold\fontsize{8}{8}\selectfont\color{popcream}\thepage\,/\,\pageref*{LastPage}};}

% ============================================================
% Titres
% ============================================================
\newcommand{\chaptitre}[2]{%
  \begin{tcolorbox}[enhanced,colback=poporange,colframe=popink,boxrule=2.6pt,arc=11pt,
      drop shadow=popink,left=15pt,right=15pt,top=12pt,bottom=11pt,before skip=3pt,after skip=18pt]
    \poppill{#2}{popink}{popcream}\par\vspace{9pt}
    {\raggedright\hyphenpenalty=10000\exhyphenpenalty=10000\popdisplay\fontsize{22}{24}\selectfont\color{popcream}\MakeUppercase{#1}\par}
  \end{tcolorbox}
}
% Section numérotée : pastille orange avec le numéro + titre Archivo
\newcounter{popsec}
\newcommand{\coursec}[1]{%
  \stepcounter{popsec}\setcounter{popsub}{0}%
  \par\vspace{16pt}\noindent
  \begin{minipage}{\linewidth}
  \tikz[baseline=-0.7ex]\node[draw=popink,line width=1.6pt,fill=poporange,rounded corners=4pt,minimum size=22pt,inner sep=0]{\popdisplay\fontsize{12}{12}\selectfont\color{popcream}\thepopsec};%
  \hspace{8pt}{\popdisplay\fontsize{15}{17}\selectfont\color{popink}\MakeUppercase{#1}}\par
  \vspace{4pt}\noindent{\color{poporange}\rule{36pt}{3.6pt}}
  \end{minipage}\par\nopagebreak\vspace{8pt}
}
\newcounter{popsub}
\newcommand{\courssub}[1]{%
  \stepcounter{popsub}%
  \par\vspace{9pt}\noindent{\popxbold\fontsize{11.5}{13}\selectfont\color{popink}{\color{poporange}\thepopsec.\thepopsub}\hspace{6pt}#1}\par\nopagebreak\vspace{4pt}
}

% ============================================================
% Boîtes pédagogiques — même recette : fond blanc, bordure épaisse colorée,
% ombre dure encre, badge plein. Titre optionnel [#1].
% ============================================================
\tcbset{popbase/.style={breakable,enhanced,colback=white,boxrule=2.3pt,arc=9pt,
  shadow={3pt}{-3pt}{0pt}{popink},left=14pt,right=14pt,top=11pt,bottom=11pt,before skip=11pt,after skip=15pt,
  fontupper=\popsemi\fontsize{10.6}{14.5}\selectfont\color{popink},
  fontlower=\popsemi\fontsize{10.6}{14.5}\selectfont\color{popink}}}
% Coche dessinée (le glyphe ✓ n'existe pas dans les polices du document)
\newcommand{\popcheck}[1]{\tikz[baseline=-0.5ex]\draw[#1,line width=1.6pt,line cap=round,line join=round](-0.13,0.0)--(-0.04,-0.09)--(0.13,0.1);}
\providecommand{\coche}{\popcheck{popgreen}}
\providecommand{\resultat}[1]{\par\smallskip\noindent\fcolorbox{popgreen}{popgreenL}{\parbox{\dimexpr\linewidth-2\fboxsep-2\fboxrule\relax}{\textbf{Résultat.} #1}}\par\smallskip}
\newcommand{\popboxhead}[3]{\poppill{#1}{#2}{white}\ifx\relax#3\relax\else\hspace{6pt}{\popxbold\fontsize{10.6}{12}\selectfont\color{popink}#3}\fi\par\vspace{7pt}}

\newtcolorbox{aretenir}[1][]{popbase,colframe=popgreen,before upper={\popboxhead{À retenir}{popgreen}{#1}}}
\newtcolorbox{exemplebox}[1][]{popbase,colframe=poporange,before upper={\popboxhead{Exemple}{poporange}{#1}}}
\newtcolorbox{definition}[1][]{popbase,colframe=popblue,before upper={\popboxhead{Définition}{popblue}{#1}}}
\newtcolorbox{propriete}[1][]{popbase,colframe=poppurple,before upper={\popboxhead{Propriété}{poppurple}{#1}}}
\newtcolorbox{methode}[1][]{popbase,colframe=popgold,before upper={\popboxhead{Méthode}{popgold}{#1}}}
\newtcolorbox{attention}[1][]{popbase,colframe=poppink,before upper={\popboxhead{Attention}{poppink}{#1}}}
\newtcolorbox{experience}[1][]{popbase,colframe=popblue,colback=popblueL,before upper={\popboxhead{Expérience}{popblue}{#1}}}
% « Le savais-tu ? » : carte violette pleine (contexte, culture, Cameroun)
\newtcolorbox{savaistu}[1][]{popbase,colback=poppurple,colframe=popink,
  fontupper=\popsemi\fontsize{10.4}{14.2}\selectfont\color{white},
  before upper={\poppill{Le savais-tu\,?}{white}{poppurple}\ifx\relax#1\relax\else\hspace{6pt}{\popxbold\color{white}#1}\fi\par\vspace{7pt}}}
% Exercice résolu : énoncé en haut, \tcblower puis la solution sur fond crème
\newcounter{popexo}\newcounter{popexr}
\newtcolorbox{exoresolu}[1][]{popbase,colframe=poporange,colbacklower=popcream,
  segmentation style={popink,line width=1.2pt,dashed},
  before upper={\stepcounter{popexr}\popboxhead{Exercice résolu \thepopexr}{poporange}{#1}},
  before lower={\poppill{Solution}{popink}{popcream}\par\vspace{6pt}}}
% Exercice (non résolu en place) — la correction est dans la partie Corrigés
\newtcolorbox{exercice}[1][]{popbase,colframe=popink,
  before upper={\stepcounter{popexo}\popboxhead{Exercice \thepopexo}{popink}{#1}}}
% Corrigé
\newtcolorbox{corrige}[1][]{popbase,colframe=popgreen,colback=popgreenL,
  before upper={\popboxhead{Corrigé}{popgreen}{#1}}}
% Objectifs / prérequis / bilan
\newtcolorbox{objectifs}{popbase,colframe=popink,colback=poporangeL,
  before upper={\popboxhead{Objectifs de la leçon}{popink}{}}}
\newtcolorbox{prerequis}{popbase,colframe=popink,colback=popblueL,
  before upper={\popboxhead{Prérequis}{popblue}{}}}
\newtcolorbox{bilan}{popbase,colframe=popink,colback=white,boxrule=2.8pt,
  before upper={\popboxhead{Fiche bilan}{poporange}{L'essentiel à connaître pour l'examen}}}
% Figure encadrée avec légende
\newtcolorbox{popfigure}[1][]{enhanced,breakable=false,colback=white,colframe=popline,boxrule=1.4pt,arc=8pt,
  left=10pt,right=10pt,top=10pt,bottom=8pt,before skip=10pt,after skip=12pt,halign=center,
  lower separated=false,halign lower=center,
  fontlower=\popsemi\fontsize{9}{11.5}\selectfont\color{popmuted},
  #1}
\newcommand{\legende}[1]{\par\vspace{6pt}{\popsemi\fontsize{9}{11.5}\selectfont\color{popmuted}#1}}

% ============================================================
% Signature OnBuch+
% ============================================================
\newcommand{\poplogoinline}[1]{{\poplogo\fontsize{#1}{#1}\selectfont\color{popink}OnBuch\color{poporange}+}}
\newcommand{\signatureonbuch}{%
  \begin{tcolorbox}[enhanced,colback=popink,colframe=popink,boxrule=2.2pt,arc=10pt,
      drop shadow=poporange,left=14pt,right=14pt,top=10pt,bottom=10pt,before skip=0pt,after skip=0pt]
    \tikz[baseline=-0.7ex]\node[circle,fill=popgreen,draw=popcream,line width=1.3pt,minimum size=22pt,inner sep=0]{\popcheck{white}};%
    \hspace{9pt}\begin{minipage}[c]{0.62\linewidth}
      {\popxbold\fontsize{12}{13}\selectfont\color{popcream}Signé\ \textbullet\ {\poplogo OnBuch\color{poporange}+}}\par\vspace{2pt}
      {\raggedright\popsemi\fontsize{8}{10.5}\selectfont\color{popline}\MakeUppercase{Équipe pédagogique OnBuch+}\\\MakeUppercase{Conforme au programme officiel MINESEC}\par}
    \end{minipage}\hfill
    {\poplogo\fontsize{22}{22}\selectfont\color{popcream}OnBuch\color{poporange}+}
  \end{tcolorbox}%
}

% ============================================================
% Page de garde
% #1 titre · #2 matière · #3 niveau · #4 module · #5 n° leçon · #6 durée ·
% #7 description / objectifs (texte ou itemize)
% ============================================================
\newcommand{\pagegarde}[7]{%
  \thispagestyle{empty}
  \newgeometry{a4paper,margin=1.7cm,top=1.6cm,bottom=1.4cm}%
  \begin{tikzpicture}[remember picture,overlay]
    \fill[poporange!18] ([xshift=-3.2cm,yshift=-3.6cm]current page.north east) circle (4.6cm);
    \fill[poppurple!14] ([xshift=2.4cm,yshift=7.6cm]current page.south west) circle (3.8cm);
  \end{tikzpicture}%
  {\poplogo\fontsize{30}{30}\selectfont\color{popink}OnBuch\color{poporange}+}\hfill
  \raisebox{0.6ex}{\poppill{Cours signé \textbullet\ Premium}{popink}{popcream}}\par
  \vspace{1pt}\popsquiggle{poppurple}\par
  \vspace{8pt}
  \begin{tcolorbox}[enhanced,colback=poporange,colframe=popink,boxrule=2.8pt,arc=13pt,
      drop shadow=popink,left=18pt,right=18pt,top=12pt,bottom=13pt,after skip=10pt]
    \poppill{#2 \textbullet\ #3}{popink}{popcream}\hspace{5pt}\poppill{#4}{white}{popink}\par\vspace{8pt}
    {\popdisplay\fontsize{14}{14}\selectfont\color{popink}LEÇON #5}\par\vspace{4pt}
    {\raggedright\hyphenpenalty=10000\exhyphenpenalty=10000\popdisplay\fontsize{25}{27}\selectfont\color{popcream}\MakeUppercase{#1}\par}
    \vspace{8pt}
    {\popxbold\fontsize{9.5}{9.5}\selectfont\color{popink}\MakeUppercase{Durée conseillée : #6}}
  \end{tcolorbox}
  \begin{tcolorbox}[enhanced,colback=white,colframe=popink,boxrule=2.2pt,arc=10pt,
      drop shadow=popink,left=16pt,right=16pt,top=10pt,bottom=10pt,after skip=8pt]
    \poppill{Dans cette leçon}{poppurple}{white}\par\vspace{5pt}
    \popsemi\fontsize{9.3}{12.6}\selectfont\color{popink}#7
  \end{tcolorbox}
  \noindent\begin{minipage}[t]{0.487\linewidth}
    \begin{tcolorbox}[enhanced,colback=popgreen,colframe=popink,boxrule=2.2pt,arc=10pt,
        drop shadow=popink,left=11pt,right=11pt,top=10pt,bottom=10pt,height=3.1cm,valign=center]
      \tcbox[colback=white,colframe=popink,boxrule=1.3pt,arc=5pt,boxsep=0pt,nobeforeafter,
        left=3pt,right=3pt,top=3pt,bottom=3pt,valign=center]{\qrcode[height=1.9cm]{https://play.google.com/store/apps/details?id=com.luvvix.onbuch&hl=fr}}%
      \hspace{8pt}\begin{minipage}[c]{0.5\linewidth}
        {\popdisplay\fontsize{11}{12.5}\selectfont\color{white}\MakeUppercase{Télécharge OnBuch}}\par\vspace{4pt}
        {\raggedright\popsemi\fontsize{8.4}{11}\selectfont\color{white}Gratuit \textbullet\ Google Play\\Tuteur IA, annales, concours, résultats\par}
      \end{minipage}
    \end{tcolorbox}
  \end{minipage}\hfill
  \begin{minipage}[t]{0.487\linewidth}
    \begin{tcolorbox}[enhanced,colback=poppurple,colframe=popink,boxrule=2.2pt,arc=10pt,
        drop shadow=popink,left=11pt,right=11pt,top=10pt,bottom=10pt,height=3.1cm,valign=center]
      \tikz[baseline=-0.7ex]\node[circle,fill=white,draw=popink,line width=1.3pt,minimum size=1.6cm,inner sep=0]{\popdisplay\fontsize{24}{24}\selectfont\color{poppurple}+};%
      \hspace{8pt}\begin{minipage}[c]{0.6\linewidth}
        {\popdisplay\fontsize{11}{12.5}\selectfont\color{white}\MakeUppercase{Passe à OnBuch+}}\par\vspace{4pt}
        {\raggedright\popsemi\fontsize{8.4}{11}\selectfont\color{white}Tous les cours signés, Tuteur IA illimité, fiches PDF — onglet OnBuch+ de l'app\par}
      \end{minipage}
    \end{tcolorbox}
  \end{minipage}
  \vspace{8pt}
  \popcontenu
  \vfill
  \signatureonbuch
  \restoregeometry
  \newpage
}

% Rangée « Ce cours contient » (page de garde)
\newcommand{\poptile}[3]{% #1 couleur #2 gros texte #3 légende
  \begin{tcolorbox}[enhanced,colback=#1,colframe=popink,boxrule=2pt,arc=9pt,shadow={2.5pt}{-2.5pt}{0pt}{popink},
      left=6pt,right=6pt,top=9pt,bottom=9pt,halign=center,valign=center,height=1.75cm,width=\linewidth,nobeforeafter]
    {\popdisplay\fontsize{12.5}{13}\selectfont\color{white}\MakeUppercase{#2}}\par\vspace{3pt}
    {\centering\popsemi\fontsize{7.8}{9.5}\selectfont\color{white}#3\par}
  \end{tcolorbox}}
\newcommand{\popcontenu}{%
  \par\noindent{\popxboldls\fontsize{7.5}{7.5}\selectfont\color{popmuted}CE COURS SIGNÉ CONTIENT}\par\vspace{7pt}
  \noindent\begin{minipage}[t]{0.235\linewidth}\poptile{popblue}{Cours}{complet, illustré, pas à pas}\end{minipage}\hfill
  \begin{minipage}[t]{0.235\linewidth}\poptile{poporange}{Méthodes}{et exercices résolus}\end{minipage}\hfill
  \begin{minipage}[t]{0.235\linewidth}\poptile{poppink}{Exercices}{type examen + corrigés}\end{minipage}\hfill
  \begin{minipage}[t]{0.235\linewidth}\poptile{popgreen}{Fiche bilan}{l'essentiel à retenir}\end{minipage}\par}

% Sommaire
\newtcolorbox{sommaire}{enhanced,colback=white,colframe=popink,boxrule=2.2pt,arc=10pt,
  drop shadow=popink,left=18pt,right=18pt,top=13pt,bottom=13pt,before skip=6pt,after skip=20pt,
  before upper={\poppill{Sommaire}{poppurple}{white}\par\vspace{9pt}\popsemi\fontsize{10.6}{14.5}\selectfont\color{popink}}}

% Signature de fin de cours — poussée tout en bas de la dernière page
\newcommand{\findecours}{%
  \par\vfill\nopagebreak
  \begin{center}\popsquiggle{poporange}\end{center}\vspace{4pt}
  \signatureonbuch
}
\AtBeginDocument{\def\llbracket{\mathopen{[\mkern-3.5mu[}}\def\rrbracket{\mathclose{]\mkern-3.5mu]}}} % intervalles d'entiers (glyphe ⟦ absent de la police)
% Glyphes absents de Fira Math : redéfinis à partir de glyphes disponibles
\AtBeginDocument{%
  \def\setminus{\mathbin{\backslash}}\let\smallsetminus\setminus
  \def\vdots{\mathord{\vbox{\baselineskip3.2pt\lineskiplimit0pt\kern4pt\hbox{.}\hbox{.}\hbox{.}}}}%
  \def\ddots{\mathinner{\mkern1mu\raise6.5pt\hbox{.}\mkern2mu\raise3.6pt\hbox{.}\mkern2mu\raise0.7pt\hbox{.}\mkern1mu}}%
  \def\triangle{\mathord{\scalebox{0.9}{$\symup{\Delta}$}}}%
  \def\nabla{\mathord{\raisebox{\depth}{\scalebox{1}[-1]{$\symup{\Delta}$}}}}%
  \def\checkmark{\popcheck{popdark}}%
  \def\star{\mathbin{\ast}}\def\bigstar{\mathord{\ast}}%
  \def\diamond{\mathbin{\scalebox{0.6}{$\Diamond$}}}%
  \def\bigcup{\mathop{\mathchoice{\raisebox{-0.3em}{\scalebox{1.7}{$\cup$}}}{\scalebox{1.25}{$\cup$}}{\cup}{\cup}}\displaylimits}%
  \def\bigcap{\mathop{\mathchoice{\raisebox{-0.3em}{\scalebox{1.7}{$\cap$}}}{\scalebox{1.25}{$\cap$}}{\cap}{\cap}}\displaylimits}%
  \def\lesseqqgtr{\mathrel{\lessgtr}}\def\bigoplus{\mathop{\oplus}\displaylimits}%
}
\providecommand{\ssi}{si et seulement si} % abréviation fréquente
\AtBeginDocument{\providecommand{\wideparen}{\overparen}} % arc de cercle

\begin{document}

\pagegarde{IV : rédiger un résumé de texte}{Français}{Tle Tech}{Français – classe de Terminale technique}{N16}{10 séances (fiche de progression harmonisée)}{Cette leçon outille l'élève de Terminale technique pour rédiger un résumé de texte : repérage des idées essentielles, reformulation, rédaction et amélioration. Elle s'appuie sur l'étude et le bilan de l'œuvre Le Vieux nègre et la médaille et mobilise les notions d'énoncé constatif et performatif, de lexique spécialisé et de tonalités satirique et polémique.
\vspace{4pt}
\begin{multicols}{2}\small
\begin{itemize}
\item À la fin de cette leçon, je sais définir le résumé et distinguer résumé, contraction de texte et compte rendu.
\item Je sais repérer et hiérarchiser les idées essentielles d'un texte en éliminant exemples, répétitions et digressions.
\item Je sais reformuler les idées essentielles avec mes propres mots en respectant la pensée de l'auteur.
\item Je sais distinguer un énoncé constatif d'un énoncé performatif et en tenir compte dans la reformulation.
\item Je sais identifier et utiliser le lexique spécialisé d'un texte sans le dénaturer.
\item Je sais reconnaître les tonalités satirique et polémique d'un texte et les rendre dans mon résumé.
\end{itemize}\end{multicols}}

\begin{sommaire}
\begin{multicols}{2}
\begin{enumerate}
\item Le résumé : définition, enjeux et cadre de l'épreuve
\item Lecture analytique : repérer les idées essentielles (à partir de Le Vieux nègre et la médaille)
\item Outils de langue 1 : énoncé constatif et performatif, lexique spécialisé
\item Outils de langue 2 : tonalités satirique et polémique
\item Rédaction du résumé : reformulation et production
\item Amélioration du résumé et bilan de l'œuvre
\item Activité d'intégration
\item Méthodes et exercices résolus
\item Exercices
\item Corrigés des exercices
\item Fiche bilan
\end{enumerate}
\end{multicols}
\end{sommaire}

\begin{objectifs}
\begin{itemize}
\item À la fin de cette leçon, je sais définir le résumé et distinguer résumé, contraction de texte et compte rendu.
\item Je sais repérer et hiérarchiser les idées essentielles d'un texte en éliminant exemples, répétitions et digressions.
\item Je sais reformuler les idées essentielles avec mes propres mots en respectant la pensée de l'auteur.
\item Je sais distinguer un énoncé constatif d'un énoncé performatif et en tenir compte dans la reformulation.
\item Je sais identifier et utiliser le lexique spécialisé d'un texte sans le dénaturer.
\item Je sais reconnaître les tonalités satirique et polémique d'un texte et les rendre dans mon résumé.
\item Je sais rédiger un résumé continu, cohérent et correct, au quart environ de la longueur du texte.
\item Je sais relire, corriger et améliorer mon résumé (orthographe, connecteurs, fidélité).
\item Je sais situer Le Vieux nègre et la médaille dans son contexte colonial et en dégager la signification pour le bilan de l'œuvre.
\end{itemize}
\end{objectifs}

\begin{prerequis}
\begin{itemize}
\item Lecture et compréhension globale d'un texte narratif et argumentatif (classes de Première)
\item Notions de thème, d'idée principale et de plan d'un texte
\item Connecteurs logiques et marqueurs de relation (étude de Première)
\item Connaissance générale de l'œuvre Le Vieux nègre et la médaille (lecture cursive)
\item Grammaire de la phrase : types et formes de phrases
\end{itemize}
\end{prerequis}

\newpage

\coursec{Le résumé : définition, enjeux et cadre de l'épreuve}

Au marché Mokolo de Yaoundé, un élève de Terminale doit expliquer en trois minutes à son oncle pressé le contenu d'un article de deux pages paru dans \emph{Cameroon Tribune} sur la grève des transporteurs. Pas question de tout lire, pas question de tout raconter : il faut aller droit à l'essentiel, sans rien inventer, sans rien déformer. À l'école, au bureau, à l'administration, celui qui sait dire l'essentiel d'un long texte sans le trahir gagne du temps et du crédit. Comment réduire un texte au quart de sa longueur sans trahir sa pensée ? Quelles étapes garantissent un résumé fidèle, clair et personnel ? C'est tout l'art de la \cle{contraction de texte}, épreuve incontournable du Baccalauréat technique, que cette leçon va nous apprendre, en nous appuyant sur \emph{Le Vieux nègre et la médaille} de Ferdinand Oyono.

\courssub{Définition de la contraction de texte}

\textbf{Intuition.} Imaginez une calebasse pleine de jus de foléré : si l'on fait bouillir ce jus, l'eau s'évapore et il reste un sirop concentré, plus petit en volume mais qui garde tout le goût du fruit. La contraction de texte fonctionne exactement ainsi : on fait « évaporer » les détails, les exemples, les répétitions, et il reste la pensée concentrée de l'auteur.

\begin{definition}[La contraction de texte]
La \cle{contraction de texte} est une opération d'écriture qui consiste à réduire un texte à une longueur imposée (généralement le \cle{quart} de sa longueur initiale) en en conservant les \cle{idées essentielles}, sans rien ajouter, sans rien déformer et sans donner son avis. Le texte obtenu s'appelle le \cle{résumé}.
\end{definition}

Trois exigences découlent de cette définition :

\begin{itemize}
\item \textbf{Réduire} : le résumé respecte un nombre de mots imposé. Ce n'est pas une paraphrase (qui dirait la même chose en aussi long), c'est une compression.
\item \textbf{Conserver} : les idées principales, la logique du raisonnement et le point de vue de l'auteur doivent subsister intacts.
\item \textbf{Reformuler} : le résumé est rédigé avec les mots du résumant, pas avec les phrases du texte. C'est une réécriture personnelle.
\end{itemize}

\begin{popfigure}
\begin{center}
\begin{tikzpicture}[font=\small]
% Entonnoir
\fill[popblue!25] (-4.5,3) -- (4.5,3) -- (1.8,1.2) -- (-1.8,1.2) -- cycle;
\fill[poporange!30] (-1.8,1.2) -- (1.8,1.2) -- (0.9,0) -- (-0.9,0) -- cycle;
\fill[popgreen!35] (-0.9,0) -- (0.9,0) -- (0.9,-1.6) -- (-0.9,-1.6) -- cycle;
\draw[popdark,thick] (-4.5,3) -- (4.5,3) -- (1.8,1.2) -- (-1.8,1.2) -- cycle;
\draw[popdark,thick] (-1.8,1.2) -- (1.8,1.2) -- (0.9,0) -- (-0.9,0) -- cycle;
\draw[popdark,thick] (-0.9,0) -- (0.9,0) -- (0.9,-1.6) -- (-0.9,-1.6) -- cycle;
\node[align=center] at (0,2.1) {\textbf{Texte de départ}\\ 400 mots (idées + exemples + détails)};
\node[align=center] at (0,0.6) {\textbf{Idées}\\ \textbf{essentielles}};
\node[align=center,text=white] at (0,-0.8) {\textbf{Résumé : 100 mots}};
\draw[-{Stealth},very thick,popink] (5.2,3) -- (5.2,-1.6) node[midway,right,align=left]{on élimine\\ les détails,\\ on garde le sens};
\end{tikzpicture}
\end{center}
\legende{Schéma en entonnoir de la contraction de texte : d'un texte de 400 mots, on extrait les idées essentielles pour produire un résumé de 100 mots (le quart).}
\end{popfigure}

\courssub{Résumé, compte rendu, synthèse, commentaire : ne pas confondre}

Au Baccalauréat technique comme dans la vie professionnelle, plusieurs exercices ressemblent au résumé sans s'y confondre. Les distinguer, c'est déjà éviter la moitié des erreurs.

\begin{definition}[Quatre exercices voisins]
\begin{itemize}
\item Le \cle{résumé} : réduction fidèle d'\textbf{un seul} texte, sans aucun apport personnel.
\item Le \cle{compte rendu} : présentation objective et structurée d'un texte, d'une réunion ou d'un événement ; il peut être aussi long que nécessaire et signale la structure du document (il annonce, il explique, il conclut).
\item La \cle{synthèse} : construction d'un texte bref et cohérent à partir de \textbf{plusieurs} documents ; elle organise et confronte les idées.
\item Le \cle{commentaire} : étude approfondie d'un texte qui \textbf{juge}, analyse et interprète ; l'opinion argumentée du lecteur y a toute sa place.
\end{itemize}
\end{definition}

\begin{center}
\begin{tabularx}{\linewidth}{|l|X|X|X|X|}
\hline
\rowcolor{popblueL} \textbf{Critère} & \textbf{Résumé} & \textbf{Compte rendu} & \textbf{Synthèse} & \textbf{Commentaire} \\
\hline
Nombre de textes & Un seul & Un seul (ou un événement) & Plusieurs & Un seul \\
\hline
Fidélité au sens & Totale & Totale & Totale & Interprétation possible \\
\hline
Opinion personnelle & Interdite & Interdite & Interdite & Attendue et argumentée \\
\hline
Longueur & Imposée (le quart) & Libre ou structurée & Imposée & Libre \\
\hline
Formulation & Personnelle & Personnelle & Personnelle & Personnelle \\
\hline
\end{tabularx}
\end{center}

\begin{attention}[Le piège du commentaire déguisé]
Beaucoup de candidats transforment leur résumé en commentaire : « l'auteur a raison de dire que... », « cette idée est très juste... ». Toute trace de jugement personnel fait basculer la copie hors de l'exercice demandé. Dans un résumé, on ne dit pas ce que l'on pense du texte : on dit ce que le texte pense.
\end{attention}

\courssub{Les qualités d'un bon résumé}

Un bon résumé se reconnaît à cinq qualités, que l'examinateur vérifie une à une dans votre copie.

\begin{aretenir}[Les cinq qualités du résumé]
\begin{enumerate}
\item \textbf{Fidélité} : le résumé respecte la pensée de l'auteur, son point de vue et l'ordre de ses idées. Aucune idée étrangère au texte ne s'y glisse.
\item \textbf{Brièveté} : il respecte le nombre de mots imposé, dans la marge de tolérance autorisée.
\item \textbf{Clarté} : il se lit facilement, en phrases simples et correctes ; un lecteur qui n'a pas lu le texte d'origine le comprend seul.
\item \textbf{Cohérence} : les idées s'enchaînent logiquement grâce à des connecteurs (d'abord, ensuite, en effet, cependant, enfin).
\item \textbf{Formulation personnelle} : il est écrit avec les mots du candidat, sans recopier les phrases du texte.
\end{enumerate}
\end{aretenir}

\begin{methode}[Les quatre étapes du résumé]
\begin{enumerate}
\item \textbf{Lire et comprendre} : lire le texte deux ou trois fois, dégager le thème et la thèse de l'auteur, vérifier le sens des mots difficiles.
\item \textbf{Repérer les idées essentielles} : souligner les idées principales (souvent une par paragraphe), éliminer exemples, répétitions, citations et détails.
\item \textbf{Reformuler} : exprimer chaque idée avec ses propres mots, en regroupant ce qui peut l'être, sans rien ajouter.
\item \textbf{Rédiger et vérifier} : rédiger un texte suivi et cohérent, compter les mots, indiquer leur nombre à la fin, relire l'orthographe.
\end{enumerate}
\end{methode}

\courssub{Le cadre de l'épreuve au Baccalauréat technique}

À l'examen, la contraction de texte obéit à des règles précises qu'il faut connaître avant d'entrer dans la salle.

\begin{propriete}[Règles de l'épreuve]
\begin{itemize}
\item Le sujet impose une longueur : le plus souvent \textbf{le quart} du texte, parfois un nombre de mots précis.
\item Une \textbf{marge de tolérance de $\pm 10\,\%$} est admise autour du nombre de mots exigé.
\item Le candidat \textbf{indique le nombre de mots} de son résumé à la fin de sa copie (par exemple : « 118 mots »).
\item Le résumé est rédigé en \textbf{texte suivi} (pas de plan apparent, pas de titres, pas de puces).
\end{itemize}
\end{propriete}

\begin{exemplebox}[Calcul de la longueur du résumé]
Un texte de 480 mots doit être résumé au quart.
\begin{itemize}
\item Longueur visée : $480 \div 4 = 120$ mots.
\item Tolérance : $120 \times 0{,}10 = 12$ mots.
\item Fourchette acceptée : de $120 - 12 = 108$ mots à $120 + 12 = 132$ mots.
\end{itemize}
Le candidat écrira donc un résumé compris entre 108 et 132 mots et notera le nombre exact obtenu à la fin.
\end{exemplebox}

\begin{attention}[Les erreurs rédhibitoires]
Certaines fautes coûtent très cher à l'examen :
\begin{itemize}
\item \textbf{Recopier des phrases entières du texte} : ce n'est pas un résumé, c'est un collage ; l'examinateur le sanctionne lourdement, parfois par la note de zéro à l'exercice.
\item \textbf{Donner son avis} (« je pense que », « l'auteur a tort ») : le résumé doit rester neutre.
\item \textbf{Ignorer la limite} : un résumé très en dessous ou très au-dessus de la fourchette prouve que l'on n'a pas su sélectionner l'essentiel.
\item \textbf{Oublier d'indiquer le nombre de mots} à la fin de la copie.
\item \textbf{Présenter le résumé en plan ou en liste} : il doit être rédigé en phrases liées.
\end{itemize}
\end{attention}

\courssub{Le résumé dans la vie courante : une compétence professionnelle}

La contraction de texte n'est pas un exercice artificiel inventé pour l'école : c'est un geste quotidien de la vie administrative et professionnelle camerounaise.

\begin{itemize}
\item Le secrétaire d'un \cle{COGES} (Comité de gestion des établissements scolaires) résume en une page un procès-verbal de réunion de trois heures pour l'affichage des parents d'élèves.
\item Un employé de mairie résume une \cle{note administrative} ministérielle de dix pages pour son chef de service qui doit décider en cinq minutes.
\item Un journaliste de \emph{Cameroon Tribune} ou de \emph{Mutations} condense un long rapport en un article de quelques lignes.
\item Un technicien résume le compte rendu d'une panne pour son responsable.
\end{itemize}

Dans tous ces cas, la règle est la même : dire l'essentiel, fidèlement, brièvement, avec ses propres mots.

\begin{exoresolu}[Repérer un vrai résumé]
\textbf{Énoncé.} Voici un extrait de texte suivi de deux propositions de résumé. Laquelle est acceptable à l'examen ? Justifiez.

\emph{Texte (extrait adapté du contexte du Vieux nègre et la médaille) : « Meka a reçu une médaille pour avoir donné ses terres à la mission et ses fils à la guerre. Mais cette médaille, loin de le récompenser vraiment, révèle l'injustice de la colonisation : on honore l'homme en apparence pendant qu'on le dépouille en réalité. » (48 mots)}

\emph{Proposition A : « Je trouve que l'auteur a parfaitement raison : la médaille de Meka, qui a donné ses terres à la mission et ses fils à la guerre, révèle l'injustice de la colonisation, car on honore l'homme en apparence pendant qu'on le dépouille en réalité. »}

\emph{Proposition B : « La médaille de Meka, censée récompenser ses sacrifices, dénonce en fait l'hypocrisie coloniale, qui honore l'homme tout en le spoliant. » (20 mots)}
\tcblower
\textbf{Solution.} La proposition B est le bon résumé. Elle est brève (proche du quart du texte), fidèle à la pensée de l'auteur, et reformulée avec des mots personnels (« hypocrisie coloniale », « spoliant »). La proposition A est fautive pour trois raisons : elle donne un avis personnel (« je trouve que l'auteur a parfaitement raison »), elle recopie mot pour mot de longs passages du texte (« donné ses terres à la mission et ses fils à la guerre », « on honore l'homme en apparence pendant qu'on le dépouille en réalité »), et elle est presque aussi longue que l'original. C'est un collage commenté, pas un résumé.
\end{exoresolu}

\begin{savaistu}[La contraction de texte au Baccalauréat technique]
Au Baccalauréat technique, la contraction de texte porte souvent sur un extrait d'une œuvre au programme — comme \emph{Le Vieux nègre et la médaille} de Ferdinand Oyono, roman camerounais publié en 1956 — ou sur un texte de presse. Elle vaut une part importante de la note de production écrite. Les jurys rapportent chaque année que les candidats qui réussissent sont ceux qui ont pris le temps de \textbf{lire deux fois} le texte avant d'écrire : dix minutes de lecture attentive valent mieux qu'une heure de rédaction précipitée.
\end{savaistu}

\begin{exercice}[À vous de jouer]
Un article de presse compte 360 mots. Le sujet demande de le résumer au quart.
\begin{enumerate}
\item Calculez la longueur visée du résumé.
\item Calculez la fourchette acceptée avec la tolérance de $\pm 10\,\%$.
\item Un candidat rend un résumé de 75 mots ; un autre en rend un de 98 mots. Lequel respecte la consigne ?
\end{enumerate}
\end{exercice}

\begin{corrige}[Exercice : À vous de jouer]
\begin{enumerate}
\item Longueur visée : $360 \div 4 = 90$ mots.
\item Tolérance : $90 \times 0{,}10 = 9$ mots ; fourchette acceptée : de $90 - 9 = 81$ mots à $90 + 9 = 99$ mots.
\item Le résumé de 75 mots est trop court (inférieur à 81) : le candidat a probablement éliminé des idées essentielles. Le résumé de 98 mots est dans la fourchette : c'est lui qui respecte la consigne.
\end{enumerate}
\end{corrige}

\begin{aretenir}[L'essentiel de la section]
La contraction de texte réduit un texte à une longueur imposée (souvent le quart) en conservant ses idées essentielles, sans rien ajouter ni déformer. Le résumé se distingue du compte rendu, de la synthèse et du commentaire par l'interdiction absolue de toute opinion personnelle. Un bon résumé est fidèle, bref, clair, cohérent et formulé avec des mots personnels. À l'examen, on respecte la marge de $\pm 10\,\%$, on indique le nombre de mots et l'on ne recopie jamais les phrases du texte. Les quatre étapes — lire, repérer, reformuler, rédiger — seront détaillées dans les sections suivantes.
\end{aretenir}

\coursec{Lecture analytique : repérer les idées essentielles (à partir de Le Vieux nègre et la médaille)}

Après avoir défini le résumé comme un exercice de condensation rigoureuse qui exige fidélité et concision, nous devons maintenant acquérir la compétence préalable indispensable : la lecture analytique. Résumer, ce n'est pas « couper » au hasard dans le texte, c'est \textbf{reconstruire} la pensée de l'auteur en ne gardant que sa charpente. Cette section vous apprend à disséquer un texte — ici l'Exposé N°1 de \emph{Le Vieux nègre et la médaille} de Ferdinand Oyono — pour en extraire la substantifique moelle.

\courssub{Première lecture globale : dégager le thème, la thèse et l'intention de l'auteur}

La première lecture ne doit pas être une lecture « au mot près ». C'est une lecture \textbf{exploratoire} : on laisse le texte venir à soi pour en saisir l'unité profonde. On cherche à répondre à trois questions fondamentales : \textbf{De quoi parle le texte ?} (le thème), \textbf{Qu'en dit l'auteur ?} (la thèse ou l'intention narrative), \textbf{Comment le dit-il ?} (le ton, le genre).

\begin{methode}[Lecture globale efficace]
\begin{enumerate}
\item \textbf{Lisez d'une traite}, sans stylo, sans dictionnaire. Acceptez de ne pas tout comprendre du premier coup.
\item \textbf{Fermez le texte} et formulez oralement ou par écrit, en \textbf{une seule phrase}, ce que vous avez retenu : « Ce texte raconte / explique / montre que... »
\item \textbf{Identifiez le genre} (roman, essai, article, discours...) et le \textbf{ton} (ironique, dramatique, neutre, polémique...). Ici, nous sommes dans un roman à tonalité satirique.
\item \textbf{Repérez les indices paratextuels} : titre de l'œuvre, titre de l'extrait (Exposé N°1), noms propres récurrents (Meka, le Commandant, le Prêtre), champs lexicaux.
\end{enumerate}
\end{methode}

\begin{exemplebox}[Application à l'Exposé N°1 de \emph{Le Vieux nègre et la médaille}]
\emph{Le Vieux nègre et la médaille} s'ouvre sur la préparation de la cérémonie de remise de la médaille à Meka, un vieux chef de village.
\begin{itemize}
\item \textbf{Thème :} La veille et le déroulement d'une cérémonie officielle coloniale (remise de médaille à un « loyal » serviteur).
\item \textbf{Thèse / Intention :} Dénoncer, par l'ironie et le regard interne du personnage, l'absurdité et l'humiliation masquée de la collaboration coloniale. La médaille n'est pas une récompense, c'est un symbole de servitude.
\item \textbf{Ton :} Satirique, distancié, parfois tendre sur le personnage de Meka (sa naïveté), féroce sur le système colonial.
\end{itemize}
\end{exemplebox}

\begin{savaistu}[Contexte historique et portée de l'œuvre]
Ferdinand Oyono publie \emph{Le Vieux nègre et la médaille} en \textbf{1956}, soit quatre ans avant l'indépendance du Cameroun (1er janvier 1960). L'œuvre fait scandale : elle ose montrer le colonisé non pas en victime passive, mais en acteur complexe, parfois ridicule, du système colonial. L'Exposé N°1 plante le décor : la naïveté de Meka qui croit à la reconnaissance de la France, alors que le lecteur, lui, perçoit l'ironie cruelle de la situation. C'est ce décalage (ironie dramatique) qui structure tout le texte.
\end{savaistu}

\courssub{Deuxième lecture analytique : repérer les idées principales et la structure du texte}

Munis de la compréhension globale, nous reprenons le texte, crayon en main. C'est la lecture \textbf{analytique} ou \textbf{microscopique}. Objectif : segmenter le texte en \textbf{unités de sens} (paragraphes ou groupes de paragraphes) et identifier l'\textbf{idée force} de chaque unité.

\begin{methode}[Repérer l'idée essentielle d'un paragraphe]
\begin{enumerate}
\item \textbf{Isolez l'unité de sens} : un paragraphe, ou un groupe de paragraphes traitant du même aspect.
\item \textbf{Cherchez la \cle{phrase-topic} (phrase nucléaire)} : c'est la phrase qui résume le paragraphe, souvent en début ou en fin de paragraphe. Elle est générale, les autres phrases la développent (exemples, explications).
\item \textbf{Si elle n'existe pas explicitement}, \textbf{formulez-la vous-même} en une phrase simple qui englobe tout le paragraphe sans en reprendre les détails.
\item \textbf{Testez} : si vous supprimez les autres phrases, l'idée reste-t-elle compréhensible ? Si oui, c'est l'essentiel.
\end{enumerate}
\end{methode}

\begin{experience}[Analyse structurée de l'Exposé N°1 (simulation sur extrait représentatif)]
Imaginons l'extrait découpé en quatre séquences narratives (correspondant aux grandes étapes de l'Exposé N°1) :
\begin{enumerate}
\item \textbf{Séquence 1 (Préparation) :} Meka, vêtu de son meilleur pagne, attend l'arrivée des Blancs. Il astique sa médaille. Le village est en effervescence.
\item \textbf{Séquence 2 (Arrivée des officiels) :} Le Commandant, le Prêtre, l'Administrateur arrivent. Description de leur attitude hautaine, de leur confort face à la misère du village.
\item \textbf{Séquence 3 (La cérémonie) :} Discours du Commandant (louanges creuses), remise de la médaille. Meka est ému, fier, il ne voit pas l'ironie.
\item \textbf{Séquence 4 (Le retour / Réaction villageoise) :} Meka rentre chez lui. Les villageois, loin de l'envier, se moquent doucement ou compatissent. La médaille est posée sur une chaise, inutile.
\end{enumerate}
\end{experience}

Voici la représentation hiérarchique de cette structure. L'arbre montre comment le thème central se ramifie en idées principales, et comment les détails (exemples, descriptions) s'y rattachent pour être ensuite éliminés.

\begin{popfigure}
\begin{tikzpicture}[
  mindmap,
  concept color=popblue!80,
  level 1/.style={sibling angle=120, level distance=4.5cm},
  level 2/.style={sibling angle=45, level distance=3.5cm},
  level 3/.style={sibling angle=30, level distance=3cm},
  every node/.style={concept, execute at begin node=\hskip0pt},
  font=\small,
  text=white,
  edge from parent/.style={draw=popline, thick}
]
\node[concept, scale=0.9, text width=3.5cm, align=center] {Thème central : \textbf{La cérémonie de la médaille}\\(Ironie coloniale)}
  child[concept color=popgreen!70] {node[concept, scale=0.8, text width=2.8cm, align=center] {Idée 1 : \textbf{Attente naïve de Meka}\\ (Préparation, fierté, propreté)}
    child[concept color=popgreen!40] {node[concept, scale=0.65, text width=2cm, align=center] {Détail : Pagne neuf\\ (À éliminer)}}
    child[concept color=popgreen!40] {node[concept, scale=0.65, text width=2cm, align=center] {Détail : Médaille astiquée\\ (À éliminer)}}
    child[concept color=popgreen!40] {node[concept, scale=0.65, text width=2cm, align=center] {Exemple : Effervescence village\\ (À éliminer)}}}
  child[concept color=poporange!70] {node[concept, scale=0.8, text width=2.8cm, align=center] {Idée 2 : \textbf{Arrivée méprisante des Blancs}\\ (Contraste confort / misère)}
    child[concept color=poporange!40] {node[concept, scale=0.65, text width=2cm, align=center] {Détail : Voiture climatisée\\ (À éliminer)}}
    child[concept color=poporange!40] {node[concept, scale=0.65, text width=2cm, align=center] {Exemple : Attitude Commandant\\ (À éliminer)}}
    child[concept color=poporange!40] {node[concept, scale=0.65, text width=2cm, align=center] {Description : Prêtre indifférent\\ (À éliminer)}}}
  child[concept color=poppurple!70] {node[concept, scale=0.8, text width=2.8cm, align=center] {Idée 3 : \textbf{Cérémonie creuse et humiliation finale}\\ (Discours, remise, ironie villageoise)}
    child[concept color=poppurple!40] {node[concept, scale=0.65, text width=2cm, align=center] {Exemple : Discours louanges\\ (À éliminer)}}
    child[concept color=poppurple!40] {node[concept, scale=0.65, text width=2cm, align=center] {Détail : Médaille sur chaise\\ (À éliminer)}}
    child[concept color=poppurple!40] {node[concept, scale=0.65, text width=2cm, align=center] {Paraphrase : Moqueries villageois\\ (À éliminer)}}};
\end{tikzpicture}
\legende{Arbre hiérarchique : du thème central aux idées principales (niveau 2) et aux éléments secondaires à éliminer (niveau 3, biffés conceptuellement).}
\end{popfigure}

\courssub{Éliminer le secondaire : exemples, répétitions, paraphrases, détails descriptifs, digressions}

C'est l'étape du \textbf{triage}. Tout ce qui ne porte pas la pensée nouvelle, l'avancée du récit ou de l'argumentation, doit être mis de côté. C'est la condition \textit{sine qua non} de la concision.

\begin{definition}[Secondaire vs Essentiel]
\begin{itemize}
\item \textbf{L'essentiel} : Idées directrices, arguments majeurs, tournants narratifs, thèse, conclusions. \emph{Indispensable à la compréhension du sens global.}
\item \textbf{Le secondaire} :
  \begin{itemize}
  \item \textbf{Exemples / Illustrations} : « Meka astique sa médaille avec un chiffon » (illustre sa fierté/soin).
  \item \textbf{Répétitions / Reformulations} : Dire deux fois que le Commandant est hautain.
  \item \textbf{Paraphrases} : Reformuler une idée déjà exprimée.
  \item \textbf{Détails descriptifs} : Couleur du pagne, marque de la voiture, chaleur du jour.
  \item \textbf{Digressions} : Apartés, souvenirs, descriptions d'atmosphère qui n'avancent pas l'intrigue principale.
  \end{itemize}
\end{itemize}
\end{definition}

\begin{attention}[Piège majeur : Confondre exemple frappant et idée essentielle]
L'image de la \textbf{médaille posée sur la chaise} à la fin de l'extrait est forte, visuelle, symbolique. L'élève a tendance à la garder dans son résumé : « Meka pose sa médaille sur une chaise. » \textbf{Erreur.} C'est un \textbf{détail illustratif} (un exemple) de l'idée essentielle : « La médaille, symbole de la collaboration, se révèle dénuée de valeur réelle et devient un objet encombrant. » Dans le résumé, on garde l'\textbf{idée} (l'inutilité/humiliation), pas l'\textbf{image} (la chaise), sauf si l'image \emph{est} l'idée (ex: description d'un tableau). Ici, l'image \emph{sert} l'idée.
\end{attention}

Appliquons ce tri sur notre extrait simulé d'Oyono.

\begin{popfigure}
\begin{center}
\begin{tabularx}{\linewidth}{|>{\columncolor{popblueL}}l|X|X|}
\hline
\rowcolor{popblueL}
\textbf{\color{white}Unité de texte (Séquence)} & \textbf{\color{white}À GARDER (Idées essentielles)} & \textbf{\color{white}À ÉLIMINER (Secondaire)} \\
\hline
\textbf{Séquence 1 : Préparation de Meka} & Meka prépare fièrement la cérémonie officielle où il doit recevoir une médaille de la part de l'administration coloniale, convaincu de la légitimité de cette distinction. & Description du pagne neuf, astiquage de la médaille, effervescence villageoise, détails météo. \\
\hline
\textbf{Séquence 2 : Arrivée des officiels} & Contraste saisissant entre le confort arrogant des colons et la misère du village ; le pouvoir colonial s'affiche. & Marque de la voiture, climatisation, description physique du Commandant/Prêtre, attitude du chauffeur. \\
\hline
\textbf{Séquence 3 : La cérémonie} & Discours hypocrite du Commandant ; remise de la médaille ; Meka, naïf, croit à l'honneur. & Contenu détaillé du discours (louanges creuses), gestuelle précise de la remise, émotion physique de Meka (larmes). \\
\hline
\textbf{Séquence 4 : Retour et réaction} & La médaille se révèle vaine (posée sur une chaise) ; le village perçoit l'ironie (moqueries/compassion) ; l'honneur est une humiliation. & Dialogues précis des villageois, position exacte de la chaise, pensées intérieures détaillées de Meka, description du coucher de soleil. \\
\hline
\end{tabularx}
\end{center}
\legende{Tableau de tri : à gauche l'unité de sens, au centre l'idée essentielle reformulée (noyau dur), à droite les éléments secondaires (exemples, détails, descriptions) à exclure du résumé.}
\end{popfigure}

\courssub{Hiérarchiser les idées retenues : respecter l'ordre chronologique ou logique}

Une fois le tri fait, les idées essentielles ne forment pas encore un résumé. Elles sont des \textbf{briques} qu'il faut assembler dans le bon ordre. Deux logiques s'imposent selon le type de texte :
\begin{itemize}
\item \textbf{Texte narratif (comme ici) :} Ordre \textbf{chronologique} impératif. On ne peut pas mettre la cérémonie avant l'arrivée des officiels. La causalité narrative (cause $\rightarrow$ effet) dicte la séquence.
\item \textbf{Texte argumentatif/explicatif :} Ordre \textbf{logique} (thèse $\rightarrow$ arguments $\rightarrow$ conclusion ; ou problème $\rightarrow$ causes $\rightarrow$ solutions).
\end{itemize}

\begin{aretenir}[Règle d'or pour le résumé narratif]
Le résumé suit \textbf{fidèlement la chronologie du texte source}. Toute inversion (anachronie) est une erreur de compréhension sanctionnée. Les connecteurs temporels (\emph{ensuite, puis, après que, lorsque, finalement}) sont les coutures qui relient les idées entre elles.
\end{aretenir}

\courssub{Construire le plan-squelette du résumé : une idée par étape, formulée en phrase repère}

Le plan-squelette (ou \textbf{plan détaillé}) est l'ossature du futur résumé. Chaque idée essentielle devient un \textbf{paragraphe} (ou une phrase complexe) du résumé. On formule chaque étape en une \textbf{phrase repère} complète, autonome, sans pronoms vagues renvoyant au texte original.

\begin{methode}[Construction du plan-squelette (5 étapes)]
\begin{enumerate}
\item \textbf{Numérotez} les idées essentielles dans l'ordre chronologique/logique.
\item \textbf{Reformulez} chaque idée en une \textbf{phrase complète} (sujet + verbe + complément), au \textbf{présent de narration} (ou passé composé si récit au passé, mais le présent de narration est la norme au Probatoire/Baccalauréat technique pour le résumé).
\item \textbf{Supprimez} tout « Il y a », « Le texte dit que », « L'auteur montre que ». Parlez \textbf{du contenu}, pas du texte. \emph{Ex: « Meka attend... » et non « Le texte raconte que Meka attend... »}
\item \textbf{Reliez} les phrases repères par des connecteurs logiques/temporels (\emph{Dans un premier temps ; Alors que ; Cependant ; Finalement}).
\item \textbf{Vérifiez} l'enchaînement : la lecture des seules phrases repères doit former un mini-récit cohérent.
\end{enumerate}
\end{methode}

\begin{exoresolu}[Construction du plan-squelette sur l'Exposé N°1]
\textbf{Énoncé :} À partir du tableau de tri ci-dessus, établir le plan-squelette (phrases repères) du résumé de l'Exposé N°1 (environ 4 phrases).

\textbf{Solution détaillée :}
\begin{enumerate}
\item \textbf{Phrase 1 (Séquence 1) :} \emph{Meka, vieux chef de village, prépare fièrement la cérémonie officielle où il doit recevoir une médaille de la part de l'administration coloniale, convaincu de la légitimité de cette distinction.}
\item \textbf{Phrase 2 (Séquence 2) :} \emph{L'arrivée des autorités françaises — Commandant, Prêtre, Administrateur — dans le confort et l'arrogance contraste violemment avec la misère et l'attente du village.}
\item \textbf{Phrase 3 (Séquence 3) :} \emph{Lors de la cérémonie, le Commandant prononce un discours hypocrite et remet la médaille à Meka, qui, naïf, y voit la consécration de sa fidélité.}
\item \textbf{Phrase 4 (Séquence 4) :} \emph{Mais de retour chez lui, Meka constate l'inutilité de la médaille, reléguée sur une chaise, tandis que les villageois, lucides, accueillent l'événement par l'ironie, révélant l'humiliation masquée de cette « récompense ».}
\end{enumerate}
\textbf{Connecteurs utilisés :} \emph{Lors de la cérémonie (temporel) ; Mais (opposition/adversatif) ; Tandis que (simultané/opposition).}
\textbf{Vérification :} La lecture de ces 4 phrases raconte l'histoire complète, sans exemple, sans détail, dans l'ordre. C'est le squelette parfait.
\end{exoresolu}

\begin{exemplebox}[L'idée essentielle globale (synthèse)]
Dans l'extrait étudié, l'idée essentielle qui traverse toutes les séquences est : \textbf{« La médaille, censée honorer Meka pour sa loyauté, révèle en réalité l'humiliation et l'ingratitude du système colonial, perçue lucidement par le village mais pas par le vieil homme. »} C'est cette thèse implicite que votre résumé doit faire transparaitre sans jamais l'énoncer explicitement comme une opinion personnelle.
\end{exemplebox}

\begin{exercice}[Entraînement : Tri et Squelette]
Reprenez un court texte narratif de votre manuel (ou un article de presse narratif). Appliquez la méthode :
\begin{enumerate}
\item Lecture globale $\rightarrow$ Thème + Thèse en 1 phrase.
\item Lecture analytique $\rightarrow$ Découpage en 3-4 séquences.
\item Tableau « À garder / À éliminer » (comme ci-dessus).
\item Plan-squelette : 3-4 phrases repères au présent de narration, reliées par connecteurs.
\end{enumerate}
\textit{Correction en classe : comparaison des phrases repères, débat sur ce qui est « essentiel » vs « secondaire ».}
\end{exercice}

\begin{corrige}[Exercice n°1 - Pistes de correction]
La correction portera sur :
\begin{itemize}
\item La justesse du thème/thèse (pas de contresens).
\item La pertinence du découpage (unités de sens cohérentes).
\item La radicalité du tri : tout exemple, citation, détail descriptif doit être dans la colonne « éliminer ».
\item La qualité des phrases repères : autonomie syntaxique, présent de narration, vocabulaire propre (reformulation), connecteurs adaptés.
\item Le respect de la chronologie.
\end{itemize}
\end{corrige}

\coursec{Outils de langue 1 : énoncé constatif et performatif, lexique spécialisé}

Après avoir appris à \textbf{repérer les idées essentielles} du \emph{Vieux nègre et la médaille} (section précédente), nous devons maintenant maîtriser les \textbf{outils linguistiques} qui permettront de les \textbf{reformuler avec exactitude}. Deux notions sont décisives pour le résumé : la distinction entre \cle{énoncé constatif} et \cle{énoncé performatif} (qui change la manière de rapporter un discours) et la gestion du \cle{lexique spécialisé} (qui garantit la précision terminologique).

\courssub{L'énoncé constatif et l'énoncé performatif : dire vs faire}

\textbf{Intuition : le langage comme action.}\par
Dans la vie courante, on croit souvent que parler, c'est seulement \emph{décrire} le monde. Pourtant, certaines paroles \emph{changent} le monde. Quand le commandant dit : « \emph{Je vous décore} », il ne décrit pas une décoration : il \textbf{la réalise} par sa parole même. C'est la découverte majeure de la \textbf{théorie des actes de langage} (J.L. Austin, \emph{Quand dire, c'est faire}, 1962) : \textbf{dire, c'est parfois faire}.

\begin{definition}[Énoncé constatif]
Un \textbf{énoncé constatif} (ou \emph{assertif}) \textbf{describe un état de fait}, une réalité extérieure. On peut le juger \textbf{vrai} ou \textbf{faux} selon qu'il correspond ou non à cette réalité.
\begin{itemize}
    \item Exemple : « Meka reçoit une médaille. » (Vérifiable : vrai si la cérémonie a eu lieu, faux sinon.)
    \item Structure typique : Sujet + verbe d'état ou d'action (indicatif) + complément.
\end{itemize}
\end{definition}

\begin{definition}[Énoncé performatif]
Un \textbf{énoncé performatif} \textbf{accomplit l'action qu'il énonce} du seul fait d'être prononcé dans des conditions appropriées. On ne dit pas qu'il est « vrai » ou « faux », mais qu'il \textbf{réussite} (est \emph{heureux}) ou \textbf{échoue} (est \emph{inheureux}).
\begin{itemize}
    \item Exemple : « Je vous décore au nom de la France. » (Le commandant \emph{fait} la décoration en le disant.)
    \item Verbes performatifs typiques (à la 1\textsuperscript{re} personne, présent, indicatif) : \emph{promettre, jurer, baptiser, nommer, condamner, déclarer ouvert, parier, léguer, décorer}.
\end{itemize}
\end{definition}

\begin{methode}[Distinguier constatif et performatif : le test de la vérité]
Pour classer un énoncé, appliquez ce test en deux temps :
\begin{enumerate}
    \item \textbf{Test de vérité :} Puis-je dire « C'est vrai » ou « C'est faux » ? $\rightarrow$ \textbf{Constatif}.
    \item \textbf{Test de réussite :} Puis-je dire « C'est réussi » ou « C'est raté / nul et non avenu » ? $\rightarrow$ \textbf{Performatif}.
    \item \textbf{Test du \emph{« hereby »} (par la présente) :} Peut-on insérer « par la présente » sans changer le sens ? « Je \textbf{par la présente} vous décore. » $\rightarrow$ \textbf{Performatif explicite}.
\end{enumerate}
\end{methode}

\begin{attention}[Piège : les performatifs déguisés]
Tous les performatifs ne portent pas un verbe performatif explicite.
\begin{itemize}
    \item \textbf{Performatif implicite / indirect :} « La séance est ouverte. » (Le président \emph{ouvre} la séance en le disant.) $\rightarrow$ Équivaut à « J'ouvre la séance. »
    \item \textbf{Constatif en apparence, performatif en force :} « Il fait froid ici. » (Peut être une \emph{demande} de fermer la fenêtre : acte directif).
\end{itemize}
Au résumé, on repère la \textbf{force illocutoire} (l'acte accompli : ordonner, promettre, décorer) pour la restituer.
\end{attention}

\textbf{Application au \emph{Vieux nègre et la médaille} : le discours colonial comme machine performative.}\par
Le roman de Ferdinand Oyono met en scène une administration coloniale dont le pouvoir repose sur la \textbf{parole performative} : nommer des chefs, décorer, condamner, baptiser. Le langage ne \emph{raconte} pas l'ordre colonial, il le \textbf{constitue}.

\begin{popfigure}
\centering
\begin{tikzpicture}[
    cell/.style={rectangle, draw=popdark, thick, minimum height=1.1cm, align=center, text width=7.5cm, inner sep=4pt},
    header/.style={fill=popblue, text=white, font=\bfseries, minimum height=0.9cm},
    headerB/.style={fill=poporange, text=white, font=\bfseries, minimum height=0.9cm},
    crit/.style={fill=popblueL, font=\small\bfseries, minimum height=0.7cm},
    ex/.style={fill=popcream, font=\small, minimum height=0.7cm},
    arrow/.style={->, thick, poporange, shorten >=2pt, shorten <=2pt}
]
% Colonnes
\node[header] (c1) at (0,0) {Énoncé \textbf{Constatif}};
\node[headerB] (c2) at (16,0) {Énoncé \textbf{Performatif}};

% Lignes : Critères
\node[crit, below=0pt of c1.south west, anchor=north west] (cr1) {Critère de reconnaissance};
\node[crit, below=0pt of c2.south west, anchor=north west] (cr2) {Critère de reconnaissance};

\node[ex, below=0pt of cr1.south west, anchor=north west] (ex1) {Se juge \textbf{VRAI} ou \textbf{FAUX}};
\node[ex, below=0pt of cr2.south west, anchor=north west] (ex2) {Se juge \textbf{RÉUSSI} (heureux) ou \textbf{RATÉ} (inheureux)};

\node[ex, below=0pt of ex1.south west, anchor=north west] (ex3) {Verbes : \emph{être, avoir, recevoir, constater, décrire...} (souvent 3\textsuperscript{e} pers.)};
\node[ex, below=0pt of ex2.south west, anchor=north west] (ex4) {Verbes : \emph{promettre, décorer, baptiser, nommer, condamner...} (1\textsuperscript{re} pers. présent + « par la présente »)};

\node[ex, below=0pt of ex3.south west, anchor=north west] (ex5) {Descrit un état \textbf{antérieur} ou \textbf{simultané}};
\node[ex, below=0pt of ex4.south west, anchor=north west] (ex6) {Crée l'état \textbf{par l'énonciation même}};

% Exemples du roman
\node[header, below=0.3cm of ex5.south west, anchor=north west, fill=popgreen, text width=7.5cm] (exR1) {Exemples dans \emph{Le Vieux nègre et la médaille}};
\node[headerB, below=0.3cm of ex6.south west, anchor=north west, fill=poppurple, text width=7.5cm] (exR2) {Exemples dans \emph{Le Vieux nègre et la médaille}};

\node[ex, below=0pt of exR1.south west, anchor=north west] (e1) {« Meka \textbf{a reçu} la médaille. » (Récit du fait accompli.)};
\node[ex, below=0pt of exR2.south west, anchor=north west] (e2) {« Je vous \textbf{décore} au nom de la France. » (Le commandant \textbf{accomplit} l'acte.)};

\node[ex, below=0pt of e1.south west, anchor=north west] (e3) {« Le Père Vandermayer \textbf{a baptisé} Meka. » (Constat a posteriori.)};
\node[ex, below=0pt of e2.south west, anchor=north west] (e4) {« Je vous \textbf{baptise} au nom du Père... » (Le prêtre \textbf{fait} le baptême.)};

\node[ex, below=0pt of e3.south west, anchor=north west] (e5) {« Le chef \textbf{a été nommé} par l'administrateur. » (Description d'une nomination passée.)};
\node[ex, below=0pt of e4.south west, anchor=north west] (e6) {« Je vous \textbf{nomme} chef de canton. » (L'administrateur \textbf{créé} le chef.)};

\node[ex, below=0pt of e5.south west, anchor=north west] (e7) {« Meka \textbf{est content} de sa médaille. » (Description d'un état psychologique.)};
\node[ex, below=0pt of e6.south west, anchor=north west] (e8) {« Je vous \textbf{félicite}, Meka. » (Acte de politesse / reconnaissance officielle.)};

% Flèches conditions de réussite
\node[below=0.3cm of e8.south west, anchor=north west, text width=15.5cm, font=\small\itshape, text=popdark] (cond) {\textbf{Conditions de réussite (félicité) d'un performatif (Austin) :} 1. Procédure conventionnelle existante (code colonial, rituel catholique). 2. Circonstances / personnes appropriées (commandant légitime, Meka présent). 3. Exécution correcte et complète. 4. Intentions sincères (le commandant veut-il vraiment honorer Meka ? $\rightarrow$ Ironie dramatique : le performatif est \textbf{détourné}, \textbf{vide} de sens véritable).};
\end{tikzpicture}
\legende{Tableau comparatif : constatif vs performatif dans \emph{Le Vieux nègre et la médaille}. Le pouvoir colonial s'exerce par la performativité (nommer, décorer), que le résumé doit restituer sans la réduire à une simple description.}
\end{popfigure}

\begin{exemplebox}[Performatif colonial et ironie]
\textbf{Énoncé :} « Je vous décerne la médaille de la Reconnaissance française. » (Commandant, p.~72 env.)
\begin{itemize}
    \item \textbf{Analyse :} Verbe \emph{décerner} (1\textsuperscript{re} pers. présent) + autorité investie + cérémonie officielle $\rightarrow$ \textbf{Performatif explicite réussi formellement}.
    \item \textbf{Portée romanesque :} Oyono montre le \textbf{décalage} entre la force performative (l'acte administratif est valide) et la \textbf{valeur symbolique} (la médaille est « en toc », Meka est manipulé). Le performatif devient \textbf{instrument de domination symbolique}.
\end{itemize}
\end{exemplebox}

\textbf{Enjeu pour le résumé : reformuler le performatif en constatif sans perte de force.}\par
Le résumé est un discours \textbf{rapporté} (souvent au passé, 3\textsuperscript{e} personne). Il transforme \textbf{systématiquement} les performatifs en constatifs.
\begin{center}
\begin{tabularx}{\linewidth}{|>{\raggedright\arraybackslash}X|>{\raggedright\arraybackslash}X|>{\raggedright\arraybackslash}X|}
\hline
\rowcolor{popblueL} \textbf{Discours original (Performatif)} & \textbf{Transformation abusive (Perte de sens)} & \textbf{Reformulation correcte (Constatif précis)} \\
\hline
« Je vous \textbf{décore}. » & « Il a donné une médaille. » (Trop vague, perd l'officialité) & « Le commandant \textbf{a décoré} Meka. » (Verbe d'action accomplie, conserve l'acte institutionnel) \\
\hline
« Je vous \textbf{nomme} chef. » & « Il est devenu chef. » (Perte de l'agent et de l'acte juridique) & « L'administrateur \textbf{l'a nommé} chef. » (Agent + verbe performatif passé) \\
\hline
« Je vous \textbf{baptise}... » & « Il est devenu chrétien. » (Perte de l'acte sacramentel) & « Le Père Vandermayer \textbf{l'a baptisé}. » \\
\hline
\end{tabularx}
\end{center}

\begin{aretenir}[Règle d'or pour le résumé]
\textbf{Conservez le verbe d'action performative au passé composé (ou imparfait) à la 3\textsuperscript{e} personne.}
Ne le remplacez \textbf{jamais} par un verbe vague (\emph{donner, faire, rendre}) ni par une périphrase qui gomme l'institution (administration, église, justice). Le résumé doit dire \textbf{QUI a FAIT QUOI} par sa parole.
\end{aretenir}

\begin{exoresolu}[Repérage et reformulation]
\textbf{Énoncé :} Extrait du discours du Commandant : « \emph{Au nom du Président de la République française, je vous fais chevalier de l'Ordre du Mérite. Recevez cette médaille comme gage de notre reconnaissance.} »
\begin{enumerate}
    \item \textbf{Identifiez les performatifs.}
    \item \textbf{Reformulez en phrase de résumé (constative).}
\end{enumerate}
\tcblower
\textbf{Solution :}
\begin{enumerate}
    \item \textbf{« Je vous fais chevalier... »} (Verbe \emph{faire} + titre honorifique = performatif de nomination/décoration). \textbf{« Recevez cette médaille... »} (Impératif performatif : acte de remise officielle).
    \item \textbf{Résumé :} « Le commandant, au nom du Président de la République, \textbf{a fait chevalier} Meka de l'Ordre du Mérite et \textbf{lui a remis} la médaille en gage de reconnaissance. »
    \begin{itemize}
        \item \emph{A fait chevalier} : conserve la force institutionnelle (création du titre).
        \item \emph{Lui a remis} : rend l'acte de « Recevez » (transfert de possession officiel).
    \end{itemize}
\end{enumerate}
\end{exoresolu}

\courssub{Le lexique spécialisé : précision terminologique et fidélité au texte}

\textbf{Observation : le texte comme lieu de savoirs techniques.}\par
\emph{Le Vieux nègre et la médaille} n'est pas seulement un roman psychologique ; il est traversé par des \textbf{langues de spécialité} (ou \emph{lectes}) : l'\textbf{administration coloniale}, le \textbf{droit coutumier / chefferie}, la \textbf{religion catholique}, l'\textbf{agriculture cacaoyère}. Chaque domaine possède son \textbf{lexique spécialisé} : des mots dont le sens est \textbf{figé, précis, non interchangeable} par un synonyme courant sans perte d'information.

\begin{definition}[Lexique spécialisé (termes techniques)]
Ensemble des \textbf{unités lexicales} (mots simples ou composés) propres à un \textbf{domaine d'activité professionnel, scientifique, technique, administratif ou religieux}.
\begin{itemize}
    \item \textbf{Caractères :} Monosémie (un sens précis dans le domaine), absence de connotation vague, nécessité pour l'expert / l'initié.
    \item \textbf{Exemples :} \emph{Arrêté, décret, canton, chef de poste, administrateur, médaille du Mérite, chevalier, ordonnance} (Admin.) ; \emph{Baptême, catéchumène, parrain, marraine, hostie, bénédicité, vicaire} (Religion) ; \emph{Plantation, cabosse, écabossage, fermentation, séchage, courtier, tonne} (Agriculture).
\end{itemize}
\end{definition}

\begin{popfigure}
\centering
\begin{tikzpicture}[
    scale=0.85,
    every node/.style={transform shape},
    cloud/.style={draw=popdark, thick, rounded corners=15pt, minimum height=2.5cm, minimum width=4.5cm, align=center, inner sep=6pt, font=\small},
    term/.style={font=\footnotesize, text=popdark},
    title/.style={font=\bfseries\small, text=white},
    admin/.style={cloud, fill=popblueL, label={[title, anchor=south, yshift=-2pt]above:Administration Coloniale}},
    relig/.style={cloud, fill=poppurpleL, label={[title, anchor=south, yshift=-2pt]above:Religion Catholique}},
    agri/.style={cloud, fill=popgreenL, label={[title, anchor=south, yshift=-2pt]above:Agriculture / Cacao}},
    custom/.style={cloud, fill=poporangeL, label={[title, anchor=south, yshift=-2pt]above:Coutume / Chefferie}},
    link/.style={-Stealth, thick, popmuted, shorten >=3pt, shorten <=3pt}
]
% Nuages principaux
\node[admin, align=center] (A) at (0,0){
    \begin{tabular}{c}
    \textbf{Arrêté} \quad \textbf{Décret} \quad \textbf{Ordonnance} \\
    \textbf{Chef de poste} \quad \textbf{Administrateur} \quad \textbf{Cercle} \\
    \textbf{Canton} \quad \textbf{Chef de canton} \quad \textbf{Indigène} \\
    \textbf{Médaille du Mérite} \quad \textbf{Chevalier} \quad \textbf{Reconnaissance} \\
    \textbf{Prestation de serment} \quad \textbf{Registre} \quad \textbf{Archives}
    \end{tabular}
};

\node[relig, align=center] (R) at (9,0){
    \begin{tabular}{c}
    \textbf{Baptême} \quad \textbf{Catéchumène} \quad \textbf{Parrain/Marraine} \\
    \textbf{Père Vandermayer} \quad \textbf{Vicaire} \quad \textbf{Mission} \\
    \textbf{Hostie} \quad \textbf{Bénédicité} \quad \textbf{Confession} \\
    \textbf{Communion} \quad \textbf{Extrême-onction} \quad \textbf{Paroisse}
    \end{tabular}
};

\node[agri, align=center] (G) at (4.5,-5){
    \begin{tabular}{c}
    \textbf{Plantation} \quad \textbf{Cacaoyer} \quad \textbf{Cabosse} \\
    \textbf{Écabossage} \quad \textbf{Fermentation} \quad \textbf{Séchage} \\
    \textbf{Courtier} \quad \textbf{Acheteur} \quad \textbf{Tonne} \\
    \textbf{Rendement} \quad \textbf{Récolte} \quad \textbf{Main-d'œuvre}
    \end{tabular}
};

\node[custom, align=center] (C) at (-4.5,-5){
    \begin{tabular}{c}
    \textbf{Médecine traditionnelle} \quad \textbf{Fétiche} \quad \textbf{Initiation} \\
    \textbf{Ancêtres} \quad \textbf{Esprits} \quad \textbf{Tabou} \\
    \textbf{Conseil des notables} \quad \textbf{Palabre} \quad \textbf{Terre} \\
    \textbf{Héritage} \quad \textbf{Polygamie} \quad \textbf{Dot}
    \end{tabular}
};

% Flèches d'interaction (Meka au centre)
\node[circle, draw=popink, thick, minimum size=2.2cm, fill=poppinkL, font=\bfseries\small, text=popdark, align=center] (M) at (4.5,-2.5){MEKA\\ \emph{(Carrefour des lexiques)}};

\draw[link] (A.south) -- (M.north);
\draw[link] (R.south) -- (M.north);
\draw[link] (G.north) -- (M.south);
\draw[link] (C.north) -- (M.south);

% Annotations
\node[below=0.2cm of G.south, font=\itshape\footnotesize, text=popgreen, text width=4cm, align=center] {Lexique de la \textbf{production} et de l'\textbf{économie}};
\node[below=0.2cm of C.south, font=\itshape\footnotesize, text=poporange, text width=4cm, align=center] {Lexique de la \textbf{tradition} et du \textbf{social}};
\node[above=0.2cm of A.north, font=\itshape\footnotesize, text=popblue, text width=4cm, align=center] {Lexique du \textbf{pouvoir} et de la \textbf{légalité}};
\node[above=0.2cm of R.north, font=\itshape\footnotesize, text=poppurple, text width=4cm, align=center] {Lexique du \textbf{sacré} et de l'\textbf{identité}};
\end{tikzpicture}
\legende{Nuage des champs lexicaux spécialisés dans \emph{Le Vieux nègre et la médaille}. Meka se situe à l'intersection de ces quatre langues techniques. Le résumé doit conserver les termes clés (en gras) qui structurent l'intrigue.}
\end{popfigure}

\begin{methode}[Repérer et traiter le lexique spécialisé dans un texte à résumer]
\begin{enumerate}
    \item \textbf{Identifiez le(s) domaine(s) :} Administration ? Droit ? Médecine ? Technique ? Religion ? Économie ?
    \item \textbf{Surlignez les termes techniques :} Mots monosémiques, souvent des noms propres de fonction (\emph{administrateur}), d'acte (\emph{arrêté}), d'objet (\emph{cabosse}), de procédure (\emph{fermentation}).
    \item \textbf{Vérifiez l'indispensabilité :} Le terme est-il \textbf{irremplaçable} pour la compréhension exacte de l'intrigue ou de l'argumentation ? (Ex. : « \emph{écabossage} » ne se résume pas par « casser les cabosses » sans perte technique).
    \item \textbf{Conservez le terme exact dans le résumé.} N'utilisez \textbf{jamais} de périphrase vague (« le papier officiel » pour « \emph{arrêté} », « le prêtre » pour « \emph{vicaire} » si le grade compte).
    \item \textbf{Contexte :} Si le terme est très pointu et rare, assurez-vous qu'il est compréhensible par le contexte immédiat du résumé (ex. : « la \emph{fermentation} des fèves »).
\end{enumerate}
\end{methode}

\begin{attention}[Erreur fatale au résumé : la « traduction » appauvrissante]
\begin{itemize}
    \item \textbf{Texte :} « Le \textbf{chef de poste} a signé l'\textbf{arrêté} de nomination. »
    \item \textbf{Mauvais résumé :} « Le \textbf{chef} a signé le \textbf{papier} de nomination. » $\rightarrow$ Perte de la \textbf{valeur juridique} (\emph{arrêté} = acte administratif réglementaire ; \emph{chef de poste} = grade précis dans la hiérarchie coloniale).
    \item \textbf{Bon résumé :} « Le \textbf{chef de poste} a signé l'\textbf{arrêté} de nomination. »
\end{itemize}
\textbf{Règle :} \textbf{Termes spécialisés = Invariants du résumé.} On ne les change pas, on ne les explique pas (sauf si le sujet l'exige explicitement), on les \textbf{transpose} tels quels.
\end{attention}

\begin{exemplebox}[Lexique spécialisé et performatif : couplage fréquent]
Dans le roman, les performatifs \textbf{mobilisent} le lexique spécialisé.
\begin{quote}
« \emph{Au vu de l'\textbf{arrêté} n° 45/23, je vous \textbf{nomme} \textbf{chef de canton} de Doum. } »
\end{quote}
\begin{itemize}
    \item \textbf{Performatifs :} \emph{nomme} (création de la fonction).
    \item \textbf{Lexique spécialisé :} \emph{arrêté} (source juridique), \emph{chef de canton} (titre administratif précis), \emph{Doum} (toponyme administratif).
    \item \textbf{Résumé correct :} « Sur la base de l'\textbf{arrêté} n° 45/23, l'administrateur \textbf{a nommé} Meka \textbf{chef de canton} de Doum. »
\end{itemize}
\end{exemplebox}

\begin{exoresolu}[Traitement conjoint : performatif + lexique spécialisé]
\textbf{Texte source (inventé, style Oyono) :}
« \emph{Le Père Vandermayer, après la \textbf{messe} de onze heures, a convoqué Meka. "Mon fils, tu as bien appris ton \textbf{catéchisme}. Je te \textbf{baptise} aujourd'hui même. Tu auras pour \textbf{parrain} le \textbf{chef de poste}." Meka a reçu le \textbf{nom de baptême} de Michel. } »

\textbf{Consigne :} Résumer en 2 phrases max en respectant performatifs et lexique spécialisé.

\tcblower
\textbf{Solution commentée :}
\begin{enumerate}
    \item \textbf{Analyse :}
        \begin{itemize}
            \item Performatifs : « Je te baptise » (Père), « Tu auras pour parrain... » (désignation performative).
            \item Lexique spécialisé (Religion/Admin) : \emph{messe, catéchisme, baptise, parrain, chef de poste, nom de baptême}.
        \end{itemize}
    \item \textbf{Rédaction :}
    \begin{quote}
    « Après la \textbf{messe}, le Père Vandermayer \textbf{a baptisé} Meka, qui \textbf{a reçu} le \textbf{nom de baptême} de Michel ; le \textbf{chef de poste} a été désigné comme \textbf{parrain}. »
    \end{quote}
    \item \textbf{Justification :}
        \begin{itemize}
            \item \emph{A baptisé} (constatif passé du performatif \emph{baptise}) : conserve l'acte sacramentel.
            \item \emph{A reçu le nom de baptême} : rend le performatif implicite de nomination chrétienne.
            \item \emph{Chef de poste, parrain, messe, nom de baptême} : termes spécialisés conservés intacts.
        \end{itemize}
\end{enumerate}
\end{exoresolu}

\courssub{Synthèse et articulation avec la rédaction du résumé}

\begin{aretenir}[Bilan des outils de langue pour le résumé (Sections 2 + 3)]
\begin{tabularx}{\linewidth}{|l|X|X|}
\hline
\rowcolor{popblueL} \textbf{Opération de résumé} & \textbf{Apport de la Section 2 (Idées essentielles)} & \textbf{Apport de la Section 3 (Outils de langue)} \\
\hline
\textbf{Sélection} & Hiérarchiser : Idées forces > Détails & Repérer les \textbf{performatifs} (actes majeurs) et le \textbf{lexique spécialisé} (termes invariants) \\
\hline
\textbf{Reformulation} & Généraliser, synthétiser & \textbf{Convertir performatif $\rightarrow$ constatif} (verbe d'action au passé) ; \textbf{Conserver termes spécialisés} \\
\hline
\textbf{Structuration} & Ordonner logiquement (chronologique / thématique) & Utiliser la \textbf{valeur illocutoire} pour ordonner les actes (nomination $\rightarrow$ décoration $\rightarrow$ réaction) \\
\hline
\end{tabularx}
\end{aretenir}

\begin{methode}[Check-list avant d'écrire le résumé (Intégration S2+S3)]
Avant de passer à la rédaction (Section 5), vérifiez votre brouillon :
\begin{enumerate}
    \item [$\square$] Ai-je listé les \textbf{actes performatifs majeurs} du texte (décorations, nominations, baptêmes, promesses, menaces) ?
    \item [$\square$] Les ai-je transformés en \textbf{constatifs précis} (Sujet institutionnel + verbe d'action passé + Complément) ?
    \item [$\square$] Ai-je identifié \textbf{tous les termes spécialisés} indispensables (titres, actes juridiques, termes techniques, grades) ?
    \item [$\square$] Me suis-je interdit toute \textbf{périphrase vague} à la place d'un terme technique ?
    \item [$\square$] La \textbf{tonalité} (ironie, satire — voir Section 4) est-elle perceptible dans le choix des verbes de rapport (\emph{prétendre, faire semblant de, décerner ironiquement}) ?
\end{enumerate}
\end{methode}

\begin{savaistu}[Austin, Searle et la pragmatique au Cameroun]
La distinction constatif/performatif, fondatrice de la \textbf{pragmatique} (étude du langage en situation), éclaire les \textbf{litiges coutumiers} et \textbf{administratifs} actuels au Cameroun.
\begin{itemize}
    \item \textbf{Investiture d'un chef :} La phrase « Je vous installe chef » prononcée par l'autorité administrative (SDO/DO) ou le roi des rois est un \textbf{performatif institutionnel}. Sa validité dépend des \textbf{conditions de félicité} (légitimité coutumière + reconnaissance administrative).
    \item \textbf{Actes de justice :} « Je vous condamne à... » (Juge) / « Je vous pardonne » (Chef traditionnel). La force obligatoire ne vient pas de la vérité du propos mais de l'\textbf{autorité} du locuteur.
    \item \textbf{Contrats oraux :} « Je te donne ma parole » (Promesse performative) engage l'honneur dans les sociétés à tradition orale, souvent plus fort qu'un écrit.
\end{itemize}
Maîtriser ces notions, c'est comprendre comment \textbf{la parole structure le pouvoir et le lien social}, hier comme aujourd'hui.
\end{savaistu}

\textbf{Transition vers la Section 4 :} Nous savons maintenant \emph{comment} reformuler les actes de langage (constatif/performatif) et \emph{quels mots} garder (lexique spécialisé). Mais le texte d'Oyono n'est pas neutre : il est traversé par l'\textbf{ironie}, la \textbf{satire} et la \textbf{polémique}. La Section 4 (« Outils de langue 2 : tonalités satirique et polémique ») nous apprendra à \textbf{restituer le ton} sans le trahir, car un résumé qui perd l'ironie de l'auteur perd la moitié du sens.

\coursec{Outils de langue 2 : tonalités satirique et polémique}

Dans la section précédente, nous avons appris à distinguer l'\cle{énoncé constatif} (qui décrit le monde) de l'\cle{énoncé performatif} (qui agit sur le monde), et à manier le \cle{lexique spécialisé}. Mais un texte ne dit pas seulement \emph{ce qu'il dit} : il le dit \emph{d'une certaine manière}. Or, dans \emph{Le Vieux nègre et la médaille}, Ferdinand Oyono ne se contente pas de raconter les mésaventures de Meka : il \emph{se moque}, il \emph{dénonce}, il \emph{attaque}. Cette manière de dire, cette attitude de l'énonciateur face à son sujet, c'est ce qu'on appelle la \cle{tonalité} (ou \cle{registre}). La comprendre est indispensable pour rédiger un résumé fidèle : un résumé qui « aplatit » le ton trahit le texte.

\courssub{La tonalité : définition et repérage}

\begin{definition}[La tonalité (ou registre)]
La \cle{tonalité} (ou \cle{registre}) est l'attitude que l'énonciateur adopte face à son sujet et face à son destinataire. Elle est perceptible à travers le \textbf{choix des mots} (lexique mélioratif ou péjoratif), les \textbf{procédés d'écriture} (ironie, exagération, interrogation) et le \textbf{ton général} (rire, indignation, émotion, admiration).
\end{definition}

\textbf{Intuition.} Imaginez deux élèves qui commentent la même note catastrophique à un devoir. Le premier dit : « Bravo, vraiment, quel chef-d'œuvre ! » en souriant. Le second dit : « Ce résultat est inacceptable, il faut réagir ! » en frappant du poing sur la table. Les deux critiquent la même chose, mais pas de la même manière : le premier \emph{raille} (tonalité satirique), le second \emph{attaque} (tonalité polémique). Le sens littéral des mots ne suffit donc pas : il faut écouter la \emph{voix} du texte.

\begin{methode}[Identifier la tonalité d'un passage]
\begin{enumerate}
\item \textbf{Relever le lexique} : les mots sont-ils \cle{péjoratifs} (ils dévalorisent : « dérision », « mascarade », « prétendue récompense ») ou \cle{mélioratifs} (ils valorisent) ? Y a-t-il un écart suspect entre des mots élogieux et une réalité désastreuse ?
\item \textbf{Repérer les figures et procédés} : \cle{ironie} et \cle{antiphrase} (dire le contraire de ce qu'on pense), \cle{hyperbole} (exagération), \cle{caricature} (déformation grotesque d'un personnage), \cle{interrogation rhétorique} (question qui n'attend pas de réponse).
\item \textbf{Déterminer le ton} : l'énonciateur rit-il (satire), s'indigne-t-il (polémique), s'émeut-il (lyrique) ?
\item \textbf{Conclure par une formule précise} : « l'auteur dénonce \emph{avec ironie}… », « le journaliste attaque \emph{violemment}… ».
\end{enumerate}
\end{methode}

\courssub{La tonalité satirique : critiquer par le rire}

\begin{definition}[La satire]
La \cle{satire} est un discours qui critique \textbf{par le rire et l'ironie} les défauts, les vices et les ridicules d'une société ou d'individus. Elle ne se contente pas de faire rire : elle veut \textbf{corriger} en dénonçant. On dit alors que l'auteur « raille » ou « fustige » les travers de son époque.
\end{definition}

\textbf{Les procédés satiriques principaux.}

\begin{itemize}
\item L'\cle{ironie} : procédé général qui consiste à laisser entendre le contraire de ce qu'on affirme, souvent avec un sourire.
\item L'\cle{antiphrase} : figure d'ironie qui consiste à \textbf{dire exactement le contraire} de ce qu'on pense (« Quelle belle récompense ! » pour désigner une dérision).
\item La \cle{caricature} : portrait déformé et grotesque qui grossit les défauts d'un personnage (le commandant vaniteux, le prêtre complaisant).
\item L'\cle{hyperbole} (ou exagération) : amplification outrancière qui rend le ridicule visible (« une médaille qui valait tous les trésors du monde »).
\end{itemize}

\begin{savaistu}[L'antiphrase, arme favorite d'Oyono]
L'\cle{antiphrase} consiste à dire le contraire de ce qu'on pense. Quand Ferdinand Oyono écrit, dans \emph{Le Vieux nègre et la médaille} (1956), que la médaille offerte à Meka est une grande « récompense », il dénonce en réalité une \textbf{dérision} : le vieil homme a perdu ses terres (données à la mission catholique) et ses deux fils (morts à la guerre pour la France), et on le « remercie » avec un simple bout de métal. L'éloge apparent est donc une accusation déguisée. Oyono, écrivain camerounais devenu plus tard diplomate et homme d'État, est avec Mongo Beti l'un des grands satiristes de la colonisation en Afrique centrale.
\end{savaistu}

\textbf{Observation guidée dans le roman.} Menons l'investigation comme au laboratoire : observons les faits du récit, puis interprétons.

\begin{experience}[Repérage de la satire dans \emph{Le Vieux nègre et la médaille}]
\textbf{Matériel} : le roman, un tableau de relevés à deux colonnes (« ce qui est dit » / « ce qui est pensé »).

\textbf{Protocole} :
\begin{enumerate}
\item Relever les « dons » faits par Meka aux Blancs : ses terres à la mission, ses deux fils à la guerre.
\item Relever ce que Meka reçoit en échange : une médaille, remise lors de la fête du 14 juillet.
\item Observer le comportement des personnages coloniaux : le commandant, le père Vandermayer, le gosier (interprète).
\item Confronter les deux colonnes et formuler l'écart.
\end{enumerate}

\textbf{Observations} : le sacrifice de Meka est immense (terres, fils, sueur pour les chantiers), la contrepartie est dérisoire (une médaille sans valeur). Le discours officiel, lui, parle de « reconnaissance » et d'« amitié franco-africaine ».

\textbf{Interprétation} : l'écart entre le discours élogieux et la réalité misérable est le moteur de l'\cle{ironie} d'Oyono. En prenant au sérieux le vocabulaire officiel, le romancier en révèle l'hypocrisie : c'est toute l'administration coloniale et le clergé qui sont tournés en dérision.
\end{experience}

\courssub{La tonalité polémique : attaquer pour réfuter}

\begin{definition}[Le registre polémique]
Le discours \cle{polémique} est un discours de \textbf{controverse} qui réfute et attaque de manière vive une thèse, une opinion, une institution ou un adversaire. Contrairement à la satire qui rit, la polémique \textbf{s'indigne et accuse} ouvertement. Ses procédés typiques sont l'\cle{accusation directe}, l'\cle{interrogation rhétorique} et le \cle{lexique péjoratif}.
\end{definition}

\textbf{Les procédés polémiques principaux.}

\begin{itemize}
\item L'\cle{accusation directe} : on nomme le coupable et on le désigne (« l'administration a spolié ce vieillard »).
\item L'\cle{interrogation rhétorique} : question dont la réponse est évidente, qui accule l'adversaire (« Est-ce là le prix d'une vie de sacrifices ? »).
\item Le \cle{lexique péjoratif} : champs lexicaux de la violence, de l'injustice, du mépris (« spoliation », « mascarade », « hypocrisie »).
\item La \cle{réfutation} : on cite la thèse adverse pour la démolir (« On prétend que la colonisation a civilisé ; or… »).
\end{itemize}

\textbf{Satire ou polémique ?} Les deux tonalités dénoncent, mais différemment : la satire utilise le \emph{rire} comme arme (elle fait sourire le lecteur qui devient complice), la polémique utilise l'\emph{indignation} (elle mobilise le lecteur contre un adversaire désigné). Un même texte peut passer de l'une à l'autre : chez Oyono, l'ironie satirique domine, mais certains passages (l'arrestation de Meka, la nuit sous la pluie) basculent vers une dénonciation plus directe et pathétique.

\begin{attention}[Le piège majeur : prendre l'ironie au premier degré]
L'erreur la plus grave, dans un résumé, est de \textbf{prendre l'ironie au premier degré} et d'attribuer à l'auteur le contraire de sa pensée. Écrire « Oyono montre que la médaille récompense dignement les services rendus par Meka à la France » est une \textbf{contre-vérité} : Oyono pense exactement l'inverse. Pour éviter ce piège : (1) vérifiez toujours la cohérence globale du texte (un éloge peut-il être sincère quand tout le reste est désastreux ?) ; (2) dans le résumé, signalez la tonalité avec un verbe d'intention : « Oyono \emph{dénonce avec ironie}… », « l'auteur \emph{tourne en dérision}… ».
\end{attention}

\courssub{Tableau de synthèse : procédés et exemples du roman}

\begin{center}
\begin{tabularx}{\linewidth}{|l|X|X|}
\hline
\rowcolor{popblueL} \textbf{Procédé} & \textbf{Définition} & \textbf{Exemple dans \emph{Le Vieux nègre et la médaille}} \\
\hline
Antiphrase (ironie) & Dire le contraire de ce qu'on pense & La médaille présentée comme une grande « récompense » alors qu'elle est dérisoire face aux terres et aux fils perdus. \\
\hline
Caricature & Portrait grotesque qui grossit les défauts & Le commandant et les notables coloniaux, vaniteux et ridicules dans leur mise en scène de la fête du 14 juillet. \\
\hline
Hyperbole & Exagération qui rend le ridicule visible & L'importance démesurée accordée par tous à un simple bout de métal, objet d'envies et de convoitises dans le village. \\
\hline
Accusation directe (polémique) & Désigner nommément le responsable & La dénonciation de l'administration coloniale et du clergé, qui profitent de la naïveté de Meka. \\
\hline
Interrogation rhétorique & Question à réponse évidente qui accule l'adversaire & « Est-ce là le prix de tant de sacrifices ? » (reformulation du questionnement qui traverse le roman). \\
\hline
Lexique péjoratif & Mots qui dévalorisent la cible & Lexique de la dérision et de la mascarade appliqué à la cérémonie de remise de la médaille. \\
\hline
\end{tabularx}
\end{center}

\courssub{Une échelle des tonalités}

Les tonalités ne sont pas des cases fermées : elles forment un \textbf{continuum}, de l'expression des sentiments (lyrique) à l'attaque frontale (polémique), en passant par la critique par le rire (satirique). Le schéma suivant positionne les passages étudiés.

\begin{popfigure}
\begin{center}
\begin{tikzpicture}[scale=1.0]
% Axe principal
\draw[-{Stealth[length=3mm]},very thick,popline] (0,0) -- (12.5,0);
% Graduations
\foreach \x/\c/\titre/\desc in {
  1.5/popgreen/{Lyrique}/{exprime les sentiments},
  6.25/poporange/{Satirique}/{critique par le rire},
  11.0/popink/{Polémique}/{attaque et réfute}} {
  \fill[\c] (\x,0) circle (0.12);
  \node[above,font=\bfseries,text=\c] at (\x,0.15) {\titre};
  \node[below,font=\footnotesize,text=popmuted,align=center] at (\x,-0.15) {\desc};
}
% Positionnement des extraits
\node[above=0.75cm,align=center,font=\footnotesize,text=popgreen] at (1.5,0.2) {éloge du village\\ de Doum (nostalgie)};
\draw[->,popgreen] (1.5,0.35) -- (1.5,0.62);
\node[above=0.75cm,align=center,font=\footnotesize,text=poporange] at (6.25,0.2) {la « récompense »\\ de la médaille (ironie)};
\draw[->,poporange] (6.25,0.35) -- (6.25,0.62);
\node[above=0.75cm,align=center,font=\footnotesize,text=popink] at (11.0,0.2) {arrestation de Meka :\\ dénonciation de l'injustice};
\draw[->,popink] (11.0,0.35) -- (11.0,0.62);
% Flèche d'intensité
\node[font=\footnotesize\itshape,text=popmuted] at (6.25,-0.95) {intensité croissante de l'attaque};
\end{tikzpicture}
\end{center}
\legende{Échelle des tonalités : du lyrique (expression des sentiments) au polémique (attaque frontale), en passant par le satirique (critique par le rire). Les passages d'\emph{Le Vieux nègre et la médaille} se répartissent sur cette échelle ; l'ironie satirique y domine.}
\end{popfigure}

\courssub{Enjeu pour le résumé : signaler la tonalité sans la reproduire}

\begin{aretenir}[La tonalité dans le résumé]
Un résumé fidèle ne doit pas \textbf{aplatir} la tonalité du texte, mais il ne doit pas non plus la \textbf{reproduire} longuement (pas de citations ironiques, pas de caricatures reprises). La bonne pratique consiste à la \textbf{signaler brièvement} par un verbe ou un adverbe d'intention :
\begin{itemize}
\item « Oyono \emph{dénonce avec ironie} l'hypocrisie coloniale… »
\item « L'auteur \emph{tourne en dérision} la prétendue reconnaissance de l'administration… »
\item « Le journaliste \emph{s'insurge contre} cette décision… »
\end{itemize}
Ainsi, le lecteur du résumé sait \emph{comment} le texte dit ce qu'il dit, en une seule formule.
\end{aretenir}

\begin{exoresolu}[Identifier la tonalité d'un éditorial]
\textbf{Énoncé.} Lisez cet extrait (inventé) d'éditorial de presse camerounaise : « Quelle merveilleuse décision ! Couper l'eau potable de tout un quartier pendant trois semaines, voilà assurément une belle manière de "servir" les populations. Faut-il rappeler à nos édiles que l'eau est un droit et non une faveur ? Les habitants d'Obili méritent mieux que ce mépris organisé. »

1. Relevez deux procédés et nommez-les. 2. Identifiez la ou les tonalités dominantes. 3. Résumez cet extrait en une phrase qui signale la tonalité.

\tcblower
\textbf{Solution détaillée.}

1. \textbf{Procédés} : « Quelle merveilleuse décision ! » est une \cle{antiphrase} (ironie) : le journaliste pense le contraire de « merveilleuse ». « Faut-il rappeler à nos édiles que… ? » est une \cle{interrogation rhétorique} : la réponse est évidemment « oui ». On relève aussi un \cle{lexique péjoratif} : « mépris organisé ».

2. \textbf{Tonalité} : l'extrait commence sur le mode \cle{satirique} (ironie, antiphrase, guillemets moqueurs autour de « servir ») puis bascule vers le \cle{polémique} (interrogation rhétorique, accusation directe : « ce mépris organisé »). La tonalité dominante est polémique, teintée d'ironie.

3. \textbf{Résumé en une phrase} : « Le journaliste dénonce avec ironie, puis s'insurge contre la coupure prolongée d'eau potable à Obili, qu'il présente comme une violation d'un droit fondamental et une marque de mépris envers les habitants. » On remarque que la tonalité est signalée (« dénonce avec ironie », « s'insurge ») sans être reproduite.
\end{exoresolu}

\begin{exercice}[Entraînement : satirique ou polémique ?]
Pour chacun des extraits suivants (inspirés de la presse et de la littérature), identifiez la tonalité dominante et justifiez par deux procédés relevés.

1. « Bravo à nos dirigeants : grâce à eux, les nids-de-poule de nos routes sont devenus de véritables piscines municipales ! »

2. « Il est inadmissible que les coupables de ce détournement restent impunis. Qui protège ces prédateurs des deniers publics ? »

3. « Le commandant, droit comme une statue dans son uniforme trop neuf, attendait que les indigènes admirent enfin sa grandeur. »

4. Rédigez le résumé de l'extrait 1 en une phrase qui signale la tonalité.
\end{exercice}

\begin{corrige}[Exercice : satirique ou polémique ?]
1. \textbf{Tonalité satirique} : antiphrase (« Bravo »), hyperbole ironique (« piscines municipales ») qui ridiculisent l'inaction des dirigeants.

2. \textbf{Tonalité polémique} : accusation directe (« détournement », « prédateurs »), interrogation rhétorique (« Qui protège… ? »), lexique péjoratif.

3. \textbf{Tonalité satirique} : caricature (le commandant « droit comme une statue », uniforme « trop neuf »), ironie sur sa vanité (« admirent enfin sa grandeur ») ; on reconnaît le regard d'Oyono sur les notables coloniaux.

4. « L'auteur de ce billet raille avec ironie l'inaction des dirigeants face à la dégradation des routes, qu'il compare moqueusement à des piscines. »
\end{corrige}

\begin{attention}[Dernières précautions avant la rédaction]
\begin{itemize}
\item Ne confondez pas \cle{satire} et \cle{humour} gratuit : la satire a toujours une \textbf{visée critique} (corriger, dénoncer).
\item Dans un résumé, n'écrivez jamais « l'auteur dit que… » pour une phrase ironique : écrivez « l'auteur \emph{fait mine de} dire que… » ou « l'auteur \emph{ironise sur}… ».
\item Une caricature légendée (dessin de presse) se lit comme un texte : repérez ce que le dessin \emph{déforme} (caricature) et ce que la légende \emph{affirme} (souvent une antiphrase).
\end{itemize}
\end{attention}

Nous disposons désormais de tous les outils de langue nécessaires : énoncés constatifs et performatifs, lexique spécialisé, tonalités satirique et polémique. La prochaine section mettra ces outils au service de la \textbf{rédaction effective du résumé} : reformuler les idées essentielles sans trahir ni le sens ni le ton du texte.

\coursec{Rédaction du résumé : reformulation et production}

Après avoir analysé les \emph{tonalités satirique et polémique} qui colorent \emph{Le Vieux nègre et la médaille}, nous disposons désormais des clés de lecture pour comprendre comment Ferdinand Oyono construit son propos critique. Mais comprendre un texte ne suffit pas : l'épreuve du résumé exige de le \emph{reconstruire} en le condensant, sans le trahir. C'est tout l'objet de cette section : passer de l'analyse à la production, en maîtrisant les gestes techniques de la reformulation et de la condensation.

\courssub{Reformuler : changer les mots et la structure sans changer le sens}

Reformuler, ce n'est pas remplacer un mot par un synonyme au hasard. C'est une \emph{opération chirurgicale} sur la phrase : on en garde le sens exact (la vérité propositionnelle), mais on en modifie l'enveloppe linguistique (lexique, syntaxe, voix, modalité). L'objectif est double : \textbf{condenser} (gagner de la place) et \textbf{neutraliser} (effacer la subjectivité de l'auteur pour ne garder que l'information).

\begin{methode}[Reformuler une phrase : la démarche en 4 temps]
\begin{enumerate}
    \item \textbf{Isoler l'idée noyau} : quelle est l'information indispensable ? (Sujet + verbe + complément essentiel).
    \item \textbf{Changer le lexique} : substituer des synonymes précis, des termes génériques (hyperonymes) ou des périphrases concises.
    \item \textbf{Transformer la syntaxe} : nominaliser les verbes, passer à la voix passive/active, fusionner deux propositions en une, supprimer les expansions secondaires.
    \item \textbf{Tester l'équivalence} : la phrase reformulée pourrait-elle remplacer l'originale dans le texte sans perte de sens ni contresens ? Si oui, c'est gagné.
\end{enumerate}
\end{methode}

\begin{exemplebox}[Avant / Après sur une phrase d'Oyono]
\textbf{Original (32 mots) :} « Meka, le vieux chef respecté de tout le village, avait donné ses meilleures terres à la mission catholique et avait vu ses deux fils unique mourir à la guerre pour la France. »

\textbf{Reformulé (16 mots) :} « Meka, chef villageois respecté, offrit ses terres à la mission et perdit ses deux fils au front. »

\textbf{Opérations visibles :}
\begin{itemize}
    \item \emph{Synonymie / hyperonymie :} « vieux chef respecté de tout le village » $\to$ « chef villageois respecté » ; « mourir à la guerre pour la France » $\to$ « périr au front ».
    \item \emph{Suppression d'expansions :} « meilleures » (terres), « catholique » (mission), « unique » (fils) — détails non essentiels à l'intrigue.
    \item \emph{Fusion / nominalisation :} « avait donné... et avait vu... mourir » $\to$ « offrit... et perdit » (deux verbes conjugués au lieu de trois formes verbales).
    \item \emph{Changement de voix / temps :} passé composé / plus-que-parfait $\to$ passé simple (temps du récit).
\end{itemize}
\end{exemplebox}

\begin{popfigure}
\begin{tikzpicture}[
    node distance=1.2cm and 1.5cm,
    box/.style={rectangle, rounded corners, draw=popblue, thick, fill=popblueL, minimum width=3.2cm, minimum height=1.1cm, align=center, font=\small},
    arrow/.style={->, >=Stealth, thick, popdark},
    opbox/.style={ellipse, draw=poporange, thick, fill=poporangeL, minimum width=2.8cm, minimum height=0.9cm, align=center, font=\small\bfseries}
]
% Original sentence
\node[box, align=center] (orig){\textbf{Phrase originale}\\(32 mots)\\\emph{Meka, le vieux chef respecté...}};
% Operations
\node[opbox, right=of orig, xshift=1.5cm] (syn) {Synonymie / Hyperonymie};
\node[opbox, below=of syn, align=center] (sup){Suppression\\(détails secondaires)};
\node[opbox, below=of sup] (fus) {Fusion / Nominalisation};
\node[opbox, below=of fus, align=center] (voi){Changement\\ de voix / temps};
% Arrows from original to operations
\draw[arrow] (orig.east) -- ++(0.5,0) |- (syn.west);
\draw[arrow] (orig.east) -- ++(0.5,0) |- (sup.west);
\draw[arrow] (orig.east) -- ++(0.5,0) |- (fus.west);
\draw[arrow] (orig.east) -- ++(0.5,0) |- (voi.west);
% Result
\node[box, right=of fus, xshift=2.5cm, fill=popgreenL, draw=popgreen, align=center] (res){\textbf{Phrase reformulée}\\(16 mots)\\\emph{Meka, chef villageois respecté,}\\\emph{offrit ses terres à la mission}\\\emph{et perdit ses deux fils au front.}};
\draw[arrow] (syn.east) -- ++(0.5,0) |- (res.north west);
\draw[arrow] (sup.east) -- ++(0.5,0) |- (res.west);
\draw[arrow] (fus.east) -- ++(0.5,0) |- (res.west);
\draw[arrow] (voi.east) -- ++(0.5,0) |- (res.south west);
% Gain label
\node[draw=popgold, thick, fill=popgoldL, rounded corners, below=0.3cm of res, font=\small\bfseries, text=popdark] {Gain : -50\% de mots, sens intégral conservé};
\end{tikzpicture}
\legende{Schéma des opérations de reformulation appliquées à une phrase de \emph{Le Vieux nègre et la médaille}. Chaque flèche représente une transformation syntaxique ou lexicale qui condense l'information.}
\end{popfigure}

\begin{savaistu}[La nominalisation, reine de la condensation]
En français, transformer un verbe en nom (ex. : « il résista » $\to$ « sa résistance » ; « il décida » $\to$ « sa décision ») permet de \textbf{supprimer le sujet, l'auxiliaire et souvent les compléments circonstanciels}, tout en gardant le noyau sémantique. C'est le levier n°1 pour passer d'une phrase de 20 mots à un groupe nominal de 4 mots. Au Probatoire, les correcteurs repèrent immédiatement les candidats qui savent nominaliser : c'est la marque d'une maîtrise syntaxique avancée.
\end{savaistu}

\courssub{Techniques de condensation : du texte au résumé}

La reformulation phrase à phrase ne suffit pas. Le résumé exige une \emph{vision globale} : il faut regrouper, hiérarchiser, généraliser. Trois opérations majeures :

\begin{propriete}[Les trois piliers de la condensation]
\begin{enumerate}
    \item \textbf{Suppression des expansions accessoires} : adjectifs qualificatifs non discriminants, compléments du nom redondants, relatives explicatives, incises.
    \item \textbf{Regroupement des idées voisines} : fusionner plusieurs phrases qui traitent du même thème (même personnage, même lieu, même action) en une seule phrase complexe.
    \item \textbf{Généralisation des énumérations} : remplacer une liste exhaustive par un terme générique (hyperonyme) ou un quantificateur.
    \begin{itemize}
        \item « Il prit son fusil, ses cartouches, son couteau, sa gourde » $\to$ « Il prit son équipement. »
        \item « Les anciens, les notables, les chefs de famille » $\to$ « Les notables du village. »
    \end{itemize}
\end{enumerate}
\end{propriete}

\begin{exoresolu}[Condensation d'un paragraphe complet (extrait fictif inspiré d'Oyono)]
\textbf{Texte original (84 mots) :}
« Le soleil brûlait déjà les toits de tôle quand Meka quitta sa case. Il marcha lentement vers la place du marché, soutenu par sa canne en bois noir. Les enfants jouaient près du grand manguier ; les femmes vendaient des tomates, des oignons, des piments, du manioc. Meka salua le chef de quartier, puis le catéchiste, puis le représentant de l'administration. Il s'assit sur le banc réservé aux anciens, en face du mât où flottait le drapeau tricolore. »

\textbf{Résumé condensé (23 mots) :}
« Sous un soleil ardent, Meka gagne la place du marché, salue les autorités et s'installe sur le banc des anciens, face au drapeau français. »

\textbf{Analyse des opérations :}
\begin{itemize}
    \item \emph{Suppression :} « soutenu par sa canne en bois noir », « les enfants jouaient près du grand manguier », description détaillée des étals (tomates, oignons, piments, manioc $\to$ implicitement « marché »).
    \item \emph{Regroupement :} « salua le chef de quartier, puis le catéchiste, puis le représentant de l'administration » $\to$ « salue les autorités ».
    \item \emph{Généralisation :} énumération des marchandises $\to$ terme générique « marché » ; liste des salutations $\to$ « autorités ».
    \item \emph{Nominalisation implicite :} « Meka quitta sa case... marcha... salua... s'assit » $\to$ « Meka gagne... salue... s'installe » (suite d'actions au présent de narration).
\end{itemize}
\textbf{Taux de compression :} 73\% (cible Probatoire : 20--25\% du texte original).
\end{exoresolu}

\begin{popfigure}
\begin{tikzpicture}[
    node distance=0.8cm and 1cm,
    txt/.style={rectangle, draw=popdark, thick, fill=popcream, minimum width=5.5cm, minimum height=3.2cm, align=left, font=\small},
    arr/.style={->, >=Stealth, thick, poppurple},
    lbl/.style={font=\small\bfseries, poppurple}
]
\node[txt, align=center] (orig){
\textbf{Texte original (84 mots)}\\
\emph{Le soleil brûlait déjà les toits de tôle quand Meka quitta sa case. Il marcha lentement vers la place du marché, soutenu par sa canne en bois noir. Les enfants jouaient près du grand manguier ; les femmes vendaient des tomates, des oignons, des piments, du manioc. Meka salua le chef de quartier, puis le catéchiste, puis le représentant de l'administration. Il s'assit sur le banc réservé aux anciens, en face du mât où flottait le drapeau tricolore.}
};
\node[txt, right=of orig, xshift=1.5cm, fill=popgreenL, draw=popgreen, align=center] (res){
\textbf{Résumé (23 mots)}\\
\emph{Sous un soleil ardent, Meka gagne la place du marché, salue les autorités et s'installe sur le banc des anciens, face au drapeau français.}
};
\draw[arr] (orig.east) -- node[lbl, above, midway] {Condensation -73\%} (res.west);
% Annotations on original
\node[lbl, anchor=north west] at (orig.north west) {$\times$ Détails spatio-temporels};
\node[lbl, anchor=south west] at ($(orig.south west)+(0,0.3)$) {$\times$ Énumérations exhaustives};
\node[lbl, anchor=east] at ($(orig.east)+(-0.1,0.5)$) {$\times$ Liste des salutations};
% Annotations on summary
\node[lbl, anchor=north west] at (res.north west) {$\checkmark$ Généricité};
\node[lbl, anchor=south west] at ($(res.south west)+(0,0.3)$) {$\checkmark$ Regroupement};
\node[lbl, anchor=east] at ($(res.east)+(-0.1,0.5)$) {$\checkmark$ Fusion d'actions};
\end{tikzpicture}
\legende{Exemple annoté : passage de 84 mots à 23 mots. Les croix rouges marquent les suppressions ; les coches vertes signalent les opérations de condensation (généralisation, regroupement, fusion).}
\end{popfigure}

\courssub{Rédiger un texte continu : cohérence, connecteurs, autonomie}

Un résumé n'est pas une liste de phrases reformulées les unes à la suite des autres. C'est un \textbf{texte à part entière}, qui doit pouvoir se lire \emph{sans avoir lu l'original}. Cela impose trois exigences :

\begin{propriete}[Un bon résumé est un texte autonome]
Le lecteur qui n'a jamais ouvert \emph{Le Vieux nègre et la médaille} doit comprendre : \textbf{qui} agit, \textbf{quoi} se passe, \textbf{où} et \textbf{quand}, \textbf{pourquoi} (causalité), et \textbf{comment} cela se termine. Aucune référence implicite (« le vieux », « la médaille », « le chef ») ne doit rester flottante : la première occurrence doit être explicite (« Meka, vieux chef villageois », « la médaille remise par l'administration »).
\end{propriete}

Pour assurer la fluidité et la logique, on utilise des \textbf{connecteurs logiques} (non pas des connecteurs de parole comme « d'ailleurs », « en fait ») :

\begin{center}
\begin{tabularx}{\linewidth}{|l|X|X|}\hline
\textbf{Fonction} & \textbf{Connecteurs recommandés} & \textbf{Exemple dans un résumé d'Oyono} \\ \hline
Ordre chronologique & d'abord, ensuite, puis, enfin, finalement & « D'abord, Meka offre ses terres ; ensuite, il perd ses fils ; enfin, il reçoit la médaille. » \\ \hline
Opposition / concession & cependant, pourtant, néanmoins, mais & « Meka a tout donné ; cependant, l'administration ne lui offre qu'une médaille. » \\ \hline
Cause / conséquence & car, parce que, donc, ainsi, par conséquent & « Il a servi la France ; par conséquent, il attend une juste récompense. » \\ \hline
Ajout / énumération & de plus, par ailleurs, également & « Il perd ses fils ; de plus, ses terres sont confisquées. » \\ \hline
Conclusion / synthèse & en définitive, au total, en somme & « En définitive, le contrat colonial est rompu. » \\ \hline
\end{tabularx}
\end{center}

\begin{attention}[Pièges de la rédaction continue]
\begin{itemize}
    \item \textbf{Plan apparent :} « I. Meka donne ses terres. II. Ses fils meurent. III. Il reçoit la médaille. » $\to$ \textbf{Interdit.} Pas de chiffres romains, pas de « Première partie », pas de sauts de ligne artificiels.
    \item \textbf{Puces / tirets :} « - Meka offre ses terres. - Ses fils meurent. » $\to$ \textbf{Interdit.} Phrases complètes, liées par connecteurs.
    \item \textbf{Références floues :} « Le vieux attend sa récompense. » (Qui ? Quel vieux ? Quelle récompense ?) $\to$ \textbf{Correction :} « Meka, le vieux chef, attend la médaille promise. »
    \item \textbf{Commentaire personnel :} « C'est injuste. » « On comprend sa colère. » $\to$ \textbf{Interdit.} Voir boîte \textsc{Attention} ci-dessous.
\end{itemize}
\end{attention}

\courssub{Respecter la personne et le ton : l'impersonnalité obligatoire}

Le résumé est un \emph{texte de l'ombre} : l'auteur du résumé s'efface totalement derrière l'auteur du texte source.

\begin{definition}[Règles d'impersonnalité au résumé]
\begin{enumerate}
    \item \textbf{Troisième personne obligatoire :} « Meka reçoit une médaille », jamais « Je vois que Meka reçoit... » ou « On voit que... ».
    \item \textbf{Présent de narration (ou temps du texte) :} Si le texte est au passé simple, on garde le passé simple ; si le texte mêle les temps, on harmonise au présent de narration (usage journalistique / académique) : « Meka offre ses terres, perd ses fils, reçoit une médaille. »
    \item \textbf{Zéro modalité subjective :} pas de verbes d'opinion (trouver, penser, croire, estimer), pas d'adjectifs évaluatifs (injuste, cruel, mérité, scandaleux), pas d'exclamations, pas de questions rhétoriques.
    \item \textbf{Fidélité au point de vue du texte :} si le texte est ironique, le résumé \emph{rapporte} l'ironie sans l'adopter : « L'administration lui remet \emph{avec solennité} une \emph{simple} médaille » (l'adverbe et l'adjectif sont du texte, pas du résumeur).
\end{enumerate}
\end{definition}

\begin{attention}[Erreur fatale : le « je » et l'avis personnel]
\begin{itemize}
    \item $\times$ « \textbf{Je trouve} que Meka a raison de refuser la médaille. »
    \item $\times$ « \textbf{À mon avis}, l'administration est hypocrite. »
    \item $\times$ « \textbf{C'est révoltant} de voir ce vieux traité ainsi. »
    \item $\checkmark$ « Meka refuse la médaille que lui tend l'administration. »
    \item $\checkmark$ « L'administration remet une médaille à Meka, après qu'il a donné ses terres et perdu ses fils. »
\end{itemize}
\textbf{Conséquence au Probatoire :} toute trace de subjectivité (je, à mon avis, c'est inadmissible, etc.) entraîne une \textbf{pénalité lourde} (souvent 3 à 5 points sur 20) car elle prouve que le candidat n'a pas compris la nature même de l'exercice.
\end{attention}

\courssub{Compter les mots et ajuster : la contrainte quantitative}

Au Probatoire / Baccalauréat technique, le résumé est \textbf{calibré} : généralement \textbf{20 à 25\% du nombre de mots du texte original} (souvent 100--120 mots pour un texte de 450--500 mots). La marge de tolérance est de $\pm 10\%$.

\begin{methode}[Compter et ajuster : le protocole d'examen]
\begin{enumerate}
    \item \textbf{Compter le texte original} : mots pleins + mots outils (tout ce qui est séparé par un espace). Noter $N_{orig}$.
    \item \textbf{Calculer la cible} : $N_{cible} = 0,22 \times N_{orig}$ (22\% = milieu de l'intervalle). Ex. : $N_{orig}=480 \to N_{cible}=106$ mots. Tolérance : $[95 ; 117]$.
    \item \textbf{Rédiger un premier jet} sans compter, en visant la concision naturelle.
    \item \textbf{Compter le jet} : $N_{jet}$.
    \item \textbf{Ajuster} :
    \begin{itemize}
        \item \textbf{Si $N_{jet} > 1,1 \times N_{cible}$ (dépassement) :} couper les adjectifs restants, fusionner deux phrases en une, remplacer une relative par un participe présent ou un GN, supprimer un exemple gardé par erreur.
        \item \textbf{Si $N_{jet} < 0,9 \times N_{cible}$ (insuffisance) :} vérifier qu'aucune idée \emph{essentielle} (personnage principal, action majeure, cause/conséquence clé, chute) n'a été oubliée ; développer la phrase qui la porte en ajoutant le complément circonstanciel ou le complément d'objet manquant.
    \end{itemize}
    \item \textbf{Recopier proprement} la version finale, recompter une dernière fois, écrire le nombre de mots en marge.
\end{enumerate}
\end{methode}

\begin{exemplebox}[Ajustement réel : dépassement puis coupe]
\textbf{Cible :} 100 mots ($\pm 10$). \textbf{Premier jet :} 128 mots.
\textbf{Phrase en trop :} « Meka, qui est un vieux chef très respecté de son village, avait donné ses terres à la mission. » (19 mots)
\textbf{Coupure :} « Meka, chef villageois respecté, offrit ses terres à la mission. » (11 mots) \textbf{Gain : -8 mots.}
\textbf{Nouveau total :} 120 mots. Encore 10 de trop.
\textbf{Deuxième coupe :} « Il salua le chef, le catéchiste et l'administrateur. » (11 mots) $\to$ « Il salua les autorités. » (4 mots) \textbf{Gain : -7 mots.}
\textbf{Total :} 113 mots. Dans la fourchette. \textbf{Validé.}
\end{exemplebox}

\courssub{Production guidée : du collectif à l'individuel}

L'apprentissage du résumé suit une \emph{progressivité étayée} : modélisation collective $\to$ entraînement guidé $\to$ autonomie.

\begin{experience}[Séance type : rédaction collective au tableau (extrait d'Oyono)]
\textbf{Matériel :} Extrait de 120 mots (début du roman : arrivée de l'administrateur, remise de la médaille, réaction de Meka). Tableau, feutres trois couleurs (noir = texte original, bleu = idées retenues, rouge = reformulation).

\textbf{Protocole (25 min) :}
\begin{enumerate}
    \item \textbf{Lecture commune} (2 min) : élève lit à voix haute ; classe suit.
    \item \textbf{Repérage des idées essentielles} (5 min) : à main levée, élèves proposent « l'info à garder ». Professeur surligne au bleu au tableau. Objectif : 5--6 idées noyau.
    \item \textbf{Reformulation collective} (12 min) : pour chaque idée, un élève propose une phrase reformulée ; classe critique (synonyme juste ? structure plus courte ? sens gardé ?) ; professeur écrit au rouge la version validée.
    \item \textbf{Chaînage} (4 min) : ajouter les connecteurs logiques entre les phrases rouges ; vérifier l'autonomie (premières mentions explicites ?).
    \item \textbf{Comptage} (2 min) : un élève compte les mots du résumé rouge ; ajustement si besoin (coupe / développement).
\end{enumerate}

\textbf{Observations attendues :}
\begin{itemize}
    \item Les élèves repèrent spontanément les expansions à couper (« très », « grand », « solennellement »).
    \item La nominalisation (« sa résistance » au lieu de « il résista ») émerge naturellement après 2--3 exemples.
    \item Le piège du « je » est évité car le professeur rappelle la règle à la première occurrence.
\end{itemize}

\textbf{Interprétation :} Cette modélisation rend visible le processus invisible de l'expert. Les élèves intègrent le \emph{dialogue interne} du résumeur : « Est-ce essentiel ? Comment le dire plus court ? Quel connecteur ici ? »
\end{experience}

\begin{exercice}[Production individuelle : second extrait (100 mots env.)]
\textbf{Consigne :} « Résumez le texte suivant en \textbf{22 à 26 mots} (texte de 110 mots). Respectez les règles : 3\textsuperscript{e} personne, présent de narration, connecteurs, autonomie, zéro avis. Écrivez le nombre de mots en fin de résumé. »

\textbf{Texte :}
« Le lendemain, au lever du soleil, Meka se rendit à la résidence de l'administrateur. Il portait son plus beau pagne et sa canne sculptée. Le gardien lui demanda d'attendre sous le manguier. Une heure plus tard, l'administrateur apparut, souriant, une petite boîte en velours à la main. Il épingla la médaille sur la poitrine de Meka, lui serra la main et prononça quelques mots en français que Meka ne comprit pas. Meka rentra chez lui, la médaille au cou, silencieux. »

\textbf{Indications pour le corrigé (enseignant) :}
\begin{itemize}
    \item Idées essentielles : (1) Meka va chez l'administrateur ; (2) attente ; (3) remise médaille ; (4) discours incompris ; (5) retour silencieux.
    \item Références à expliciter : « l'administrateur » (première mention), « la médaille » (objet central).
    \item Connecteurs : ensuite / puis / cependant / enfin.
\end{itemize}
\end{exercice}

\begin{corrige}[Exercice n°1 -- Production individuelle]
\textbf{Résumé type (24 mots) :}
« Le lendemain, Meka se rend chez l'administrateur, attend, reçoit une médaille lors d'une cérémonie brève dont il ne comprend pas le discours, et rentre silencieux. (24 mots) »

\textbf{Analyse du corrigé :}
\begin{itemize}
    \item \textbf{Autonomie :} « Meka », « l'administrateur », « une médaille » explicités dès la première occurrence.
    \item \textbf{Condensation :} « au lever du soleil » $\to$ « Le lendemain » ; « plus beau pagne et canne sculptée » $\to$ supprimé (détail vestimentaire) ; « gardien », « manguier », « une heure », « boîte en velours », « serra la main » $\to$ supprimés ; « quelques mots en français que Meka ne comprit pas » $\to$ « dont il ne comprend pas le discours » (nominalisation + relatif).
    \item \textbf{Connecteurs :} « et » (ajout), « dont » (relatif = cause/conséquence implicite), « et » (chronologie).
    \item \textbf{Impersonnalité :} 3\textsuperscript{e} personne, présent de narration, zéro modalité.
    \item \textbf{Comptage :} 24 mots $\in [22; 26]$. \textbf{Validé.}
\end{itemize}
\end{corrige}

\begin{aretenir}[Check-list finale avant de rendre la copie]
\begin{itemize}
    \item [ ] Nombre de mots dans la fourchette ($\pm 10\%$ de la cible) \textbf{écrit en marge}.
    \item [ ] Texte \textbf{continu} : pas de plan, pas de puces, phrases liées par connecteurs logiques.
    \item [ ] \textbf{Autonomie} : toute référence (personnage, lieu, objet) explicite à sa première occurrence.
    \item [ ] \textbf{Impersonnalité} : 3\textsuperscript{e} personne, présent de narration (ou temps du texte), aucun « je », aucun avis, aucune émotion du résumeur.
    \item [ ] \textbf{Fidélité} : idées essentielles toutes présentes, dans l'ordre logique/chronologique du texte, sans contresens.
    \item [ ] \textbf{Reformulation visible} : synonymes, nominalisations, fusions, suppressions — pas de copier-coller de segments du texte original (sauf termes techniques / noms propres).
    \item [ ] \textbf{Orthographe / syntaxe} : relecture à l'envers (dernière phrase $\to$ première) pour traquer les accords, les homophones, la ponctuation.
\end{itemize}
\end{aretenir}

La section suivante — \emph{Amélioration du résumé et bilan de l'œuvre} — nous fera travailler sur la \emph{révision raisonnée} (auto-évaluation, correction par les pairs, grille critériée) et sur la \emph{synthèse globale} de \emph{Le Vieux nègre et la médaille} : thèmes, personnages, portée historique et littéraire, pour boucler l'unité dans sa cohérence.

\coursec{Amélioration du résumé et bilan de l'œuvre}

\courssub{Transition : de la production à la perfection}

Après avoir rédigé une première version de votre résumé (section précédente), vous détenez une matière brute. Comme le sculpteur qui vient d'extraire la forme grossière du bloc de marbre, vous devez maintenant polir, affiner, vérifier chaque arête. Cette étape finale — l'amélioration — est celle qui sépare le travail « correct » du travail « excellent », celui qui obtient la note maximale au Probatoire. C'est aussi le moment de refermer le livre : nous ferons le bilan complet de \emph{Le Vieux nègre et la médaille}, pour en retenir la portée littéraire, historique et humaine.

\courssub{La relecture critique : trois passes pour un texte parfait}

Ne relisez jamais votre résumé d'une seule traite en espérant tout voir. L'œil humain ne peut pas contrôler simultanément le sens, la langue et la forme. La méthode professionnelle impose \textbf{trois passes distinctes}, chacune avec un objectif unique.

\begin{methode}[La relecture en 3 passes]
\textbf{1re passe — Le sens (fidélité et exhaustivité).} Lisez uniquement pour vérifier le contenu. Posez-vous trois questions :
\begin{enumerate}
    \item Ai-je déformé une idée du texte original ? (Fidélité)
    \item Ai-je oublié une idée \cle{essentielle} (nœud narratif, argument majeur, thèse) ? (Exhaustivité)
    \item L'ordre des idées respecte-t-il la logique du texte source ? (Cohérence)
\end{enumerate}
\textbf{2e passe — La langue (correction et style).} Ne regardez plus le texte source. Concentrez-vous sur votre phrase.
\begin{enumerate}
    \item Orthographe lexicale et grammaticale (accords, homonymes : \emph{leur/leurs}, \emph{a/à}, \emph{et/est}).
    \item Ponctuation : les virgules découpent-elles correctement les propositions ? Les points séparent-ils bien les phrases ?
    \item Style : ai-je évité les répétitions (mots, structures) ? Les lourdeurs (trop de verbes \emph{être/avoir}, trop de passifs) ? Le vocabulaire est-il précis (lexique spécialisé le cas échéant) ?
\end{enumerate}
\textbf{3e passe — Le format (contrainte d'examen).}
\begin{enumerate}
    \item Comptez les mots \textbf{un par un} (mots outils compris).
    \item Vérifiez l'indication finale : \emph{[Nombre de mots : X]}.
    \item Vérifiez la présentation : paragraphe unique, pas de citation, pas de « je », pas de commentaire personnel.
\end{enumerate}
\end{methode}

\begin{attention}[Erreur fatale à l'examen]
Négliger l'indication du nombre de mots à la fin du résumé. Elle est \textbf{exigée} au Probatoire / Baccalauréat technique. Un résumé parfait sans cette mention perd systématiquement des points (souvent 1 à 2 points sur la note de forme). Comptez, recomptez, écrivez-le lisiblement entre crochets.
\end{attention}

\courssub{Correction linguistique et stylistique : les chantiers prioritaires}

\textbf{1. La chasse aux fautes d'accord (le « radar » du correcteur).}\par
Le correcteur du Probatoire traque prioritairement :
\begin{itemize}
    \item L'accord du participe passé (avec \emph{avoir} : COD placé avant ; avec \emph{être} : sujet).
    \item L'accord des adjectifs et pronoms (genre, nombre).
    \item Les homonymes grammaticaux : \emph{ce/se, la/l'a, ou/où, on/ont, son/sont, mes/mais/mets/met}.
\end{itemize}
\emph{Astuce :} Soulignez chaque verbe, identifiez son sujet, vérifiez l'accord. Soulignez chaque participe passé, appliquez la règle.

\textbf{2. La ponctuation au service de la clarté.}\par
Un résumé dense nécessite une ponctuation rigoureuse pour éviter l'ambiguïté.
\begin{itemize}
    \item \textbf{Virgule :} isole une proposition incidente, sépare des éléments d'énumération, marque l'apposition.
    \item \textbf{Point-virgule :} sépare deux propositions indépendantes liées par le sens (évite le « et » répétitif).
    \item \textbf{Deux-points :} annoncent une explication, une énumération ou une citation (interdite ici).
\end{itemize}

\textbf{3. L'élégance stylistique : concision et fluidité.}\par
\begin{itemize}
    \item \textbf{Remplacez les périphrases par des verbes précis :} \emph{« Il est en train de manger » $\to$ « Il mange » ; « Il fait la connaissance de » $\to$ « Il rencontre »}.
    \item \textbf{Évitez le passif systématique :} \emph{« La médaille est remise à Meka par le commandant » $\to$ « Le commandant remet la médaille à Meka »} (plus court, plus vif).
    \item \textbf{Nominalisez avec modération :} \emph{« Le fait qu'il refuse » $\to$ « Son refus »} (gain de mots), mais attention à ne pas créer un « style nominal » lourd et abstrait.
    \item \textbf{Connecteurs logiques variés :} \emph{donc, par conséquent, cependant, en outre, de surcroît, dès lors}... Évitez la chaîne « Et... Et... Et... ».
\end{itemize}

\begin{exoresolu}[Nettoyage d'un brouillon]
\textbf{Brouillon d'élève (62 mots) :}
\emph{« Meka est un vieux nègre qui a bien servi les blancs. Il reçoit une médaille. Il est très content. Mais sa femme Kelara n'est pas d'accord. Elle dit que c'est une honte. À la fin Meka comprend qu'il a été dupé. Il jette la médaille. Il retrouve sa dignité. C'est une critique de la colonisation. [62 mots] »}

\textbf{Analyse des défauts :} Phrases hachées (style télégraphique), répétitions (\emph{Meka, Il, Il, Il}), vocabulaire basique (\emph{bien servi, content, pas d'accord, honte, comprend, jette}), absence de connecteurs logiques, point de vue personnel (\emph{C'est une critique...}).

\textbf{Version améliorée (48 mots) :}
\emph{« Fidèle serviteur de l'administration coloniale, Meka reçoit la médaille de la Légion d'honneur. Si l'ancien combattant exulte, son épouse Kelara y voit une insulte. La cérémonie tourne à la farce : ivre, Meka réalise l'inanité de cette distinction. Il rejette la médaille, recouvrant ainsi sa dignité bafouée par le système assimilationniste. [48 mots] »}

\textbf{Gains :} Subordination (participes présents, relatives), vocabulaire précis (\emph{exulte, insulte, farce, inanité, assimilationniste, recouvrant}), logique cause/effet marquée, suppression du commentaire final explicite (l'ironie du texte parle d'elle-même).
\tcblower
\textbf{Correction détaillée :}
\begin{enumerate}
    \item \emph{Fidèle serviteur...} : apposition détachée (gain de place, élégance).
    \item \emph{Si l'ancien combattant exulte, son épouse...} : concession (Si = Bien que) + opposition implicite.
    \item \emph{La cérémonie tourne à la farce :} verbe fort + nom précis.
    \item \emph{ivre, Meka réalise...} : participe présent adverbial (cause).
    \item \emph{Il rejette la médaille, recouvrant ainsi...} : gérondif (simultanéité/conséquence) + adverbe de liaison.
    \item \emph{dignité bafouée par le système assimilationniste} : participe passé adjectival + complément d'agent + terme conceptuel précis.
\end{enumerate}
\end{exoresolu}

\courssub{Grille d'auto-évaluation et d'évaluation croisée (pair-à-pair)}

Avant de rendre votre copie, notez-vous vous-même, puis échangez avec un camarade. La grille ci-dessous est calquée sur les attendus officiels de l'examen (total 20 points, répartis sur 5 critères).

\begin{center}
\begin{tabularx}{\linewidth}{|>{\centering\arraybackslash}X|>{\centering\arraybackslash}X|>{\centering\arraybackslash}X|>{\centering\arraybackslash}X|>{\centering\arraybackslash}X|}\hline
\rowcolor{popblueL}
\textbf{Critère} & \textbf{4 pts (Excellent)} & \textbf{3 pts (Bon)} & \textbf{2 pts (Moyen)} & \textbf{1 pt (Insuffisant)} \\ \hline
\textbf{Fidélité} & Aucune déformation. Idées respectées à la lettre. & 1 infidélité mineure (détail). & 1 déformation majeure ou plusieurs mineures. & Contenu trahi, hors-sujet partiel. \\ \hline
\textbf{Condensation} & Ratio respecté (1/3 $\pm$ 10\%). Idées fusionnées sans perte. & Ratio respecté $\pm$ 15\%. Quelques redondances. & Ratio non respecté ($>$ 1/2 ou $<$ 1/4). Coupes arbitraires. & Quasi-copie ou résumé squelettique. \\ \hline
\textbf{Cohérence} & Enchaînement fluide, connecteurs variés, logique impeccable. & Liens logiques présents, quelques lourdeurs. & Ruptures logiques, connecteurs absents/erronés. & Incompréhensible, idées juxtaposées. \\ \hline
\textbf{Langue} & Zéro faute. Style élégant, vocabulaire précis. & 1 à 2 fautes mineures. Style correct. & 3 à 5 fautes (dont accords). Style lourd. & $>$ 5 fautes. Style défaillant. \\ \hline
\textbf{Nb de mots} & Indication exacte, présente, bien formatée. & Indication présente, écart $\pm$ 2 mots. & Indication absente ou écart $>$ 5 mots. & Indication absente + non-respect ratio. \\ \hline
\end{tabularx}
\end{center}

\begin{popfigure}
\begin{tikzpicture}[scale=0.9, every node/.style={font=\small}]
% Axe principal
\draw[popdark, thick] (0.5,0) -- (13.5,0);
% Événements
% 1916
\draw[popblue, thick] (1.5, -0.15) -- (1.5, 0.15);
\node[below=0.2cm, popdark, font=\bfseries] at (1.5, -0.4) {1916};
\node[above=0.5cm, align=center, text width=3.5cm, poppurple] at (1.5, 0.5) {Arrivée des Français};
\node[above=2.2cm, align=center, text width=3.5cm, popmuted, font=\footnotesize] at (1.5, 0.5) {Fin de l'allemand, mandat français};
% 1956
\draw[popblue, thick] (7, -0.15) -- (7, 0.15);
\node[below=0.2cm, popdark, font=\bfseries] at (7, -0.4) {1956};
\node[above=0.5cm, align=center, text width=3.5cm, poppurple] at (7, 0.5) {Publication du roman};
\node[above=2.2cm, align=center, text width=3.5cm, popmuted, font=\footnotesize] at (7, 0.5) {Oyono, littérature engagée};
% 1960
\draw[popblue, thick] (12.5, -0.15) -- (12.5, 0.15);
\node[below=0.2cm, popdark, font=\bfseries] at (12.5, -0.4) {1960};
\node[above=0.5cm, align=center, text width=3.5cm, poppurple] at (12.5, 0.5) {Indépendance};
\node[above=2.2cm, align=center, text width=3.5cm, popmuted, font=\footnotesize] at (12.5, 0.5) {Cameroun souverain};
% Zones contextuelles (remplacement des braces par des lignes et étiquettes)
\draw[poporange, thick] (1.5, -1.0) -- (7, -1.0);
\node[poporange, font=\footnotesize\bfseries] at (4.25, -1.4) {Période coloniale française};
\draw[popgreen, thick] (7, -1.0) -- (12.5, -1.0);
\node[popgreen, font=\footnotesize\bfseries] at (9.75, -1.4) {Décolonisation};
% Titre
\node[popdark, font=\bfseries\large] at (7, 3.5) {Frise contextuelle : Le Vieux nègre et la médaille};
\end{tikzpicture}
\legende{Frise chronologique situant l'œuvre (1956) entre l'installation de l'administration française (1916) et l'indépendance (1960).}
\end{popfigure}

\courssub{Inscription de l'œuvre dans son contexte : le Cameroun des années 1950}

Pour réussir la présentation d'œuvre (souvent demandée en introduction de commentaire ou en question de cours), vous devez maîtriser le « carnet d'identité » contextuel.

\begin{exemplebox}[Situer l'œuvre — Modèle de présentation en 5 lignes]
\textbf{Auteur :} Ferdinand Oyono (1929--2010), haut fonctionnaire et diplomate camerounais.
\textbf{Date :} 1956 (année charnière : loi-cadre Defferre, début de l'autonomie interne).
\textbf{Genre :} Roman court / Roman satirique / Roman anti-colonial.
\textbf{Thème central :} Dénonciation de la politique d'assimilation et de la complicité de l'Église missionnaire ; quête de la dignité africaine.
\textbf{Contexte :} Cameroun sous administration française (mandat puis tutelle ONU). Montée des nationalismes (UPC, Ruben Um Nyobé). Littérature de la « négritude » et de la contestation (Mongo Beti, \emph{Le Pauvre Christ de Bomba}, 1956).
\end{exemplebox}

\begin{definition}[Contexte historique : le Cameroun sous administration française (1916--1960)]
Après la défaite allemande (1916), la SDN confie le Cameroun à la France (mandat, puis tutelle ONU en 1946). L'administration met en place le \cle{régime de l'indigénat} (travail forcé, impôt de capitation, justice sommaire). La politique officielle est l'\cle{assimilation} : l'Africain « évolué » peut devenir citoyen français s'il renonce à son statut coutumier. L'Église catholique (missionnaires du Saint-Esprit) est partenaire de l'école et de l'encadrement social. Dans les années 1950, l'Union des Populations du Cameroun (UPC) réclame l'indépendance immédiate ; la répression est féroce (guerre du maquis, 1955--1960). Oyono écrit \emph{dans} cette tension.
\end{definition}

\begin{savaistu}[Oyono et Beti : les deux piliers de la contestation camerounaise]
En 1956 paraissent \textbf{deux} chefs-d'œuvre fondateurs de la littérature camerounaise francophone :
\begin{itemize}
    \item \emph{Le Vieux nègre et la médaille} (Ferdinand Oyono) : satire du vieil assimilé, ironie douce-amère, fin sur le rire libérateur.
    \item \emph{Le Pauvre Christ de Bomba} (Mongo Beti) : roman à thèse, violence de la dénonciation, échec du prêtre missionnaire, tragédie collective.
\end{itemize}
Tous deux publient chez Julliard (collection « Lettres nouvelles »). Ils seront diplomates (Oyono) ou enseignants (Beti) après l'indépendance. Leur œuvre commune a façonné l'imaginaire postcolonial camerounais.
\end{savaistu}

\courssub{Bilan de l'étude de l'œuvre : thèmes, personnage, écriture}

\textbf{1. Thèmes majeurs.}\par
\begin{itemize}
    \item \textbf{La dénonciation de la colonisation :} Non pas par la violence graphique, mais par l'absurde du quotidien (cérémonie ridicule, médailles en chocolat, promesses non tenues). Le système colonial est dévoilé comme une \cle{machine à broyer les hommes} tout en leur faisant croire qu'ils sont honorés.
    \item \textbf{L'illusion de l'assimilation :} Meka \emph{croit} au contrat : « Je donne ma force, ma loyauté, mes fils ; la France me donne l'honneur, la citoyenneté ». La médaille révèle la supercherie : il reste un « sujet », un « vieux nègre », jamais un citoyen. L'assimilation est un \cle{mirage}.
    \item \textbf{La dignité retrouvée :} Le geste final (jeter la médaille dans la latrine) n'est pas une défaite, c'est une \cle{victoire symbolique}. Meka redevient \emph{Meka}, l'ancêtre, le père, l'homme libre dans sa tête. Le rire final (« Il rit, il rit... ») est l'arme absolue du dominé qui a compris le jeu.
\end{itemize}

\textbf{2. Le personnage de Meka : archétype et individu.}\par
Meka n'est pas un héros de roman réaliste psychologique ; c'est un \cle{type} (le vieux combattant, l'évolué, le chef de clan) qui devient \cle{individu} par la prise de conscience.
\begin{itemize}
    \item \textbf{Naïveté lucide :} Il connaît les rouages (il a « mangé » le blanc), mais il veut croire au mythe.
    \item \textbf{Double langage :} Devant le commandant, il joue le rôle (humilité, gratitude) ; en aparté (ou en pensée), il juge, calcule, ironise.
    \item \textbf{Le basculement :} L'ivresse (vin de palme + vin blanc) lève l'inhibition. La vérité jaillit : « \emph{Moi, Meka, j'ai été dupé !} »
\end{itemize}

\textbf{3. L'art de l'ironie chez Oyono : une esthétique de la résistance.}\par
L'ironie n'est pas un ornement, c'est l'\cle{arme} du roman.
\begin{itemize}
    \item \textbf{Ironie situationnelle :} La médaille (symbole suprême) remise dans une cour boueuse, par un commandant ivre, à un vieillard ivre, sous le regard méprisant de l'épouse.
    \item \textbf{Ironie verbale (antiphrase) :} Les discours officiels (commandant, prêtre) disent l'inverse de la réalité. Le narrateur reprend ce lexique pour le vider de son sens.
    \item \textbf{Point de vue décalé (focalisation interne) :} On voit la scène \emph{par les yeux de Meka}. Ses pensées (en italique souvent) contredisent ses paroles. Le lecteur est complice de la lucidité du personnage.
    \item \textbf{Le rire final :} C'est l'ironie suprême. Le colonisateur croit avoir honoré ; l'honoré rit de l'honneur. Le pouvoir symbolique s'effondre.
\end{itemize}

\begin{methode}[Présentation d'œuvre en 5 lignes : compétence Bac]
À l'écrit du Probatoire/Bac, on vous demande souvent : « Présentez brièvement l'œuvre étudiée. » Apprenez par cœur ce canevas (5 phrases types) :
\begin{enumerate}
    \item \emph{Le Vieux nègre et la médaille} est un roman satirique de Ferdinand Oyono, publié en 1956.
    \item Il se déroule au Cameroun sous administration française, dans les années 1950.
    \item Il met en scène Meka, ancien combattant et chef de canton, qui reçoit la Légion d'honneur.
    \item L'œuvre dénonce l'illusion de l'assimilation et la complicité de l'Église coloniale.
    \item Par l'ironie et le rire final, elle célèbre la reconquête de la dignité africaine.
\end{enumerate}
Entraînez-vous à le réciter en moins de 30 secondes.
\end{methode}

\courssub{Compte rendu et remédiation : analyser les copies pour progresser}

La dernière séance de la séquence est consacrée au retour sur les productions des élèves (résumés, commentaires, dissertations). C'est le moment de la \cle{métacognition} : comprendre \emph{pourquoi} on a réussi ou échoué.

\begin{methode}[Séance type de remédiation collective (1h)]
\begin{enumerate}
    \item \textbf{Anonymisation :} Le professeur sélectionne 3 à 4 copies types (Excellent / Moyen / Insuffisant / Hors-sujet). Il les photocopie sans noms.
    \item \textbf{Évaluation par les pairs (grille) :} En binômes, les élèves notent chaque copie avec la grille à 5 critères (voir tableau ci-dessus). Ils doivent justifier chaque note par une citation précise de la copie.
    \item \textbf{Mise en commun :} Le professeur projette les copies (ou les redistribue). Chaque binôme donne sa note et son argumentaire. Le professeur valide, corrige, explicite le barème officiel.
    \item \textbf{Identification des erreurs récurrentes (Top 5) :}
    \begin{enumerate}
        \item Dépassement du nombre de mots (souvent 1/2 au lieu de 1/3).
        \item Présence de « Je pense que », « À mon avis », « Ce texte montre que... » (commentaire au lieu de résumé).
        \item Copie de phrases entières du texte (pas de reformulation).
        \item Oubli d'idées essentielles (ex: le rôle de Kelara, le refus final, l'ironie).
        \item Fautes d'accords majeures (participes passés, sujets complexes).
    \end{enumerate}
    \item \textbf{Réécriture ciblée :} Chaque élève reprend \emph{sa propre} copie et corrige \textbf{un seul} des défauts identifiés (ex: « Je réduis mon texte de 20\% en fusionnant les phrases » OU « Je remplace tous les "il y a" par des verbes d'action »).
    \item \textbf{Bilan personnel écrit :} Sur un carnet de progression : « Mon point fort : ... / Mon point à travailler : ... / Mon action pour la prochaine fois : ... ».
\end{enumerate}
\end{methode}

\begin{aretenir}[Check-list finale avant l'examen (Résumé)]
\begin{itemize}
    \item [ ] J'ai lu le texte 2 fois (1. global, 2. découpage/idées forces).
    \item [ ] J'ai listé les idées essentielles (mots-clés, pas phrases).
    \item [ ] J'ai regroupé/fusionné les idées voisines.
    \item [ ] J'ai rédigé en \textbf{mes mots} (reformulation), à la 3e personne, temps du récit (imparfait/passé simple) ou présent de narration.
    \item [ ] J'ai utilisé des connecteurs logiques variés.
    \item [ ] J'ai relu : 1. Sens (fidélité) 2. Langue (fautes, style) 3. Mots (comptage + indication \emph{[X mots]}).
    \item [ ] Je connais la fiche d'identité de l'œuvre (Auteur, Date, Genre, Thème, Contexte) en 5 lignes.
\end{itemize}
\end{aretenir}

\begin{exercice}[Entraînement final : Remédiation sur un résumé défaillant]
Voici un résumé d'élève (78 mots) pour un texte source de 240 mots (cible : 80 mots $\pm$ 10\% $\to$ 72--88 mots). Le texte source raconte la cérémonie de remise de médaille à Meka, l'opposition de Kelara, l'ivresse, le rejet final.
\emph{« Dans ce texte, on voit Meka qui est content car il va avoir une médaille. Le commandant arrive. Il fait un discours. Meka boit du vin. Kelara sa femme n'est pas contente elle dit que c'est une honte pour les noirs. Meka devient ivre. Il se rend compte que la médaille ne vaut rien. Il la jette dans les toilettes. Il est fier de lui. C'est une critique de la colonisation française. [78 mots] »}

\textbf{Travail demandé :}
\begin{enumerate}
    \item Évaluez ce résumé sur la grille à 5 critères (note sur 20).
    \item Identifiez les 3 défauts majeurs.
    \item Réécrivez-le correctement (cible 75 mots).
\end{enumerate}
\end{exercice}

\begin{corrige}[Exercice n°16 -- Remédiation]
\textbf{1. Évaluation (sur 20) :}
\begin{itemize}
    \item Fidélité : 2/4 (Idées présentes mais déformées : « content » au lieu de fierté complexe ; « critique de la colonisation » explicite = commentaire).
    \item Condensation : 3/4 (78 mots pour cible 80 : OK, mais ratio source/résumé non vérifié ici).
    \item Cohérence : 1/4 (Style télégraphique, « Dans ce texte on voit », « Il fait un discours » vide, pas de connecteurs logiques).
    \item Langue : 2/4 (Fautes : « Kelara sa femme » (virgule), « il est fier » (style), répétitions « Meka/Il/Il/Il »).
    \item Nb mots : 2/4 (Indication absente).
    \item \textbf{Total : 10/20 (Moyen).}
\end{itemize}

\textbf{2. Trois défauts majeurs :}
\begin{enumerate}
    \item \textbf{Commentaire intrusif :} « C'est une critique de la colonisation... » (Interdit en résumé).
    \item \textbf{Absence de reformulation / Style oral :} « Dans ce texte on voit », « Meka qui est content », « Kelara sa femme n'est pas contente ».
    \item \textbf{Pauvreté lexicale et syntaxique :} Verbes faibles (\emph{est, a, dit, devient, est}), pas de subordination, pas d'ironie rendue.
\end{enumerate}

\textbf{3. Proposition de réécriture (76 mots) :}
\emph{« À l'occasion d'une cérémonie officielle, l'administrateur colonial remet la Légion d'honneur à Meka, vieux chef de canton et ancien combattant. Tandis que le récipiendaire savoure cet honneur, son épouse Kelara dénonce une humiliation pour leur race. Grisé par l'alcool, Meka perçoit soudain la vanité de cette distinction : il la jette aux latrines. Par ce geste de rejet, l'ancien loyaliste retrouve sa dignité d'homme libre, tournant en dérision le mythe assimilationniste. [76 mots] »}
\end{corrige}

\courssub{Synthèse de la séquence : ce qu'il faut emporter pour l'examen}

\begin{propriete}[Compétences certificatives (Probatoire / Bac Tech) -- Résumé]
À l'issue de cette unité, vous devez être capable de :
\begin{enumerate}
    \item Lire un texte de 200--300 mots en 10 min et en extraire la structure (thèse / arguments ou intrigue / nœuds).
    \item Sélectionner les idées \textbf{essentielles seulement} (élaguer exemples, descriptions, détails).
    \item Reformuler ces idées dans une langue \textbf{correcte, fluide, impersonnelle}, en respectant l'ordre logique du texte.
    \item Produire un texte de \textbf{longueur imposée} (généralement 1/3 du texte source $\pm$ 10\%), \textbf{en un seul paragraphe}, \textbf{sans citation}, \textbf{sans commentaire}.
    \item Indiquer \textbf{exactement} le nombre de mots utilisés entre crochets à la fin.
    \item Situer une œuvre étudiée (Oyono, Beti, etc.) en 5 lignes : Auteur, Date, Genre, Thème, Contexte.
\end{enumerate}
Ces compétences sont transverses : elles vous serviront en Philosophie (synthèse de textes), en Histoire-Géographie (compte-rendu de document), en EPS (compte-rendu d'activité), et dans toute votre vie professionnelle (notes de synthèse, rapports).
\end{propriete}

\begin{popfigure}
\begin{tikzpicture}[node distance=1.2cm and 1.5cm, every node/.style={font=\small}]
% Nœuds
\node[draw, rounded corners, fill=popblueL, text width=3.5cm, align=center, minimum height=1.8cm] (start) {\textbf{TEXTE SOURCE}\\(200--300 mots)};
\node[draw, rounded corners, fill=poporangeL, text width=3.5cm, align=center, minimum height=1.8cm, right=of start] (lecture) {\textbf{LECTURE ACTIVE}\\(Repérage idées forces\\Découpage en séquences)};
\node[draw, rounded corners, fill=popgreenL, text width=3.5cm, align=center, minimum height=1.8cm, right=of lecture] (selection) {\textbf{SÉLECTION / FUSION}\\(Liste mots-clés\\Regroupement logique)};
\node[draw, rounded corners, fill=poppurpleL, text width=3.5cm, align=center, minimum height=1.8cm, right=of selection] (redaction) {\textbf{RÉDACTION}\\(Reformulation\\Connecteurs\\Impersonnel\\1 paragraphe)};
\node[draw, rounded corners, fill=poppinkL, text width=3.5cm, align=center, minimum height=1.8cm, right=of redaction] (controle) {\textbf{AUTO-CONTRÔLE 3 PASSES}\\(1. Sens / 2. Langue / 3. Mots\\Nb mots : X)};
\node[draw, rounded corners, fill=popgoldL, text width=3.5cm, align=center, minimum height=1.8cm, below=of redaction] (final) {\textbf{RÉSUMÉ FINAL}\\(Prêt pour l'examen)};
% Flèches
\draw[->, thick, popblue] (start) -- (lecture) node[midway, above] {10 min};
\draw[->, thick, poporange] (lecture) -- (selection) node[midway, above] {5 min};
\draw[->, thick, popgreen] (selection) -- (redaction) node[midway, above] {15 min};
\draw[->, thick, poppurple] (redaction) -- (controle) node[midway, above] {10 min};
\draw[->, thick, poppink] (controle) -- (final) node[midway, right] {Validé};
\draw[->, thick, dashed, popmuted] (redaction) |- (final) node[midway, left] {Si erreur : retour};
\end{tikzpicture}
\legende{Organigramme de la méthode du résumé : gestion du temps (40 min type examen) et boucles de rétroaction.}
\end{popfigure}

\textbf{Le mot de la fin.}\par Vous avez maintenant entre les mains la méthode complète, les outils de langue (constatif/performatif, lexique spécialisé, ironie/satire), la connaissance de l'œuvre référence (\emph{Le Vieux nègre et la médaille}) et la grille d'excellence. Le résumé n'est pas un exercice scolaire artificiel : c'est l'art de \textbf{distinguer l'essentiel du bruit}, de \textbf{restituer la pensée d'autrui sans la trahir}, et de \textbf{maîtriser sa langue} sous contrainte. C'est une compétence de pouvoir. Travaillez-la. Elle ne vous quittera plus.

\newpage

\coursec{Activité d'intégration}

\coursec{Leçon 16 : Rédiger un résumé de texte — Activité d'intégration}

\courssub{Objectifs de l'activité}

À l'issue de cette activité d'intégration, l'élève sera capable de :
\begin{itemize}
\item lire un texte de presse à tonalité ironique et en dégager le \cle{thème}, la \cle{thèse} et la \cle{tonalité} ;
\item distinguer les \cle{idées essentielles} des idées secondaires (exemples, répétitions, anecdotes) ;
\item construire un \cle{plan-squelette} fidèle à la progression du texte ;
\item produire un \cle{résumé} conforme aux exigences du Baccalauréat technique : respect de la longueur imposée (125 mots $\pm$ 10 \%), reformulation personnelle, objectivité, cohérence et correction de la langue ;
\item évaluer la production d'un pair à l'aide d'une \cle{grille critériée} et améliorer sa propre copie.
\end{itemize}

\courssub{Situation}

Le club de journalisme du lycée technique de Douala prépare son dossier de presse de fin d'année. La rédaction a reçu de \emph{Cameroon Tribune} l'article ci-dessous (environ 500 mots), relatant une cérémonie de remise de distinctions à d'anciens combattants. La tonalité ironique de l'article n'est pas sans rappeler \emph{Le Vieux nègre et la médaille} de Ferdinand Oyono. Pour le dossier, chaque article doit être accompagné d'un résumé de 125 mots ($\pm$ 10 \%, soit entre 113 et 137 mots).

\begin{exemplebox}[Texte à résumer : « Des médailles pour des oublis »]
\textbf{Des médailles pour des oublis}

Douala, quartier Bonabéri. Sous un chapiteau dressé pour l'occasion, une cinquantaine d'anciens combattants attendent depuis sept heures du matin. Il est onze heures passées lorsque les officiels arrivent enfin, précédés d'une fanfare qui couvre à peine le murmure des estomacs vides.

La cérémonie peut commencer. Le discours d'ouverture célèbre « ces héros qui ont donné leur jeunesse pour la patrie ». Héros, certes, mais héros à qui l'on n'avait pas adressé la parole depuis la dernière commémoration, il y a exactement un an, au même endroit, sous le même chapiteau.

Monsieur Etoundi, quatre-vingt-trois ans, avance en s'appuyant sur sa canne. On lui épingle au revers de sa veste râpée une médaille étincelante. Il sourit. Que pèse une médaille, se demandera-t-on plus tard, face à une retraite qui n'arrive pas, face à un dossier de pension « en cours de traitement » depuis 1998 ? Monsieur Etoundi, lui, ne se demande rien. Il serre sa médaille comme d'autres serrent un billet de banque.

Le maître de cérémonie égrène les noms. À chaque nom, des applaudissements nourris ; à chaque nom, un flash de photographe. Demain, les journaux publieront les images. Demain, tout le monde aura oublié les noms.

Dans son allocution, l'autorité promet que « l'État n'oublie pas ses enfants ». La phrase est belle. Elle l'est d'autant plus qu'elle revient chaque année, fidèle comme la médaille, immobile comme le dossier de pension. Un vieillard, au troisième rang, s'est assoupi. Son voisin le secoue : « Réveille-toi, on distribue la gloire. »

Vient le moment tant attendu du banquet. Enfin, presque : le traiteur n'a livré que la moitié des plateaux prévus. On se partage donc les sandwichs, avec la même discipline qu'autrefois on se partageait les rations. L'analogie fait sourire les anciens ; elle devrait faire rougir les organisateurs.

Interrogé à la sortie, monsieur Etoundi déclare être « très honoré ». Sa veste porte désormais trois médailles : une pour 1960, une pour 1995, une pour aujourd'hui. Trois médailles, zéro pension. Le compte est vite fait, même pour un vieux soldat qui a appris les mathématiques sur les champs de bataille.

La fanfare reprend. Le chapiteau sera démonté avant la nuit. Les officiels regagnent leurs voitures climatisées. Les anciens combattants, eux, attendent le car qui doit les ramener. Il arrivera avec deux heures de retard, mais ils sont habitués : toute leur vie n'a été qu'une longue attente, ponctuée de courtes cérémonies.

On dira que la patrie est reconnaissante. La reconnaissance, à Douala, se mesure en grammes de métal doré et en minutes de discours. Elle ne se mesure jamais en francs versés.

Pourtant, personne ne proteste. Les anciens combattants rentrent chez eux, la médaille dans la poche, la dignité sur le visage. Ils reviendront l'an prochain, s'ils sont encore de ce monde, recevoir une médaille de plus et une promesse de moins.

Ainsi va la reconnaissance nationale : éclatante en surface, creuse en dessous, comme ces médailles que l'on épingle sur des vestes que l'on ne remplacera jamais.
\end{exemplebox}

\courssub{Consignes}

\textbf{Travail de groupe (étapes 1 à 3), puis travail individuel (étape 4), puis évaluation croisée (étapes 5 et 6).}

\begin{enumerate}
\item \textbf{Lecture globale.} Lisez l'article deux fois. Dégagez le thème, la thèse implicite de l'auteur et la tonalité dominante. Justifiez votre réponse par trois procédés précis (ironie, antithèse, hyperbole, gradation...).
\item \textbf{Repérage des idées essentielles.} Soulignez dans le texte les idées essentielles (4 à 6 idées maximum). Barrez les exemples, les anecdotes et les répétitions. Que faites-vous du cas de monsieur Etoundi : idée essentielle ou illustration ?
\item \textbf{Plan-squelette.} Rédigez en une phrase chacune les grandes étapes du texte (situation initiale, déroulement, point de vue de l'auteur, conclusion).
\item \textbf{Rédaction individuelle.} Rédigez un résumé de 125 mots ($\pm$ 10 \%) : phrases complètes, présent de vérité générale, aucune citation, aucun jugement personnel, reformulation de toutes les idées essentielles. Indiquez le nombre de mots à la fin.
\item \textbf{Évaluation croisée.} Échangez les résumés au sein du groupe et notez la copie reçue sur 20 à l'aide de la grille critériée ci-dessous. Rédigez un commentaire justificatif de deux lignes.
\item \textbf{Amélioration.} À partir des remarques reçues, réécrivez une version définitive de votre résumé. La meilleure production de chaque groupe sera lue en classe et comparée au corrigé du professeur.
\end{enumerate}

\begin{center}
\begin{tabularx}{\linewidth}{|X|c|}
\hline
\rowcolor{popblueL} \textbf{Critère d'évaluation} & \textbf{Barème} \\
\hline
Fidélité au texte (aucune idée ajoutée, aucune déformation) & /4 \\
\hline
Présence de toutes les idées essentielles & /4 \\
\hline
Reformulation personnelle (pas de phrases recopiées) & /4 \\
\hline
Cohérence et enchaînement logique & /3 \\
\hline
Respect de la longueur (113 à 137 mots) & /2 \\
\hline
Correction de la langue (orthographe, grammaire, ponctuation) & /3 \\
\hline
\textbf{Total} & \textbf{/20} \\
\hline
\end{tabularx}
\end{center}

\courssub{Ressources à mobiliser}

\begin{aretenir}[Boîte à outils du résumé]
\begin{itemize}
\item \textbf{Méthode en cinq étapes :} lire $\rightarrow$ dégager thème et thèse $\rightarrow$ sélectionner les idées essentielles $\rightarrow$ construire le plan-squelette $\rightarrow$ rédiger, compter, relire.
\item \textbf{Idée essentielle :} idée dont la suppression rend le texte incompréhensible. L'exemple, l'anecdote (monsieur Etoundi), la répétition sont éliminés ou fondus dans l'idée générale.
\item \textbf{Reformulation :} remplacer les mots du texte par des synonymes ou des périphrases, condenser une phrase en un groupe nominal, transformer le concret en général.
\item \textbf{Énoncé constatif et performatif :} repérer les énoncés performatifs du texte (« l'État n'oublie pas ses enfants », promesse) et les neutraliser dans le résumé par un constat objectif (« l'autorité promet que... »).
\item \textbf{Lexique spécialisé :} conserver les termes techniques nécessaires au sens (médaille, pension, ancien combattant, commémoration) ; ils ne se reformulent pas.
\item \textbf{Tonalités satirique et polémique :} l'ironie de l'article doit être signalée dans le résumé (« l'auteur dénonce avec ironie... ») sans être imitée : le résumé reste objectif.
\item \textbf{Contraintes formelles :} présent de vérité générale, phrases complètes, pas de citation, pas d'avis personnel, nombre de mots indiqué.
\end{itemize}
\end{aretenir}

\begin{attention}[Pièges fréquents]
Recopier des phrases entières du texte ; donner son avis (« je trouve que... ») ; garder l'anecdote d'Etoundi comme si elle était une idée ; reproduire l'ironie au lieu de la signaler ; oublier de compter les mots ; dépasser la marge des 10 \%.
\end{attention}

\courssub{Démarche et corrigé commenté}

\begin{exoresolu}[Production guidée du résumé]
Énoncé : produire le résumé de l'article « Des médailles pour des oublis » en 125 mots ($\pm$ 10 \%).
\tcblower
\textbf{Étape 1 — Thème, thèse, tonalité.} Thème : la reconnaissance de la patrie envers les anciens combattants. Thèse implicite : cette reconnaissance est purement symbolique et creuse, car elle ne s'accompagne d'aucune aide concrète (pensions impayées). Tonalité : ironique et satirique (antithèses médaille/pension, hyperbole « toute leur vie n'a été qu'une longue attente », formule ironique « on distribue la gloire »).

\textbf{Étape 2 — Idées essentielles retenues (5) :}
\begin{enumerate}
\item une cérémonie de remise de médailles à des anciens combattants se déroule dans la négligence (retards, banquet réduit) ;
\item les discours officiels célèbrent des héros que l'on oublie le reste de l'année ;
\item les promesses de l'État se répètent sans jamais se concrétiser (pensions non versées) ;
\item les anciens combattants accueillent ces honneurs avec dignité et sans protester ;
\item l'auteur dénonce une reconnaissance purement symbolique, éclatante en surface mais vide.
\end{enumerate}
Le cas d'Etoundi est une \emph{illustration} : il sera fondu dans les idées 2 et 3.

\textbf{Étape 3 — Plan-squelette :} cérémonie négligente $\rightarrow$ éloges officiels contrastant avec l'oubli $\rightarrow$ promesses non tenues $\rightarrow$ dignité silencieuse des anciens $\rightarrow$ dénonciation ironique d'une reconnaissance creuse.

\textbf{Étape 4 — Résumé rédigé (134 mots) :}

\emph{L'article relate, sur un ton ironique et satirique, une cérémonie de remise de médailles à des anciens combattants à Douala, marquée par un retard considérable et une organisation négligente, symbolisée par un banquet réduit. Les discours officiels célèbrent avec emphase des héros que l'on ignore délibérément le reste de l'année. Les promesses de l'État, renouvelées chaque année avec une fidélité dérisoire, ne se concrétisent jamais : les pensions restent impayées et les dossiers administratifs sans suite depuis des décennies. Malgré cette indifférence institutionnelle, les anciens combattants accueillent ces honneurs vides avec une dignité silencieuse, sans protester, prêts à revenir l'an prochain. L'auteur dénonce ainsi une reconnaissance nationale purement symbolique et creuse : éclatante dans les discours et le métal doré des médailles, elle reste vide de toute aide concrète envers ceux qu'elle prétend honorer.} (134 mots)

\textbf{Étape 5 — Vérification critériée :} fidélité (aucune idée ajoutée) ; les 5 idées essentielles présentes ; aucune phrase recopiée ; connecteurs logiques (« pourtant », « malgré cela », « ainsi ») ; 134 mots, dans la marge (113-137) ; tonalité signalée (« sur un ton ironique et satirique », « l'auteur dénonce ») sans être imitée ; langue correcte. Note attendue : 20/20.
\end{exoresolu}

\courssub{Bilan de l'activité}

Cette activité a permis de mettre en évidence trois acquis majeurs. D'abord, le résumé n'est pas un raccourci mécanique : il exige une \cle{lecture analytique} capable de hiérarchiser les idées et de distinguer l'essentiel de l'illustration, même lorsque celle-ci est attachante comme le personnage d'Etoundi. Ensuite, la présence d'une \cle{tonalité ironique} — comme dans \emph{Le Vieux nègre et la médaille} — impose un traitement particulier : le résumé signale l'ironie sans la reproduire, car il doit rester objectif et neutre. Enfin, l'\cle{évaluation croisée} et la réécriture montrent que le résumé est un texte perfectible : la relecture critique, guidée par une grille, transforme une première version en une production conforme aux exigences du Baccalauréat technique. Ces compétences sont directement transférables à la vie courante : rédiger un compte rendu, condenser un rapport, informer sans déformer.

\newpage

\coursec{Méthodes et exercices résolus}

\courssub{Méthode 1 : Appliquer les notions de l'unité — repérer les idées essentielles d'un texte}

\begin{methode}[Repérer les idées essentielles d'un texte]
\begin{enumerate}
\item Lire le texte \textbf{deux fois} : une première lecture globale pour saisir le thème général, une seconde lecture active (crayon à la main) pour souligner.
\item Repérer la \cle{structure} du texte : introduction, développement (articulations logiques : d'abord, ensuite, enfin ; mais, cependant, donc), conclusion.
\item Dégager l'\cle{idée principale} de chaque paragraphe ou groupe de paragraphes : c'est souvent la phrase qui résume les autres ou que les exemples illustrent.
\item Éliminer ce qui est secondaire : les \cle{exemples}, les répétitions, les anecdotes, les comparaisons, les citations, les précisions chiffrées.
\item Noter au brouillon le \cle{plan des idées} sous forme de titres courts (une phrase par étape).
\item Vérifier que les idées retenues, mises bout à bout, restituent fidèlement la pensée de l'auteur \textbf{sans rien ajouter}.
\end{enumerate}
\textbf{Réflexe :} se poser pour chaque phrase la question : « Si je supprime cette phrase, le sens général du texte disparaît-il ? » Si non, elle est secondaire.
\end{methode}

\begin{exoresolu}[Repérage des idées essentielles dans un extrait]
\textbf{Énoncé.} On donne l'extrait suivant, inspiré de la thématique du \emph{Vieux nègre et la médaille} de Ferdinand Oyono :

« Meka avait reçu une médaille. Une médaille ! Lui, le vieux nègre, décoré par les Blancs pour avoir donné ses terres à la mission. Tout le village l'enviait. Ses amis le félicitaient, les femmes chantaient, on dansa toute la nuit. Pourtant, au fond de lui-même, Meka sentait confusément que cette médaille ne changerait rien à sa vie : ses terres étaient perdues, ses fils étaient morts à la guerre, et il restait seul avec sa vieille épouse. »

Question : dégagez les idées essentielles de ce passage en trois phrases maximum.

\tcblower
\textbf{Solution détaillée.}

\textbf{Étape 1 — Thème général.} Le passage évoque la médaille reçue par Meka et le contraste entre l'apparence (la fête) et la réalité (le vide de sa vie).

\textbf{Étape 2 — Classement des informations.}
\begin{itemize}
\item Idée essentielle 1 : Meka a été décoré d'une médaille par les Blancs pour avoir cédé ses terres.
\item Idée essentielle 2 : le village célèbre cet honneur (félicitations, chants, danses) — c'est l'admiration collective.
\item Idée essentielle 3 : Meka éprouve pourtant un malaise, car la médaille ne compense pas ses pertes (terres, fils).
\end{itemize}

\textbf{Étape 3 — Éléments secondaires à éliminer.} Les détails de la fête (« les femmes chantaient, on dansa toute la nuit ») sont des illustrations de l'idée 2 ; la mention de la vieille épouse est une précision accessoire.

\textbf{Étape 4 — Rédaction des idées essentielles.}
\begin{enumerate}
\item Meka a reçu une médaille des Blancs en échange de ses terres.
\item Tout le village admire et fête cet honneur.
\item Meka ressent cependant que cette récompense ne compense pas la perte de ses terres et de ses fils.
\end{enumerate}

\textbf{Conclusion.} Ces trois phrases restituent la pensée complète du passage : l'honneur apparent, l'enthousiasme collectif et l'amertume intime du personnage. Aucune idée personnelle n'a été ajoutée.
\end{exoresolu}

\courssub{Méthode 2 : Rédiger et justifier une démarche — reformuler sans copier}

\begin{methode}[Reformuler les idées essentielles]
\begin{enumerate}
\item Changer les \cle{mots} : utiliser des synonymes (par exemple « recevoir » $\rightarrow$ « être gratifié de » ; « donner » $\rightarrow$ « céder »).
\item Changer les \cle{constructions} : transformer une phrase active en passive, une subordonnée en proposition principale, un nom en verbe (« l'admiration du village » $\rightarrow$ « le village admire »).
\item Changer le \cle{point de vue} si nécessaire : passer du récit à la première personne à un résumé à la troisième personne.
\item Regrouper les phrases courtes qui expriment la même idée ; scinder les phrases trop longues.
\item Employer des \cle{connecteurs logiques} pour articuler les idées : d'abord, ensuite, cependant, enfin, c'est pourquoi.
\item Interdits absolus : recopier des phrases entières du texte, garder les exemples, introduire son opinion (« je pense que », « à mon avis »).
\end{enumerate}
\textbf{Formule à retenir :} résumé = mêmes idées + autres mots + ton neutre + longueur imposée.
\end{methode}

\begin{exoresolu}[Reformulation d'un passage]
\textbf{Énoncé.} Reformulez en une ou deux phrases, sans recopier aucune phrase du texte, l'idée suivante extraite d'un texte sur l'éducation :

« L'école technique donne aux jeunes les moyens de gagner leur vie : un mécanicien, un électricien, un couturier diplômé trouvent plus facilement du travail que celui qui n'a aucune formation. »

\tcblower
\textbf{Solution détaillée.}

\textbf{Étape 1 — Idée essentielle.} La formation technique facilite l'insertion professionnelle des jeunes.

\textbf{Étape 2 — Éléments à supprimer.} Les exemples de métiers (mécanicien, électricien, couturier) illustrent l'idée : ils doivent disparaître ou être généralisés (« les techniciens diplômés »).

\textbf{Étape 3 — Reformulation par synonymes et changement de construction.}
\begin{itemize}
\item « donne les moyens de gagner leur vie » $\rightarrow$ « permet d'accéder à l'emploi » ;
\item « trouvent plus facilement du travail » $\rightarrow$ « s'insèrent plus aisément dans la vie active » ;
\item comparaison transformée en cause : « parce qu'ils possèdent une qualification ».
\end{itemize}

\textbf{Étape 4 — Rédaction.}

\emph{Proposition de reformulation :} « La formation technique, en dotant les jeunes d'une qualification professionnelle, leur permet de s'insérer plus aisément dans la vie active que s'ils n'avaient reçu aucun apprentissage. »

\textbf{Vérification.} Aucune phrase du texte n'est recopiée ; l'idée est identique ; le ton est neutre ; les exemples ont disparu.

\textbf{Conclusion.} La reformulation respecte les trois règles du résumé : fidélité au sens, vocabulaire personnel, neutralité du ton.
\end{exoresolu}

\courssub{Méthode 3 : Respecter la contrainte de longueur (contraction)}

\begin{methode}[Contracter un texte à la longueur imposée]
\begin{enumerate}
\item Lire la \cle{consigne} : le sujet impose un nombre de mots (par exemple « environ 80 mots ») ou une fraction du texte (le quart, le tiers).
\item Compter les mots du texte original si une fraction est demandée : un texte de 400 mots résumé au quart donne un résumé d'environ 100 mots.
\item Rédiger une première version, puis \cle{compter les mots} de sa production (chaque mot séparé par une espace compte pour un ; les mots composés à trait d'union comptent pour un).
\item Si le texte est trop long : supprimer les adjectifs non essentiels, fusionner deux phrases, remplacer une proposition par un nom (« parce qu'il était pauvre » $\rightarrow$ « par pauvreté »).
\item Si le texte est trop court : une idée essentielle a probablement été oubliée ; relire le plan des idées.
\item Indiquer le \cle{nombre de mots} à la fin du résumé, entre parenthèses, comme l'exige souvent l'épreuve.
\end{enumerate}
\textbf{Tolérance :} on admet généralement une marge d'environ 10 \% autour du nombre imposé ; viser juste reste la règle.
\end{methode}

\begin{exoresolu}[Contraction à longueur imposée]
\textbf{Énoncé.} Résumez en environ 50 mots (plus ou moins 5 mots) le texte suivant (120 mots) :

« Dans beaucoup de villages africains, l'arbre à palabres était le centre de la vie sociale. C'est sous son ombre généreuse que les anciens se réunissaient pour régler les conflits entre familles, juger les disputes de terres, marier les jeunes gens et transmettre aux enfants l'histoire du lignage. Personne ne parlait sans y être invité par le chef, et la parole de l'aîné était écoutée dans un silence respectueux. Aujourd'hui, les palabres ont presque disparu : les tribunaux modernes, les mairies et les écoles ont remplacé l'arbre, et les jeunes préfèrent la ville au village. Pourtant, certains regrettent cette époque où la parole réglait tout. »

\tcblower
\textbf{Solution détaillée.}

\textbf{Étape 1 — Calcul de la longueur.} 120 mots $\times$ environ 40 \% $\approx$ 50 mots : la consigne est cohérente.

\textbf{Étape 2 — Idées essentielles (plan).}
\begin{enumerate}
\item L'arbre à palabres était le centre de la vie sociale du village.
\item On y réglait les conflits et on y transmettait les traditions, dans le respect de la parole des anciens.
\item Les institutions modernes l'ont fait disparaître.
\item Certains le regrettent.
\end{enumerate}

\textbf{Étape 3 — Éléments éliminés.} Les exemples de conflits (disputes de terres, mariages), les détails du protocole (silence respectueux), la mention des mairies et écoles.

\textbf{Étape 4 — Rédaction et comptage.}

\emph{Version proposée :} « L'arbre à palabres, centre de la vie sociale du village africain, servait à régler les conflits et à transmettre les traditions, dans le respect de la parole des anciens. Les institutions modernes l'ont fait disparaître, ce que certains regrettent. »

\textbf{Comptage :} L'(1) arbre(2) à(3) palabres(4), centre(5) de(6) la(7) vie(8) sociale(9) du(10) village(11) africain(12), servait(13) à(14) régler(15) les(16) conflits(17) et(18) à(19) transmettre(20) les(21) traditions(22), dans(23) le(24) respect(25) de(26) la(27) parole(28) des(29) anciens(30). Les(31) institutions(32) modernes(33) l'(34)ont(35) fait(36) disparaître(37), ce(38) que(39) certains(40) regrettent(41). Total : \textbf{41 mots}.

\textbf{Étape 5 — Ajustement.} 41 mots est légèrement sous la fourchette (45--55). On complète avec une idée du plan : « ...dans le respect de la parole des anciens, qui décidait de tout. » (4 mots de plus) et « Les institutions modernes de la ville l'ont fait disparaître... » (3 mots de plus). Nouvelle version à 48 mots, conforme à la consigne.

\textbf{Conclusion.} Le résumé final (48 mots) respecte la consigne de longueur, conserve les quatre idées essentielles et n'ajoute aucun commentaire personnel. On écrit en fin de copie : (48 mots).
\end{exoresolu}

\courssub{Méthode 4 : Mobiliser les outils de langue — énoncé constatif et performatif, lexique spécialisé}

\begin{methode}[Distinguer énoncé constatif et énoncé performatif]
\begin{enumerate}
\item Un énoncé \cle{constatif} décrit une réalité : il peut être vrai ou faux. Exemple : « Le chef a décoré Meka. » (on peut vérifier).
\item Un énoncé \cle{performatif} accomplit l'action qu'il énonce : le dire, c'est le faire. Exemples : « Je te décore », « Je promets », « La séance est ouverte ».
\item Indices du performatif : verbe à la \textbf{première personne du présent}, souvent accompagné d'une formule solennelle ; il ne peut être ni vrai ni faux, seulement valide ou non.
\item Pour le résumé : rapporter les performatifs par des verbes d'action (« l'officier décore Meka », « le chef promet ») et les constatifs par des verbes de description.
\item Repérer le \cle{lexique spécialisé} d'un texte (vocabulaire d'un domaine : colonial, médical, juridique) : dans un résumé, on le conserve s'il est indispensable au sens, sinon on le généralise.
\end{enumerate}
\end{methode}

\begin{exoresolu}[Constatif ou performatif ?]
\textbf{Énoncé.} Classez les énoncés suivants en constatifs ou performatifs, en justifiant :
\begin{enumerate}
\item « Je vous décore de la médaille de l'amitié franco-africaine. » (l'officier à Meka)
\item « Meka portait sa médaille sur sa poitrine. »
\item « Je déclare la réunion ouverte. »
\item « Le vieux nègre avait donné ses terres à la mission. »
\end{enumerate}

\tcblower
\textbf{Solution détaillée.}

\textbf{Critère appliqué :} un énoncé est performatif si le fait de le prononcer réalise l'action ; il est constatif s'il décrit une situation vérifiable (vraie ou fausse).

\textbf{Énoncé 1 : performatif.} Le verbe « décorer » est à la première personne du présent ; en prononçant cette phrase dans les conditions officielles, l'officier \emph{accomplit} la décoration. Dire « je vous décore », c'est décorer. L'énoncé n'est ni vrai ni faux : il est valide.

\textbf{Énoncé 2 : constatif.} La phrase décrit une réalité observable : on peut vérifier si Meka porte ou non sa médaille. Elle est vraie ou fausse.

\textbf{Énoncé 3 : performatif.} « Je déclare » à la première personne du présent : prononcer cette phrase ouvre effectivement la réunion. L'action est réalisée par la parole.

\textbf{Énoncé 4 : constatif.} Il s'agit d'un fait passé, rapporté au plus-que-parfait : la donation des terres est vérifiable dans l'histoire du personnage.

\textbf{Conclusion.} Les énoncés 1 et 3 sont performatifs (la parole accomplit l'acte) ; les énoncés 2 et 4 sont constatifs (la parole décrit un fait). Dans un résumé, on rapportera : « L'officier décore Meka » (action) et « Meka avait cédé ses terres » (fait).
\end{exoresolu}

\courssub{Méthode 5 : Identifier les tonalités satirique et polémique pour restituer le ton}

\begin{methode}[Reconnaître et restituer une tonalité]
\begin{enumerate}
\item La tonalité \cle{satirique} se moque d'un personnage, d'une institution ou d'une société pour dénoncer ses défauts ; procédés : ironie, exagération, ridiculisation, contrastes.
\item La tonalité \cle{polémique} attaque violemment et directement une idée ou un adversaire ; procédés : accusations, questions rhétoriques, vocabulaire péjoratif, ton véhément.
\item Indices à relever dans le texte : adjectifs ironiques, écart entre ce qui est dit et ce qui est pensé, accumulation de critiques.
\item Dans le résumé, on \textbf{ne reproduit pas} la moquerie, mais on la \textbf{signale objectivement} : « l'auteur tourne en dérision... », « le narrateur dénonce avec ironie... », « le texte critique vivement... ».
\item Ne jamais confondre la position de l'auteur avec la sienne : le résumé rapporte, il ne juge pas.
\end{enumerate}
\end{methode}

\begin{exoresolu}[Analyse de la tonalité d'un passage]
\textbf{Énoncé.} Relevez la tonalité dominante de ce passage inspiré de l'univers d'Oyono, puis formulez la phrase de résumé qui la restitue :

« Quelle gloire ! Meka, décoré ! Le Blanc lui avait tapé sur l'épaule comme on flatte un bon chien, et le voilà qui rayonnait, fier de sa médaille de fer, tandis que ses terres, ses belles terres, appartenaient désormais à la mission. »

\tcblower
\textbf{Solution détaillée.}

\textbf{Étape 1 — Relevé des indices.}
\begin{itemize}
\item Exclamation ironique : « Quelle gloire ! » — l'enthousiasme affiché est en réalité moqué.
\item Comparaison dévalorisante : « comme on flatte un bon chien » — le geste du Blanc est rabaissant.
\item Ironie sur l'objet : « médaille de fer » — la récompense est sans valeur.
\item Contraste dénonciateur : la fierté de Meka est mise en regard de la perte réelle (« ses belles terres »).
\end{itemize}

\textbf{Étape 2 — Identification.} Le passage se moque de la naïveté de Meka et dénonce l'hypocrisie de la récompense coloniale : la tonalité est \textbf{satirique} (moquerie qui dénonce), avec une pointe polémique contre l'administration coloniale.

\textbf{Étape 3 — Phrase de résumé restituant le ton.}

\emph{Proposition :} « Le narrateur tourne en dérision la fierté de Meka, décoré d'une médaille sans valeur, et dénonce ainsi l'hypocrisie d'une récompense qui masque la spoliation de ses terres. »

\textbf{Vérification.} La phrase rapporte la critique sans la reproduire directement, sans recopier les images du texte et sans ajouter d'opinion personnelle.

\textbf{Conclusion.} La tonalité satirique est correctement identifiée et restituée de manière neutre, conformément aux exigences du résumé.
\end{exoresolu}

\courssub{Méthode 6 : Améliorer et justifier sa production (relecture et auto-correction)}

\begin{methode}[Améliorer un résumé : la grille de relecture]
\begin{enumerate}
\item \textbf{Fidélité} : relire le texte original et vérifier qu'aucune idée essentielle ne manque et qu'aucune idée étrangère n'a été ajoutée.
\item \textbf{Reformulation} : surligner les groupes de mots repris tels quels au texte ; s'il y en a plus de trois ou quatre consécutifs, reformuler.
\item \textbf{Cohérence} : vérifier les connecteurs logiques et la progression (le résumé doit se lire comme un texte autonome, pas comme une liste).
\item \textbf{Langue} : corriger l'orthographe, les accords, la concordance des temps (le résumé s'écrit en général au présent de vérité générale ou aux temps du récit).
\item \textbf{Longueur} : recompter les mots et ajuster ; inscrire le total.
\item \textbf{Présentation} : écrire « Résumé » en titre si demandé, sauter une ligne, soigner l'écriture.
\end{enumerate}
\textbf{Justifier sa démarche} (savoir-faire exigé) : savoir expliquer oralement ou par écrit pourquoi telle phrase a été supprimée (exemple, répétition) et telle autre conservée (idée directrice).
\end{methode}

\begin{exoresolu}[Correction d'un résumé fautif]
\textbf{Énoncé.} Voici le résumé produit par un élève à partir du texte sur l'arbre à palabres (consigne : environ 50 mots). Relevez quatre défauts et proposez une version améliorée.

« Je trouve que ce texte est très intéressant. L'arbre à palabres était le centre de la vie sociale et c'est sous son ombre généreuse que les anciens se réunissaient pour régler les conflits entre familles et juger les disputes de terres. Les tribunaux modernes ont remplacé l'arbre. C'est dommage car les jeunes préfèrent la ville. » (59 mots)

\tcblower
\textbf{Solution détaillée.}

\textbf{Étape 1 — Relevé des défauts.}
\begin{enumerate}
\item \textbf{Opinion personnelle} : « Je trouve que ce texte est très intéressant » et « C'est dommage » — le résumé doit être neutre ; l'élève s'exprime à la première personne, ce qui est interdit.
\item \textbf{Recopie du texte} : « le centre de la vie sociale », « sous son ombre généreuse », « régler les conflits entre familles » sont repris mot pour mot.
\item \textbf{Exemples conservés} : « les disputes de terres » est un détail secondaire qui aurait dû être généralisé.
\item \textbf{Idée essentielle déformée} : la dernière phrase mélange la disparition des palabres et le regret des anciens ; la transmission des traditions (idée du texte) a disparu.
\end{enumerate}

\textbf{Étape 2 — Justification des suppressions.} Les exemples (disputes de terres) illustrent l'idée « régler les conflits » : ils sont supprimables sans perte de sens. Les jugements personnels n'appartiennent pas au texte : ils violent la règle de fidélité.

\textbf{Étape 3 — Version améliorée.}

\emph{Proposition courte :} « Centre de la vie sociale du village, l'arbre à palabres permettait aux anciens de régler les conflits et de transmettre les traditions. Les institutions modernes l'ont fait disparaître, au regret de certains. » (31 mots — version acceptable seulement si la consigne fixe une fourchette large ; ici elle est trop courte.)

\emph{Version conforme à la consigne :} « Centre de la vie sociale du village africain, l'arbre à palabres permettait aux anciens de régler les conflits entre familles et de transmettre l'histoire et les traditions aux jeunes générations. Aujourd'hui, les institutions modernes l'ont fait disparaître, ce que certains regrettent. » (44 mots, dans la tolérance de 45--55 si l'on admet la marge de 10 \% ; on peut ajouter « de la ville » après « institutions modernes » pour atteindre 47 mots.)

\textbf{Conclusion.} La version améliorée est neutre, reformulée, fidèle aux quatre idées essentielles et conforme à la longueur demandée : elle satisferait aux critères de notation de l'épreuve.
\end{exoresolu}

\begin{attention}[Les 5 erreurs qui coûtent des points au Bac]
\begin{enumerate}
\item \textbf{Recopier des phrases du texte.} Toute reprise de plus de trois ou quatre mots consécutifs est sanctionnée : le résumé exige une reformulation complète (synonymes, changement de constructions).
\item \textbf{Donner son avis.} « Je pense que », « à mon avis », « ce texte est beau » : le résumé rapporte la pensée de l'auteur, jamais celle du candidat. La première personne est interdite.
\item \textbf{Conserver les exemples et les détails.} Anecdotes, chiffres, comparaisons, citations : ce sont des illustrations à éliminer. Ne garder que les idées directrices.
\item \textbf{Négliger la consigne de longueur.} Ne pas compter ses mots, ignorer la fraction imposée ou oublier d'inscrire le total entre parenthèses fait perdre des points faciles. Compter deux fois : en rédigeant, puis en relisant.
\item \textbf{Produire une liste d'idées sans liens.} Un résumé est un texte continu et cohérent : sans connecteurs logiques (d'abord, cependant, c'est pourquoi, enfin), la copie ressemble à un plan et pénalise la note d'expression. Ajouter à cela les fautes d'accord et de conjugaison non relues, qui alourdissent encore la sanction.
\end{enumerate}
\end{attention}

\newpage

\coursec{Exercices}

\coursec{Leçon 16 : Rédiger un résumé de texte}

\courssub{Le résumé : définition, enjeux et cadre de l'épreuve}

\begin{aretenir}[Définition du résumé]
Le \cle{résumé} (ou \cle{contraction de texte}) est un exercice de réécriture qui consiste à restituer les \textbf{idées essentielles} d'un texte source dans un nombre de mots imposé (généralement le \textbf{quart} du texte original, soit 25\,\%), \textbf{sans en modifier le sens}, \textbf{sans y ajouter d'idées personnelles} et \textbf{sans recopier} des phrases entières. Il exige une \cle{reformulation} (changement de vocabulaire, de structure syntaxique, passage au style indirect) et une \cle{hiérarchisation} de l'information.
\end{aretenir}

\begin{aretenir}[Cadre de l'épreuve au Baccalauréat technique]
\begin{itemize}
\item \textbf{Longueur du texte support :} environ 400 mots.
\item \textbf{Longueur attendue :} 100 mots (\textbf{au quart}).
\item \textbf{Marge de tolérance :} $\pm 10\,\%$ (soit entre 90 et 110 mots pour un quart de 400 mots).
\item \textbf{Consignes types :} « Résumez ce texte en 100 mots $\pm 10\,\%$ », « Indiquez le nombre de mots employés », « Soulignez deux connecteurs logiques et expliquez leur rôle ».
\item \textbf{Compétences évaluées :} compréhension globale, sélection de l'information, capacité de synthèse, maîtrise de la reformulation, correction linguistique (syntaxe, vocabulaire, connecteurs), respect des contraintes formelles.
\end{itemize}
\end{aretenir}

\begin{methode}[Étapes de la contraction de texte]
\begin{enumerate}
\item \textbf{Lecture active :} Lire 2 à 3 fois pour saisir le sujet, la thèse, le plan, le ton.
\item \textbf{Découpage / Macrostructure :} Diviser le texte en séquences logiques (introduction, développement, conclusion) ou en étapes narratives. Donner un titre à chaque séquence.
\item \textbf{Sélection des idées essentielles :} Souligner les mots-clés, les arguments principaux, les tournants narratifs. \textbf{Barre}r les illustrations, exemples, détails, répétitions, digressions.
\item \textbf{Prise de notes / Schéma :} Noter les idées retenues sous forme de mots-clés ou phrases très courtes, dans l'ordre logique. Utiliser un arbre ou un tableau.
\item \textbf{Rédaction du brouillon :} Enchaîner les idées par des \cle{connecteurs logiques} (donc, mais, ensuite, cependant, en effet...). Passer au \textbf{style indirect} (ou présent de narration), 3\textsuperscript{e} personne. Reformuler le vocabulaire (synonymes, généralisation).
\item \textbf{Comptage et ajustement :} Compter les mots. Si $> 110$ : couper le secondaire, condenser les tournures. Si $< 90$ : développer une idée clé, ajouter un connecteur manquant, préciser un agent.
\item \textbf{Copie finale et relecture :} Vérifier orthographe, grammaire, cohérence, nombre de mots exact noté à la fin.
\end{enumerate}
\end{methode}

\begin{attention}[Pièges à éviter absolument]
\begin{itemize}
\item \textbf{Recopie (plagiat) :} Copier-coller des phrases du texte $\rightarrow$ note éliminatoire.
\item \textbf{Avis personnel / Jugement :} « Je pense que... », « C'est injuste... », « L'auteur a raison... » $\rightarrow$ hors sujet.
\item \textbf{Hors quota :} $< 90$ ou $> 110$ mots $\rightarrow$ pénalité lourde (souvent note sur 5 ou 0 sur la forme).
\item \textbf{Style direct / Dialogue :} Garder les guillemets ou le « je » du personnage $\rightarrow$ non reformulé.
\item \textbf{Oubli d'idées majeures :} Zapper une étape narrative ou un argument central $\rightarrow$ perte de points sur le fond.
\item \textbf{Connecteurs absents ou erronés :} Suite de phrases juxtaposées sans logique $\rightarrow$ incohérence.
\end{itemize}
\end{attention}

\courssub{Lecture analytique : repérer les idées essentielles (à partir de \emph{Le Vieux nègre et la médaille})}

\begin{savaistu}[Contexte de l'œuvre : \emph{Le Vieux nègre et la médaille}]
\begin{itemize}
\item \textbf{Auteur :} Ferdinand Oyono (Cameroun, 1929--2010), diplomate et écrivain majeur de la littérature africaine d'expression française.
\item \textbf{Publication :} 1956 (éd. René Julliard). Période coloniale tardive, veille des indépendances (Cameroun : 1960).
\item \textbf{Genre :} Roman court / Nouvelle longue, \cle{roman satirique} et \cle{polémique}.
\item \textbf{Intrigue :} Meka, vieux chef villageois « évolué », loyal à l'administration coloniale, attend la médaille d'honneur. La cérémonie tourne à la farce : pluie, fuite des Blancs, maladie de Meka, moqueries des siens. Désillusion totale.
\item \textbf{Thèmes :} Aliénation coloniale, illusion de l'assimilation, trahison des pères, ridicule du pouvoir colonial, lucidité tardive.
\end{itemize}
\end{savaistu}

\begin{methode}[Repérer les idées essentielles dans un texte narratif]
\begin{enumerate}
\item Identifier le \textbf{schéma actanciel} : Qui fait quoi ? (Sujet / Objet / Destinateur / Destinataire / Adjuvant / Opposant).
\item Repérer les \textbf{nœuds narratifs} : Situation initiale $\rightarrow$ Événement déclencheur $\rightarrow$ Péripéties $\rightarrow$ Climax $\rightarrow$ Dénouement $\rightarrow$ Situation finale.
\item Distinguer \textbf{Intrigue principale} (Meka / Médaille / Blancs) et \textbf{Intrigue secondaire} (Villageois / Fils / Terres).
\item Ne retenir que les actions qui font \textbf{avancer l'intrigue principale} et les \textbf{révélations} sur le sens de l'œuvre (ironie, critique).
\end{enumerate}
\end{methode}

\courssub{Outils de langue 1 : Énoncé constatif et performatif, Lexique spécialisé}

\begin{definition}[Énoncé constatif]
Un \cle{énoncé constatif} (ou assertif) \textbf{describe un état de choses}, une réalité extérieure au langage. Il porte sur le \textbf{vrai} ou le \textbf{faux} (valeur de vérité). Il \textbf{informe}.
\begin{itemize}
\item \textit{Exemple :} « Le commandant de cercle arriva au village. » (Vrai ou faux selon la réalité).
\item \textit{Marqueurs :} Verbes d'état, de perception, de déclaration neutre au passé, présent, futur. Sujet $\neq$ locuteur (souvent 3\textsuperscript{e} pers.).
\end{itemize}
\end{definition}

\begin{definition}[Énoncé performatif]
Un \cle{énoncé performatif} (J.L. Austin) \textbf{accomplit l'action qu'il énonce} par le seul fait d'être prononcé dans des conditions appropriées (autorité, contexte, formule rituelle). Il ne se juge pas sur le vrai/faux mais sur le \textbf{réussi/raté} (heurteux / infélicieux).
\begin{itemize}
\item \textit{Exemple :} « Je vous déclare mari et femme. » (Le maire \textbf{marie} en le disant).
\item \textit{Test :} On peut souvent préfixer « \textbf{Par la présente, je...} » ou remplacer par « Je \textbf{prononce}... ».
\item \textit{Classes (Searle) :} \textbf{Déclaratifs} (changent le monde : nommer, baptiser, déclarer la guerre), \textbf{Engageants} (promettre, jurer), \textbf{Directifs} (ordonner, supplier), \textbf{Expressifs} (remercier, féliciter, s'excuser), \textbf{Assertifs} (affirmer, constater -- limite constatif/performatif).
\end{itemize}
\end{definition}

\begin{propriete}[Critères de distinction simplifiés pour le Bac]
\begin{center}
\begin{tabularx}{\linewidth}{|l|X|X|}\hline
\textbf{Critère} & \textbf{Constatif} & \textbf{Performatif} \\ \hline
\textbf{Fonction} & Informer, décrire & Agir, transformer la réalité \\ \hline
\textbf{Vérité} & Vrai / Faux & Heureux / Infélicieux (réussi / raté) \\ \hline
\textbf{Sujet / Verbe} & Souvent 3\textsuperscript{e} pers. / Verbes d'état/action & \textbf{1\textsuperscript{re} pers. sing. ou plur.} / \textbf{Verbe performatif} (promets, déclare, ordonne, baptise, décerne...) \\ \hline
\textbf{Temps} & Tous (souvent passé narratif) & \textbf{Présent indicatif} (souvent) \\ \hline
\textbf{Contexte} & N'importe quel énonciateur & \textbf{Autorité / Qualité} requise (juge, maire, prêtre, commandant...) \\ \hline
\end{tabularx}
\end{center}
\end{propriete}

\begin{aretenir}[Lexique spécialisé (ou vocabulaire technique)]
Le \cle{lexique spécialisé} est l'ensemble des termes propres à un \textbf{domaine professionnel, scientifique, technique, administratif} ou \textbf{institutionnel}. Il assure la \textbf{précision}, l'\textbf{univocité} (un sens unique) et l'\textbf{efficacité} communicationnelle entre experts.
\begin{itemize}
\item \textit{Exemples (Administration coloniale) :} \textbf{Commandant de cercle, Administrateur, Résident, Interprète, Canton, Notables, Citation, Décoration, Médaille, Garde, Cercle, Subdivision.}
\item \textit{Exemples (Technique/Industriel) :} \textbf{Usinage, Tolérance, Nomenclature, Cahier des charges, ISO, Maintenance préventive, Tournage, Fraisage.}
\end{itemize}
\textbf{Au Bac :} Savoir identifier ces termes, les définir (par synonymie, hyperonymie, contexte) et expliquer leur effet de réel / de crédibilité / de pouvoir dans le texte.
\end{aretenir}

\courssub{Outils de langue 2 : Tonalités satirique et polémique}

\begin{definition}[Tonalité satirique]
La \cle{tonalité satirique} vise à \textbf{corriger les mœurs} par le \textbf{rire} et la \cle{moquerie}. Elle dénonce les vices, la bêtise, l'hypocrisie, le ridicule en les \textbf{grossissant} pour les rendre visibles. Elle suppose une \textbf{distance ironique} de l'auteur vis-à-vis de son sujet.
\begin{itemize}
\item \textbf{Procédés clés :} \cle{Ironie} (dire le contraire de ce qu'on pense), \cle{Caricature} (exagération physique/morale), \cle{Parodie} (imitation déformée), \cle{Antiphrase}, \cle{Understatement} (litote), \cle{Zeugme}, \cle{Comparaison dévalorisante}.
\item \textit{Effet :} Rire critique, complicité lecteur/auteur, prise de recul.
\end{itemize}
\end{definition}

\begin{definition}[Tonalité polémique]
La \cle{tonalité polémique} (du grec \textit{polemos}, guerre) vise à \textbf{combattre un adversaire}, une idée, un système avec \textbf{violence verbale} et \textbf{indignation}. Elle ne rit pas, elle \textbf{accuse}, \textbf{condamne}, \textbf{révolte}. Le ton est \textbf{grave, pathétique, agressif, direct}.
\begin{itemize}
\item \textbf{Procédés clés :} \cle{Dénonciation directe} (verbes forts : voler, assassiner, mentir, exploiter), \cle{Accumulation} d'horreurs, \cle{Anaphore} martelée, \cle{Hyperbole} dramatique, \cle{Vocabulaire péjoratif/chargé}, \cle{Apostrophe} à l'ennemi, \cle{Questions rhétoriques} accusatrices.
\item \textit{Effet :} Choc, indignation, adhésion émotionnelle, appel à l'action/révolte.
\end{itemize}
\end{definition}

\begin{aretenir}[Différences clés : Satirique vs Polémique]
\begin{center}
\begin{tabularx}{\linewidth}{|l|X|X|}\hline
\textbf{Aspect} & \textbf{Satirique} & \textbf{Polémique} \\ \hline
\textbf{Objectif} & Corriger par le rire / Montrer le ridicule & Combattre / Détruire l'adversaire / Révolter \\ \hline
\textbf{Émotion} & Amusement, ironie, distance & Colère, indignation, empathie victime, urgence \\ \hline
\textbf{Arme} & L'intelligence, l'esprit, le second degré & La force, la vérité crue, le pathos \\ \hline
\textbf{Procédés majeurs} & Ironie, caricature, parodie, litote & Dénonciation directe, anaphore, hyperbole, apostrophe \\ \hline
\textbf{Exemple \emph{Vieux nègre}} & Description du commandant « ventre tendu », cérémonie grotesque & Cri de Meka / Villageois : « Ils ont pris nos terres... » \\ \hline
\end{tabularx}
\end{center}
\textbf{Note :} Une œuvre (comme \emph{Le Vieux nègre}) peut \textbf{alterner} les deux : satire sur le pouvoir colonial (ridicule), polémique sur les conséquences (douleur, vol, mort).
\end{aretenir}

\courssub{Rédaction du résumé : reformulation et production}

\begin{methode}[Techniques de reformulation (outils de condensation)]
\begin{enumerate}
\item \textbf{Substitution lexicale :} Synonymes, termes plus généraux (hyponymes $\rightarrow$ hyperonymes). \textit{Ex:} « donna ses terres aux missionnaires et ses fils à la guerre » $\rightarrow$ « sacrifia ses biens et sa descendance ».
\item \textbf{Nominalisation / Verbalisation :} Transformer verbe $\rightarrow$ nom ou nom $\rightarrow$ verbe pour gagner en concision. \textit{Ex:} « Il attendait avec impatience » $\rightarrow$ « Son attente impatiente » / « Il attendit fébrilement ».
\item \textbf{Suppression du superflu :} Adjectifs qualificatifs non essentiels, compléments de détail, relativelles explicatives. \textit{Ex:} « Son boubou le plus propre, celui des grandes occasions » $\rightarrow$ « Son boubou de cérémonie ».
\item \textbf{Passage au style indirect / Infinitif / Participe :} « Le commandant arriva et il lut la citation » $\rightarrow$ « À l'arrivée du commandant, la lecture de la citation... » / « Le commandant arriva, lut la citation... »
\item \textbf{Fusion de phrases / Subordination :} Relier par des connecteurs logiques (cause, conséquence, concession, chronologie). \textit{Ex:} « Meka était fier. Il montrait sa médaille. » $\rightarrow$ « Fier, Meka exhibait sa médaille. »
\item \textbf{Généralisation / Résumé de séquences :} Remplacer une scène détaillée par un nom d'action. \textit{Ex:} Tout le paragraphe de la cérémonie $\rightarrow$ « La cérémonie officielle se déroule. »
\end{enumerate}
\end{methode}

\begin{exemplebox}[Exemple de reformulation guidée]
\textbf{Texte source (32 mots) :} « Le vieux Meka, qui avait donné ses deux fils à la France pendant la guerre et offert ses terres aux missionnaires, attendait avec une fierté immense la médaille promise. »
\textbf{Idées clés :} Meka -- don fils/guerre -- don terres/missionnaires -- attente fière -- médaille promise.
\textbf{Résumé (11 mots) :} \textit{Meka, ayant sacrifié fils et terres, attend fièrement la médaille promise.}
\textbf{Gains :} Relative $\rightarrow$ Participe passé ; « avec une fierté immense » $\rightarrow$ « fièrement » ; fusion causes/conséquence.
\end{exemplebox}

\courssub{Amélioration du résumé et bilan de l'œuvre}

\begin{methode}[Auto-correction du résumé (Check-list finale)]
\begin{enumerate}
\item \textbf{Compter les mots} (mots pleins + mots outils). Noter le total.
\item \textbf{Vérifier le quota :} $90 \le N \le 110$ ? Sinon $\rightarrow$ Coupe / Développe.
\item \textbf{Vérifier le fond :} Les 3--4 idées majeures sont-elles là ? Ordre logique respecté ? Pas d'invention ? Pas d'oubli crucial ?
\item \textbf{Vérifier la forme :} 3\textsuperscript{e} personne ? Présent (ou passé narratif cohérent) ? Style indirect ? Pas de guillemets ? Pas de « je » / « à mon avis » ?
\item \textbf{Vérifier la cohésion :} Connecteurs logiques variés et justes ? Références anaphoriques claires (il, ce dernier, celui-ci) ?
\item \textbf{Vérifier la langue :} Accords, orthographe, ponctuation, vocabulaire précis (lexique spécialisé si pertinent).
\item \textbf{Souligner 2 connecteurs} et annoter leur rôle (ex: \textit{Cependant} = opposition ; \textit{Ensuite} = chronologie).
\end{enumerate}
\end{methode}

\begin{savaistu}[Bilan de l'œuvre : \emph{Le Vieux nègre et la médaille}]
\begin{itemize}
\item \textbf{Portée historique :} Publié 1956, l'œuvre anticipe la décolonisation. Elle démonte le mythe de la « mission civilisatrice » et de l'assimilation « à la française ».
\item \textbf{Portée littéraire :} Premier roman camerounais d'envergure internationale. Style sobre, oralité maîtrisée, structure circulaire (attente $\rightarrow$ cérémonie $\rightarrow$ chute $\rightarrow$ lucidité).
\item \textbf{Tonalité dominante :} \textbf{Satirique} (portrait du commandant, cérémonie grotesque, naïveté de Meka) teintée de \textbf{polémique} (cri final, vol des terres/vies, injustice structurelle).
\item \textbf{Message universel :} La dignité ne s'achète pas avec des médailles. La lucidité vient trop tard pour les « collaborateurs » du système oppresseur. Critique de l'aliénation mentale.
\end{itemize}
\end{savaistu}

\courssub{Je vérifie mes connaissances}

\begin{exercice}[QCM : les règles du résumé]
Pour chaque question, choisis la bonne réponse.
\begin{enumerate}
\item Le résumé au quart d'un texte de 400 mots doit comporter environ :
\begin{enumerate}
\item 200 mots ; \item 100 mots ; \item 50 mots ; \item 25 mots.
\end{enumerate}
\item Dans un résumé, il est interdit de :
\begin{enumerate}
\item reformuler les idées de l'auteur ; \item donner son avis personnel ; \item utiliser des connecteurs logiques ; \item écrire au présent.
\end{enumerate}
\item Un énoncé performatif est un énoncé qui :
\begin{enumerate}
\item décrit une réalité ; \item peut être vrai ou faux ; \item accomplit une action par le seul fait d'être prononcé ; \item exprime un sentiment.
\end{enumerate}
\item La tonalité qui consiste à dénoncer avec violence et indignation est la tonalité :
\begin{enumerate}
\item satirique ; \item polémique ; \item lyrique ; \item tragique.
\end{enumerate}
\end{enumerate}
\end{exercice}

\begin{exercice}[Vrai ou faux justifié]
Dis si chaque affirmation est vraie ou fausse, puis justifie ta réponse en une phrase.
\begin{enumerate}
\item On peut recopier des phrases entières du texte dans un résumé.
\item Le résumé doit conserver toutes les idées essentielles du texte.
\item « Je vous déclare mari et femme » est un énoncé constatif.
\item Le lexique spécialisé est un vocabulaire propre à un domaine (administration, médecine, colonisation\ldots).
\item L'ironie est un procédé typique de la tonalité satirique.
\end{enumerate}
\end{exercice}

\begin{exercice}[Définitions]
Définis en deux ou trois phrases chacun des termes suivants, en donnant un exemple pour chacun :
\begin{enumerate}
\item la contraction de texte ;
\item l'énoncé constatif ;
\item la reformulation ;
\item la tonalité satirique.
\end{enumerate}
\end{exercice}

\begin{exercice}[Calculs directs : le respect des quotas de mots]
\begin{enumerate}
\item Un texte compte 600 mots. Combien de mots doit comporter son résumé au quart ? Donne la marge tolérée de $\pm 10\,\%$.
\item Un candidat a rédigé un résumé de 160 mots pour un texte de 400 mots. Respecte-t-il la consigne ? Justifie par un calcul.
\item Un résumé doit comporter 100 mots $\pm 10\,\%$. Entre quelles valeurs extrêmes le nombre de mots doit-il se situer ?
\end{enumerate}
\end{exercice}

\courssub{Je m'entraîne}

\begin{exercice}[Constatifs ou performatifs ?]
Classe les huit énoncés suivants, inspirés de l'univers du \emph{Vieux nègre et la médaille} de Ferdinand Oyono, en énoncés constatifs ou performatifs. Justifie chaque classement.
\begin{enumerate}
\item « Meka attendait sa médaille depuis de longues années. »
\item « Je vous décore au nom de la France. »
\item « Je promets de venir à la cérémonie. »
\item « Le commandant de cercle arriva au village. »
\item « Je vous remercie de votre fidélité. »
\item « La pluie tomba pendant toute la nuit. »
\item « Je déclare la séance ouverte. »
\item « Les villageois dansèrent jusqu'au matin. »
\end{enumerate}
\end{exercice}

\begin{exercice}[Lexique spécialisé colonial]
Dans le passage suivant, relève cinq termes appartenant au lexique spécialisé de l'administration coloniale, puis explique le sens de chacun dans le contexte.
« Le commandant de cercle, accompagné de l'interprète et des gardes, présida la cérémonie de remise de la médaille. L'administrateur lut la citation, épingla la décoration sur la poitrine du vieux Meka, puis le résident prononça un discours devant les notables du canton rassemblés. »
\end{exercice}

\begin{exercice}[Reformulation guidée]
Condense chacune des phrases suivantes (environ 30 mots) en une phrase d'environ 10 mots, sans perte du sens essentiel.
\begin{enumerate}
\item « Le vieux Meka, qui avait donné ses deux fils à la France pendant la guerre et offert ses terres aux missionnaires, attendait avec une fierté immense la médaille promise. »
\item « Pendant la nuit qui suivit la cérémonie, une pluie violente et incessante s'abattit sur le village, dispersant les invités et transformant la fête en désastre. »
\item « Les Blancs, qui avaient comblé Meka de promesses et de belles paroles, l'abandonnèrent sans hésiter dès que la cérémonie fut terminée. »
\item « Meka, déçu et humilié, comprit enfin que la médaille n'était qu'un morceau de métal sans valeur qui ne changerait rien à sa condition. »
\item « Les anciens du village, réunis autour du feu, se moquaient ouvertement de la naïveté du vieil homme qui avait cru à l'amitié des Blancs. »
\end{enumerate}
\end{exercice}

\begin{exercice}[Repérage des idées essentielles]
Lis le texte de presse suivant (environ 120 mots) : « Depuis l'ouverture de l'année scolaire, les lycées techniques de la région du Centre font face à un manque criard de matériel pédagogique. Les ateliers de mécanique, pourtant essentiels à la formation des élèves de Terminale, fonctionnent avec des machines vieilles de vingt ans. Certains établissements ont reçu des promesses d'équipement, mais rien de concret n'est arrivé. Les enseignants dénoncent une situation qui compromet la préparation du Baccalauréat technique. Les parents d'élèves, réunis samedi dernier à Yaoundé, ont demandé au ministère des mesures urgentes. Une délégation doit être reçue la semaine prochaine. »
\begin{enumerate}
\item Souligne les cinq idées essentielles du texte.
\item Barre les éléments secondaires (exemples, précisions, détails).
\item Justifie chacun de tes choix en une phrase.
\end{enumerate}
\end{exercice}

\courssub{Je me prépare au Bac}

\begin{exercice}[Résumé au quart — type Baccalauréat technique]
Lis attentivement l'extrait suivant, adapté du \emph{Vieux nègre et la médaille} de Ferdinand Oyono (environ 400 mots) :

« Meka avait attendu ce jour pendant des années. La médaille ! La médaille de l'amitié française que le commandant lui avait promise ! Il avait donné ses terres aux missionnaires, il avait donné ses deux fils à la guerre de la France, ses deux fils qui n'étaient jamais revenus. Qu'avait-il reçu en échange ? Des sourires, des poignées de main, des promesses. Mais aujourd'hui, c'était différent : aujourd'hui, on allait lui épingler sur la poitrine ce morceau de métal brillant qui prouverait à tout le village que les Blancs l'aimaient.

Le matin de la cérémonie, il s'était levé avant l'aube. Il avait enfilé son boubou le plus propre, celui des grandes occasions. Ses vieilles jambes tremblaient d'émotion tandis qu'il marchait vers la place du village où l'on avait dressé une estrade. Les femmes chantaient, les tam-tams résonnaient, les enfants couraient partout. Tout le village était là, même les anciens qui se moquaient de lui d'habitude.

Le commandant arriva dans sa voiture couverte de poussière. Il était accompagné de l'interprète et de deux gardes. La cérémonie commença. L'administrateur lut une longue citation que Meka ne comprit pas, car elle était en français. Puis on lui épingla la médaille. La foule applaudit. Meka rayonnait. Il se sentait important, presque l'égal des Blancs.

Mais le soir, après la fête, la pluie se mit à tomber. Une pluie terrible, une pluie de fin du monde. Les Blancs remontèrent vite dans leurs voitures et partirent sans même se retourner. Meka resta seul sous l'averse, sa médaille collée contre sa poitrine mouillée. Le lendemain, il tomba malade. Alité, il regardait la médaille accrochée au mur de sa case et il pensait à ses fils morts, à ses terres perdues. "Qu'est-ce que cela me rapporte, cette médaille ?" se demandait-il. Les anciens du village, venus le visiter, riaient doucement de sa naïveté. »

\begin{enumerate}
\item Releve dans le texte les idées essentielles et classe-les en trois ou quatre grandes étapes du récit.
\item Rédige le résumé de cet extrait en 100 mots ($\pm 10\,\%$), en reformulant les idées sans recopier les phrases du texte et sans donner ton avis.
\item Indique à la fin de ta production le nombre exact de mots utilisés.
\item Souligne dans ton résumé deux connecteurs logiques et explique leur rôle.
\end{enumerate}
\end{exercice}

\begin{exercice}[Correction d'un mauvais résumé]
Voici le résumé produit par un candidat à partir d'un texte de 400 mots sur la colonisation :

« "Meka avait donné ses terres aux missionnaires et ses deux fils à la guerre de la France." Personnellement, je trouve que les Blancs ont été très injustes avec ce vieil homme qui méritait mieux. "Le matin de la cérémonie, il s'était levé avant l'aube et avait enfilé son boubou le plus propre." C'est vraiment scandaleux ce que les colons ont fait en Afrique ! Le commandant arriva dans sa voiture couverte de poussière, accompagné de l'interprète et de deux gardes, et la cérémonie commença pendant que les femmes chantaient et que les tam-tams résonnaient dans tout le village. Moi aussi, j'ai déjà assisté à une cérémonie de remise de prix dans mon école et je comprends ce que Meka a ressenti. En conclusion, ce texte de Ferdinand Oyono est très intéressant et je le recommande à tous les élèves de Terminale. » (environ 150 mots)

\begin{enumerate}
\item Repère et relève quatre types d'erreurs commises par ce candidat (recopie, avis personnel, dépassement du quota, éléments secondaires\ldots).
\item Pour chaque erreur, rappelle la règle du résumé qui a été violée.
\item Propose une version corrigée du résumé en 100 mots ($\pm 10\,\%$), en indiquant le nombre de mots.
\end{enumerate}
\end{exercice}

\begin{exercice}[Tonalités et bilan de l'œuvre — type question de synthèse]
\textbf{Partie A.} Pour chacun des trois courts extraits suivants, inspirés du \emph{Vieux nègre et la médaille}, identifie la tonalité dominante (satirique ou polémique) et relève les procédés utilisés (ironie, caricature, exagération, dénonciation directe\ldots) :
\begin{enumerate}
\item « Le commandant, le ventre tendu dans son uniforme trop serré, souriait de toutes ses dents en épinglant la médaille, comme s'il accomplissait le geste le plus généreux du siècle. »
\item « Ils ont pris nos terres, ils ont pris nos fils, et ils nous donnent en échange un morceau de métal qui ne vaut même pas une poule ! »
\item « Meka, si fier de sa médaille, la montrait à tous comme un enfant montre un jouet, sans voir que les Blancs riaient dans son dos. »
\end{enumerate}
\textbf{Partie B.} Rédige en une dizaine de lignes la présentation de l'œuvre \emph{Le Vieux nègre et la médaille} : auteur, date et contexte de publication, thèmes principaux, tonalité dominante et portée de l'œuvre. Ce travail pourrait servir d'introduction à un commentaire de texte au Baccalauréat technique.
\end{exercice}

\newpage

\coursec{Corrigés des exercices}

\coursec{Corrigés des exercices — Le résumé de texte (Le Vieux nègre et la médaille)}

\begin{corrige}[Exercice 1 --- QCM : les règles du résumé]
\begin{enumerate}
\item \textbf{Réponse b : 100 mots.} Le quart de 400 mots se calcule ainsi : $400 \div 4 = 100$ mots.
\item \textbf{Réponse b : donner son avis personnel.} Le résumé restitue les idées de l'auteur, jamais celles du candidat. Reformuler, utiliser des connecteurs et écrire au présent sont au contraire recommandés.
\item \textbf{Réponse c : accomplit une action par le seul fait d'être prononcé.} C'est la définition même du performatif (Austin) : dire, c'est faire (ex. « je vous déclare mari et femme »).
\item \textbf{Réponse b : polémique.} La tonalité polémique combat et dénonce avec violence et indignation ; la tonalité satirique, elle, corrige par le rire et l'ironie.
\end{enumerate}
\end{corrige}

\begin{corrige}[Exercice 2 --- Vrai ou faux justifié]
\begin{enumerate}
\item \textbf{Faux.} Recopier des phrases entières du texte constitue du plagiat et entraîne une note éliminatoire : le résumé exige la reformulation.
\item \textbf{Vrai.} Le résumé doit conserver toutes les idées essentielles (thèse, arguments principaux, étapes narratives), seuls les exemples, détails et répétitions sont supprimés.
\item \textbf{Faux.} « Je vous déclare mari et femme » est un énoncé \cle{performatif} : le maire accomplit l'acte de mariage en prononçant ces mots, dans des conditions d'autorité appropriées.
\item \textbf{Vrai.} Le lexique spécialisé est l'ensemble des termes propres à un domaine (administration, médecine, technique\ldots) ; il assure précision et univocité (ex. « commandant de cercle », « canton »).
\item \textbf{Vrai.} L'ironie (dire le contraire de ce qu'on pense), avec la caricature et la parodie, est le procédé majeur de la tonalité satirique, qui corrige les mœurs par le rire.
\end{enumerate}
\end{corrige}

\begin{corrige}[Exercice 3 --- Définitions]
\begin{enumerate}
\item \textbf{La contraction de texte} est un exercice de réécriture qui consiste à restituer les idées essentielles d'un texte source dans un nombre de mots imposé (généralement le quart), sans modifier le sens, sans ajouter d'idées personnelles et sans recopier. \textit{Exemple :} résumer en 100 mots un article de 400 mots sur la colonisation.
\item \textbf{L'énoncé constatif} est un énoncé qui décrit un état de choses, une réalité extérieure au langage ; il peut être vrai ou faux. \textit{Exemple :} « Le commandant de cercle arriva au village. »
\item \textbf{La reformulation} est la technique qui consiste à exprimer une idée avec d'autres mots et d'autres structures (synonymes, nominalisation, style indirect) sans en changer le sens. \textit{Exemple :} « il attendait avec une fierté immense » devient « il attendait fièrement ».
\item \textbf{La tonalité satirique} est le ton d'un texte qui vise à corriger les mœurs par le rire et la moquerie, en dénonçant les vices et le ridicule grâce à l'ironie, la caricature ou la parodie. \textit{Exemple :} le portrait du commandant « au ventre tendu » dans \emph{Le Vieux nègre et la médaille}.
\end{enumerate}
\end{corrige}

\begin{corrige}[Exercice 4 --- Calculs directs : le respect des quotas de mots]
\begin{enumerate}
\item Résumé au quart d'un texte de 600 mots :
\[ 600 \div 4 = 150 \text{ mots.} \]
Marge de $\pm 10\,\%$ : $10\,\%$ de 150 $= 150 \times 0,10 = 15$ mots. Donc :
\[ 150 - 15 = 135 \quad \text{et} \quad 150 + 15 = 165. \]
Le résumé doit comporter \textbf{150 mots}, avec une tolérance comprise entre \textbf{135 et 165 mots}.
\item Le candidat a écrit 160 mots pour un texte de 400 mots. Quota attendu : $400 \div 4 = 100$ mots, tolérance : de 90 à 110 mots. Or $160 > 110$. \textbf{La consigne n'est pas respectée} : le résumé dépasse de 50 mots la borne supérieure ($160 - 110 = 50$). Le candidat sera pénalisé sur le respect des contraintes.
\item Pour 100 mots $\pm 10\,\%$ : $10\,\%$ de 100 $= 10$ mots. Donc $100 - 10 = 90$ et $100 + 10 = 110$. Le nombre de mots doit se situer \textbf{entre 90 et 110 mots inclus}.
\end{enumerate}
\textbf{Points de vigilance :} toujours écrire le calcul du quart ($\div 4$), puis celui des deux bornes ; penser à conclure par une phrase (« la consigne est / n'est pas respectée »).
\end{corrige}

\begin{corrige}[Exercice 5 --- Constatifs ou performatifs ?]
\begin{enumerate}
\item \textbf{Constatif.} « Meka attendait sa médaille\ldots » décrit un état (3\textsuperscript{e} personne, imparfait) ; l'énoncé est vrai ou faux.
\item \textbf{Performatif.} « Je vous décore au nom de la France » : 1\textsuperscript{re} personne, présent, verbe performatif (\textit{décorer}) ; le commandant décore en le disant (déclaratif).
\item \textbf{Performatif.} « Je promets de venir\ldots » : le verbe \textit{promettre} accomplit l'engagement par le seul fait d'être prononcé (engageant).
\item \textbf{Constatif.} « Le commandant de cercle arriva au village » décrit un événement passé (passé simple, 3\textsuperscript{e} personne) ; vérifiable comme vrai ou faux.
\item \textbf{Performatif.} « Je vous remercie de votre fidélité » : remercier, c'est accomplir l'acte de gratitude en le disant (expressif).
\item \textbf{Constatif.} « La pluie tomba pendant toute la nuit » décrit un fait naturel ; énoncé descriptif, vrai ou faux.
\item \textbf{Performatif.} « Je déclare la séance ouverte » : l'énoncé ouvre effectivement la séance (déclaratif ; on peut préfixer « par la présente, je\ldots »).
\item \textbf{Constatif.} « Les villageois dansèrent jusqu'au matin » rapporte une action passée sans l'accomplir par le langage.
\end{enumerate}
\textbf{Points de vigilance :} la justification doit citer au moins deux critères : personne (1\textsuperscript{re} vs 3\textsuperscript{e}), temps (présent vs passé), nature du verbe (performatif ou non), et le test « dire = faire ».
\end{corrige}

\begin{corrige}[Exercice 6 --- Lexique spécialisé colonial]
Cinq termes du lexique spécialisé de l'administration coloniale (au choix parmi) :
\begin{enumerate}
\item \textbf{Commandant de cercle :} administrateur colonial français qui dirigeait un « cercle », circonscription administrative du territoire colonisé.
\item \textbf{Interprète :} auxiliaire africain chargé de traduire les propos entre l'administrateur français et les populations locales.
\item \textbf{Citation :} texte officiel lu lors d'une cérémonie, énumérant les services rendus qui justifient l'attribution d'une distinction.
\item \textbf{Résident :} représentant de l'autorité coloniale dans une circonscription, supérieur ou représentant local de l'administration.
\item \textbf{Canton :} subdivision administrative coloniale regroupant plusieurs villages, placée sous l'autorité d'un chef auxiliaire.
\end{enumerate}
(Autres réponses acceptées : \textbf{administrateur} = fonctionnaire colonial chargé de l'administration d'un territoire ; \textbf{notables} = personnalités locales reconnues et associées à l'administration ; \textbf{décoration} = distinction honorifique officielle.)
\textbf{Points de vigilance :} ne pas se contenter de recopier le mot : il faut le \textbf{définir dans le contexte colonial} (fonction, hiérarchie). Au Bac, on peut aussi demander l'effet produit : ce lexique crée un effet de réel et souligne le pouvoir de l'administration.
\end{corrige}

\begin{corrige}[Exercice 7 --- Reformulation guidée]
Propositions de condensation (le sens essentiel est conservé ; toute formulation équivalente est acceptée) :
\begin{enumerate}
\item \textit{Meka, ayant sacrifié fils et terres, attend fièrement sa médaille.} (10 mots --- relative transformée en participe, « fierté immense » $\rightarrow$ « fièrement ».)
\item \textit{Une pluie diluvienne nocturne ruina la fête et dispersa les invités.} (10 mots --- suppression des détails, nominalisation implicite.)
\item \textit{Après tant de promesses, les Blancs abandonnèrent Meka sans scrupule.} (10 mots --- subordonnée relative condensée en groupe prépositionnel.)
\item \textit{Humilié, Meka comprit enfin l'inutilité de sa médaille.} (9 mots --- la comparaison « morceau de métal » est résumée par « inutilité ».)
\item \textit{Les anciens raillaient la naïveté de Meka envers les Blancs.} (10 mots --- suppression du cadre « autour du feu », détail secondaire.)
\end{enumerate}
\textbf{Points de vigilance :} vérifier que l'idée essentielle (agent, action, enjeu) subsiste ; compter les mots ; ne garder ni guillemets ni style direct.
\end{corrige}

\begin{corrige}[Exercice 8 --- Repérage des idées essentielles]
\textbf{1. Les cinq idées essentielles (à souligner) :}
\begin{enumerate}
\item Manque criard de matériel pédagogique dans les lycées techniques du Centre.
\item Les ateliers de mécanique fonctionnent avec des machines vétustes.
\item Les promesses d'équipement n'ont pas été tenues.
\item Les enseignants dénoncent une préparation compromise du Baccalauréat technique.
\item Les parents ont demandé des mesures urgentes ; une délégation sera reçue.
\end{enumerate}
\textbf{2. Éléments secondaires (à barrer) :} « Depuis l'ouverture de l'année scolaire » (précision chronologique), « pourtant essentiels à la formation des élèves de Terminale » (précision), « vieilles de vingt ans » (détail chiffré illustratif), « réunis samedi dernier à Yaoundé » (détail de circonstance), « la semaine prochaine » (précision).
\textbf{3. Justifications :}
\begin{itemize}
\item Idée 1 : c'est le \textbf{sujet du texte} (thème annoncé dès la première phrase).
\item Idée 2 : elle \textbf{illustre et confirme} le manque par le cas des ateliers, cœur du problème.
\item Idée 3 : elle constitue un \textbf{argument} (paroles non suivies d'actes).
\item Idée 4 : elle exprime la \textbf{conséquence majeure} (compromission de l'examen).
\item Idée 5 : elle présente la \textbf{réaction et la perspective} (revendication, délégation) : c'est la chute informative du texte.
\end{itemize}
\textbf{Points de vigilance :} une idée essentielle est celle dont la suppression rend le texte incompréhensible ; un détail est ce qui peut disparaître sans nuire à la compréhension globale.
\end{corrige}

\begin{corrige}[Exercice 9 --- Résumé au quart --- type Baccalauréat technique]
\textbf{1. Idées essentielles classées en quatre étapes :}
\begin{itemize}
\item \textbf{Étape 1 (situation initiale) :} Meka a tout sacrifié pour la France (terres aux missionnaires, deux fils morts à la guerre) et attend la médaille promise.
\item \textbf{Étape 2 (préparatifs) :} Le matin de la cérémonie, Meka se prépare avec émotion ; tout le village est rassemblé.
\item \textbf{Étape 3 (cérémonie) :} Le commandant arrive ; on lit une citation incomprise et on épingle la médaille ; Meka rayonne, se croit l'égal des Blancs.
\item \textbf{Étape 4 (dénouement / désillusion) :} La pluie disperse la fête, les Blancs partent ; Meka tombe malade, mesure la vanité de la médaille ; les anciens se moquent de lui.
\end{itemize}
\textbf{2. Proposition de résumé (98 mots, dans la tolérance 90--110) :}

\emph{Meka, vieux chef, a tout sacrifié pour la France : ses terres données aux missionnaires, ses deux fils morts à la guerre. Il attend donc fièrement la médaille promise. Le jour venu, ému, il revêt son plus beau boubou devant le village rassemblé. L'administrateur lit une citation incompréhensible, puis épingle la décoration ; Meka rayonne, se croyant l'égal des Blancs. Cependant, une pluie diluvienne disperse la fête et les Blancs repartent sans un regard. Tombé malade, Meka mesure la vanité de ce morceau de métal, tandis que les anciens raillent sa naïveté.}

\textbf{3. Nombre de mots : 98 mots.} (Le candidat doit compter scrupuleusement et noter le total exact.)

\textbf{4. Connecteurs soulignés :}
\begin{itemize}
\item \emph{donc} : exprime la \textbf{conséquence} (les sacrifices justifient l'attente fière).
\item \emph{Cependant} : marque l'\textbf{opposition / le retournement} (de la gloire à la désillusion).
\end{itemize}
\textbf{Barème indicatif (sur 10) :} respect du quota 90--110 mots : 2 pts ; fidélité au texte (aucune idée inventée, aucune idée majeure omise) : 3 pts ; reformulation (pas de recopie, style indirect, 3\textsuperscript{e} personne) : 2 pts ; cohésion (connecteurs pertinents) : 2 pts ; correction linguistique : 1 pt.
\textbf{Points de vigilance :} ne pas écrire « je » ni donner son avis ; ne pas garder les guillemets (« Qu'est-ce que cela me rapporte ? » doit devenir une interrogation indirecte) ; indiquer le nombre de mots à la fin, sous peine de pénalité.
\end{corrige}

\begin{corrige}[Exercice 10 --- Correction d'un mauvais résumé]
\textbf{1. et 2. Erreurs relevées et règles violées :}
\begin{itemize}
\item \textbf{Recopie de phrases entre guillemets} (« Meka avait donné ses terres\ldots », « Le matin de la cérémonie\ldots ») : violation de la règle de \textbf{reformulation} --- toute recopie est assimilée à du plagiat.
\item \textbf{Avis personnel} (« Personnellement, je trouve que\ldots », « C'est vraiment scandaleux », « je le recommande ») : violation de la règle de \textbf{neutralité} --- le résumé ne contient que les idées de l'auteur.
\item \textbf{Dépassement du quota} (environ 150 mots au lieu de 100 $\pm 10\,\%$, soit 90--110) : violation de la \textbf{contrainte quantitative}.
\item \textbf{Présence d'éléments secondaires} (la voiture couverte de poussière, les tam-tams, l'anecdote personnelle de l'école du candidat) : violation de la règle de \textbf{sélection des idées essentielles} ; l'anecdote personnelle est en outre une invention hors texte.
\end{itemize}
\textbf{3. Version corrigée (91 mots, dans la tolérance) :}

\emph{Meka a tout donné à la France : ses terres aux missionnaires, ses deux fils morts à la guerre. Il attend donc avec fierté la médaille promise. Le jour venu, il se prépare avec émotion devant le village rassemblé. L'administrateur lit une citation, puis épingle la décoration ; Meka rayonne. Cependant, une pluie violente interrompt la fête et les Blancs s'en vont sans se retourner. Tombé malade, Meka comprend alors que cette médaille ne vaut rien face à ses fils disparus et à ses terres perdues, pendant que les anciens du village se moquent de sa crédulité.}

\textbf{Nombre de mots : 91.}
\textbf{Barème indicatif (sur 10) :} repérage des 4 erreurs : 4 pts (1 pt par erreur correctement identifiée et rattachée à sa règle) ; résumé corrigé : 6 pts (quota 1 pt, fond 2 pts, reformulation 2 pts, langue 1 pt).
\textbf{Points de vigilance :} dans la question 2, il ne suffit pas de nommer l'erreur : il faut \textbf{énoncer la règle violée} (« le résumé interdit\ldots »). Dans la version corrigée, compter les mots et vérifier l'absence totale de « je ».
\end{corrige}

\begin{corrige}[Exercice 11 --- Tonalités et bilan de l'œuvre --- type question de synthèse]
\textbf{Partie A. Identification des tonalités :}
\begin{enumerate}
\item \textbf{Tonalité satirique.} Procédés : \textbf{caricature} (« le ventre tendu dans son uniforme trop serré »), \textbf{ironie} (« comme s'il accomplissait le geste le plus généreux du siècle » : exagération qui souligne la vanité du personnage). Le rire vise à ridiculiser le représentant du pouvoir colonial.
\item \textbf{Tonalité polémique.} Procédés : \textbf{dénonciation directe} (« Ils ont pris nos terres, ils ont pris nos fils »), \textbf{anaphore} (« ils ont pris\ldots »), \textbf{comparaison dévalorisante} (« un morceau de métal qui ne vaut même pas une poule »), \textbf{exclamation} indignée. Le ton est celui de la colère et de l'accusation.
\item \textbf{Tonalité satirique} (teintée d'amertume). Procédés : \textbf{comparaison dévalorisante} (« comme un enfant montre un jouet »), \textbf{ironie} (le contraste entre la fierté de Meka et les rires des Blancs « dans son dos »). La naïveté du héros est montrée avec une distance moqueuse et pitoyable.
\end{enumerate}
\textbf{Partie B. Proposition de présentation de l'œuvre (une dizaine de lignes) :}

\emph{Le Vieux nègre et la médaille}, publié en 1956, est un roman du Camerounais Ferdinand Oyono (1929--2010), figure majeure de la littérature africaine d'expression française. Écrit à la veille des indépendances, il raconte l'histoire de Meka, vieux chef villageois qui a tout sacrifié pour la France --- ses terres, ses deux fils --- et qui attend en récompense une médaille. La cérémonie tourne à la farce et le héros, abandonné par les Blancs, découvre trop tard la vanité de sa fidélité. À travers cette désillusion, Oyono dénonce l'aliénation coloniale et l'illusion de l'assimilation. La tonalité dominante est satirique : l'ironie et la caricature ridiculisent l'administration coloniale ; mais le récit bascule par moments dans la polémique, lorsque la douleur et l'injustice sont dénoncées avec violence. Œuvre fondatrice du roman camerounais, elle affirme que la dignité ne s'achète pas avec des médailles.

\textbf{Barème indicatif (sur 10) :} Partie A : 4 pts (identification juste de chaque tonalité : 0,5 pt ; relevé et nom exact des procédés : 0,5 pt par extrait, arrondi). Partie B : 6 pts (auteur/date/contexte : 1,5 pt ; résumé de l'intrigue : 1 pt ; thèmes : 1 pt ; tonalité dominante justifiée : 1,5 pt ; portée de l'œuvre : 1 pt).
\textbf{Points de vigilance :} ne pas confondre satirique (rire, ironie, distance) et polémique (colère, dénonciation directe) ; nommer précisément les procédés (anaphore, caricature\ldots) et non seulement les citer ; dans la présentation, respecter l'ordre attendu : auteur $\rightarrow$ date/contexte $\rightarrow$ intrigue $\rightarrow$ thèmes $\rightarrow$ tonalité $\rightarrow$ portée.
\end{corrige}

\newpage

\coursec{Fiche bilan}

\begin{popfigure}
\begin{tikzpicture}[mindmap, grow cyclic, every node/.style={concept, text=white, font=\small}, concept color=popblue, level 1/.append style={level distance=3.6cm, sibling angle=60}, level 2/.append style={level distance=2.2cm, font=\scriptsize}]
\node[concept, font=\small\bfseries]{Le résumé de texte}
  child[concept color=poporange]{ node{Définition et enjeux}
    child{ node{Reformuler fidèlement} }
    child{ node{Respecter la longueur} }
  }
  child[concept color=popgreen]{ node{Lecture analytique}
    child{ node{Idées essentielles} }
    child{ node{Plan du texte} }
  }
  child[concept color=poppurple]{ node{Énoncés}
    child{ node{Constatif} }
    child{ node{Performatif} }
  }
  child[concept color=poppink]{ node{Lexique spécialisé}
    child{ node{Champ lexical} }
    child{ node{Termes techniques} }
  }
  child[concept color=popgold]{ node{Tonalités}
    child{ node{Satirique} }
    child{ node{Polémique} }
  }
  child[concept color=popblue!70]{ node{Rédaction et amélioration}
    child{ node{Reformulation} }
    child{ node{Relecture, correction} }
  };
\end{tikzpicture}
\legende{Carte mentale de la leçon : les six compétences mobilisées pour rédiger un résumé de texte au Probatoire / Baccalauréat technique.}
\end{popfigure}

\begin{aretenir}[L'essentiel de la leçon]
\textbf{Définitions clés}
\begin{itemize}
\item Le \cle{résumé} : reformulation brève, fidèle et personnelle des idées essentielles d'un texte, sans commentaire ni jugement.
\item La \cle{contraction de texte} : technique de réduction qui supprime exemples, répétitions et détails pour ne garder que l'essentiel.
\item L'\cle{énoncé constatif} : énoncé qui décrit un fait et peut être vrai ou faux (\emph{« Il pleut »}).
\item L'\cle{énoncé performatif} : énoncé qui accomplit l'acte qu'il dit (\emph{« Je promets », « Je déclare la séance ouverte »}).
\item Le \cle{lexique spécialisé} : ensemble de termes propres à un domaine (droit, médecine, technique...).
\item La tonalité \cle{satirique} : critique par la moquerie et l'ironie ; la tonalité \cle{polémique} : critique par l'attaque frontale et vive.
\end{itemize}

\textbf{Repères à connaître par cœur}
\begin{center}
\begin{tabularx}{\linewidth}{|l|X|}
\hline
\rowcolor{popblueL} \textbf{Notion} & \textbf{À retenir} \\
\hline
Longueur du résumé & Environ le quart du texte (souvent 100 mots au Probatoire / Baccalauréat technique), marge de 10 \% tolérée. \\
\hline
Étapes & Lecture analytique $\rightarrow$ repérage des idées $\rightarrow$ reformulation $\rightarrow$ rédaction $\rightarrow$ relecture. \\
\hline
Interdits & Copier des phrases, ajouter son avis, dépasser la limite de mots. \\
\hline
Performatif & Souvent à la 1\textsuperscript{re} personne du présent ; vérifiable par « par ces mots, je... ». \\
\hline
\end{tabularx}
\end{center}

\textbf{Méthodes express}
\begin{itemize}
\item Souligner une idée par paragraphe ; la formuler en une phrase personnelle.
\item Remplacer les exemples par un terme générique ; supprimer citations et chiffres secondaires.
\item Compter les mots et ajuster ; relire pour la grammaire et la cohérence.
\end{itemize}
\end{aretenir}

\begin{aretenir}[Checklist avant le Probatoire / Baccalauréat technique]
\begin{itemize}
\item[$\square$] Je sais définir le résumé et distinguer contraction et commentaire.
\item[$\square$] Je repère le thème et la thèse du texte dès la première lecture.
\item[$\square$] Je dégage une idée essentielle par partie du texte.
\item[$\square$] Je reformule avec mes propres mots, sans recopier de phrases.
\item[$\square$] Je supprime exemples, répétitions et détails secondaires.
\item[$\square$] Je distingue un énoncé constatif d'un énoncé performatif.
\item[$\square$] Je reconnais le lexique spécialisé d'un texte et son domaine.
\item[$\square$] Je sais identifier les tonalités satirique et polémique.
\item[$\square$] Je respecte la limite de mots imposée (et je les compte).
\item[$\square$] Je n'exprime jamais mon opinion personnelle dans le résumé.
\item[$\square$] Je rédige au présent de vérité générale, avec des connecteurs logiques.
\item[$\square$] Je relis pour corriger orthographe, accords et ponctuation.
\end{itemize}
\end{aretenir}

\findecours

\end{document}
